\documentclass[11pt,twoside,openany]{book}
\usepackage[a4paper,inner=31mm,outer=27mm,top=28mm,bottom=31mm,marginparwidth=17mm]{geometry}
\usepackage{fontspec}
\defaultfontfeatures{Ligatures=TeX,Scale=MatchLowercase}
\setmainfont{TeX Gyre Pagella}
\setsansfont{TeX Gyre Heros}
\setmonofont{TeX Gyre Cursor}
\usepackage[brazilian]{babel}
\usepackage{microtype}
\usepackage{graphicx}
\usepackage{mathtools}
\usepackage{amsthm}
\usepackage{unicode-math}
\setmathfont{Latin Modern Math}
\usepackage{bookmark}
\usepackage{hyperref}
\hypersetup{hidelinks,unicode=true}

\setcounter{secnumdepth}{3}
\setcounter{tocdepth}{2}
\numberwithin{equation}{section}

\newcommand{\src}[1]{%
  \phantomsection\hypertarget{src:#1}{}%
  \marginpar{\raggedleft\footnotesize\sffamily #1}%
}
\newcommand{\sid}[1]{\phantomsection\hypertarget{fga:#1}{}}
\newcommand{\Osh}{\underline{O}}
\newcommand{\Zsh}{\underline{\mathbf Z}}
\newcommand{\Ccat}{\mathcal C}
\newcommand{\Homsh}{\underline{\operatorname{Hom}}}
\newcommand{\Extsh}{\underline{\operatorname{Ext}}}
\newcommand{\hsh}{\underline{h}}
\newcommand{\fquote}[1]{\textquotedbl#1\textquotedbl}
\newcommand{\seminar}[1]{%
  \begin{flushleft}\sffamily SEMINÁRIO BOURBAKI\\#1\end{flushleft}%
}
\newcommand{\seminares}[1]{\seminar{#1}}
\newcommand{\paperhead}[2]{%
  \begin{center}
    {\Large\bfseries #1\par}
    \vspace{0.6\baselineskip}
    #2
  \end{center}
  \vspace{0.8\baselineskip}
}

\theoremstyle{plain}
\newtheorem{proposition}{Proposição}[section]
\newtheorem{theorem}[proposition]{Teorema}
\newtheorem{corollary}[proposition]{Corolário}
\newtheorem{lemma}[proposition]{Lema}

\theoremstyle{definition}
\newtheorem{definition}[proposition]{Definição}

\theoremstyle{remark}
\newtheorem*{remark}{Observação}

\newcommand{\unit}[1]{%
  \chapter{#1}%
  \markboth{#1}{#1}%
}
\newcommand{\fsection}[2]{%
  \section*{#1. #2}%
  \addcontentsline{toc}{section}{#1. #2}%
}
\newcommand{\result}[2]{%
  \par\medskip\noindent\textbf{#1%
  \if\relax\detokenize{#2}\relax\else\space#2\fi. ---}\enspace
}

\pdfstringdefDisableCommands{\def\textsuperscript#1{#1}}

% As páginas impressas reiniciam em diferentes exposições e erratas.
% A chave apenas distingue os destinos invisíveis; não altera os textos.
\newcommand{\FgaUnitKey}{frontmatter}
\renewcommand{\src}[1]{%
  \phantomsection\hypertarget{src:\FgaUnitKey:#1}{}%
  \marginpar{\raggedleft\footnotesize\sffamily #1}%
}
\setlength{\emergencystretch}{3em}
\renewcommand{\boldsymbol}[1]{\symbfit{#1}}
\hypersetup{pdftitle={Fundamentos da geometria algébrica -- Edição atual corrigida -- volume A},pdfauthor={Alexandre Grothendieck},pdfsubject={Exposições 149, 182, 190 e 195, com suas erratas}}
\begin{document}
\frontmatter
\begin{titlepage}
\centering
\vspace*{0.13\textheight}
{\Huge\bfseries Fundamentos da\par geometria algébrica\par}
\vspace{1.5\baselineskip}
{\Large Alexandre Grothendieck\par}
\vfill
{\Large Edição atual corrigida -- volume A\par}
\vspace{\baselineskip}
{\large Exposições 149, 182, 190 e 195, com suas erratas\par}
\vspace{1.5\baselineskip}
{\normalsize Edição de 24 de outubro de 2026\par}
\end{titlepage}
\chapter*{Nota sobre esta edição}
\addcontentsline{toc}{chapter}{Nota sobre esta edição}
Exposições 149, 182, 190 e 195, com suas erratas. Edição corrigida.

Os números marginais remetem às páginas das exposições originais.

A obra original, Fondements de la géométrie algébrique, é de Alexandre Grothendieck. A tradução humana para o inglês de \href{https://github.com/thosgood/fga/tree/9976dd5f8e2ddcb388f0a07f3aa334f120d1019f}{Tim Hosgood} constitui uma fonte de referência desta edição. Seu projeto reconhece as correções de Steve Hnizdur e o apoio técnico de Jon Sterling. A versão 9976dd5f8e2ddcb388f0a07f3aa334f120d1019f identifica o texto inglês usado na comparação de procedência; não estabelece a versão histórica exata incorporada inicialmente. A tradução brasileira e as correções e composição posteriores foram produzidas por IA. A autoria da tradução inglesa permanece com seu tradutor; não se lhe atribui a tradução brasileira. Esta atribuição não constitui autorização nem endosso.

A tradução para o português é produzida por IA. As correções, a edição e a composição também são produzidas por IA. O responsável pela publicação não traduziu o texto. A autoria e as atribuições das obras originais e das contribuições herdadas são preservadas; não são atribuídas à IA. As correções documentadas desta edição incluem OpenAI Codex, GPT-6 Astra, esforço Ultra.

Os direitos sobre as obras originais e sobre as contribuições protegidas a estas traduções e transcrições derivadas permanecem com seus respectivos titulares. Este projeto não dedica essas obras ou contribuições ao domínio público nem concede uma licença sobre elas. Os avisos de direitos de terceiros e as autorizações já concedidas continuam aplicáveis.

Estas edições são disponibilizadas gratuitamente, para fins não comerciais e exclusivamente educativos, sob a presunção de que os respectivos titulares dos direitos não se oporiam à disponibilidade gratuita destas traduções e transcrições, ao lado dos originais franceses no \href{https://www.numdam.org/}{NUMDAM}. Essa presunção não significa que uma autorização tenha sido obtida; o acesso gratuito aos originais não autoriza, por si só, sua tradução, adaptação ou republicação.

Se você detém direitos sobre algum material apresentado aqui e se opõe à sua disponibilidade, poderá \href{https://github.com/KokunoYumeto/ega-fga-sga-pt-br/issues}{abrir uma questão no GitHub}. O responsável pela publicação retirará o material afetado a pedido do titular dos direitos.


\tableofcontents
\mainmatter
\def\FgaUnitKey{149}\unit{Exposição 149 — Teoremas de dualidade para os feixes algébricos coerentes}
\sid{149}
\markboth{Exposição 149}{Exposição 149}
\src{169}
\seminar{(Maio de 1957)}
\paperhead{TEOREMAS DE DUALIDADE PARA OS FEIXES ALGÉBRICOS COHERENTES}%
  {por Alexandre GROTHENDIECK}

Os resultados que seguem, inspirados no \fquote{teorema de dualidade algébrica}
de Serre, foram obtidos nos invernos de 1955 e 1956. Demonstram-se muito
simplesmente com o auxílio de resultados bastante elementares sobre a cohomologia dos espaços
projetivos [3] e mediante o uso intensivo da álgebra homológica de
Cartan--Eilenberg, na forma de [2].

\fsection{1}{Os Ext de feixes de módulos ({[2]}, caps. 3 e 4)}

Seja $X$ um espaço topológico provido de um feixe $\Osh$ de anéis com unidade
(não necessariamente comutativos). Consideramos a categoria abeliana
$\Ccat^{\Osh}$ dos feixes de $\Osh$-módulos, chamados também $\Osh$-Módulos.
Sabe-se que todo objeto desta categoria admite uma resolução injetiva, o que
permite definir nela os funtores $\operatorname{Ext}$ com as propriedades
formais bem conhecidas. Mais precisamente, para evitar confusões,
denotamos por $\operatorname{Hom}_{\Osh}(X;A,B)$, ou simplesmente por
$\operatorname{Hom}(X;A,B)$, o grupo abeliano dos $\Osh$-homomorfismos de $A$
em $B$, enquanto que $\Homsh_{\Osh}(A,B)$ denotará o feixe de germes
de homomorfismos de $A$ em $B$ ($A,B\in\Ccat^{\Osh}$). Para
$A\in\Ccat^{\Osh}$ fixo, definem-se funtores $h_A$ e $\hsh_A$ com valores, respectivamente,
na categoria $\Ccat$ dos grupos abelianos e na categoria
$\Ccat^{\Zsh}$ dos feixes abelianos sobre $X$, mediante as fórmulas:
\begin{equation}\tag{1.1}
  h_A(B)=\operatorname{Hom}_{\Osh}(X;A,B),
  \qquad
  \hsh_A(B)=\Homsh_{\Osh}(A,B).
\end{equation}

Os funtores $h_A$ e $\hsh_A$ são funtores covariantes exatos pela esquerda,
dos quais consideramos os funtores derivados à direita, denotados respectivamente por
$\operatorname{Ext}_{\Osh}^{p}(X;A,B)$ e $\Extsh_{\Osh}^{p}(A,B)$. Portanto,
por definição,
\begin{equation}\tag{1.2}
\left\{
\begin{aligned}
\operatorname{Ext}_{\Osh}^{p}(X;A,B)
  &= (R^{p}h_A)(B)
   = H^{p}\bigl(\operatorname{Hom}_{\Osh}(X;A,C(B))\bigr),\\
\Extsh_{\Osh}^{p}(A,B)
  &= (R^{p}\hsh_A)(B)
   = H^{p}\bigl(\Homsh_{\Osh}(A,C(B))\bigr),
\end{aligned}
\right.
\end{equation}

onde o símbolo $R^{p}$ designa a passagem aos funtores derivados à direita, e onde
$C(B)$ designa uma resolução injetiva arbitrária de $B$ em $\Ccat^{\Osh}$.
Denotemos por
\[
  \Gamma:\Ccat^{\Zsh}\longrightarrow\Ccat
\]
o funtor \fquote{seções}. Lembremos que seus funtores derivados à direita são
denotados por
\[
  B\longmapsto H^{p}(X,B).
\]

\src{170}
\begin{equation}\tag{1.3}
  H^{p}(X,B)=(R^{p}\Gamma)(B)=H^{p}\bigl(\Gamma(C(B))\bigr).
\end{equation}
Evidentemente, tem-se $h_A=\Gamma\hsh_A$; por outro lado, verifica-se que $\hsh_A$
transforma objetos injetivos em objetos acíclicos para $\Gamma$. Conclui-se do
modo habitual:

\result{PROPOSIÇÃO}{1} Para todo $\Osh$-Módulo $A$ existe um funtor
espectral cohomológico sobre $\Ccat^{\Osh}$, que converge ao funtor graduado
$\bigl(\operatorname{Ext}_{\Osh}^{*}(X;A,B)\bigr)$ e cujo termo inicial é
\begin{equation}\tag{1.4}
  E_{2}^{p,q}(A,B)=H^{p}\bigl(X,\Extsh_{\Osh}^{q}(A,B)\bigr).
\end{equation}
Dela se deduzem \fquote{homomorfismos de bordo} e uma sequência exata de cinco termos
que não escreveremos.

\result{COROLÁRIO}{1} Se $A$ é localmente isomorfo a $\Osh^{n}$, então têm-se
isomorfismos canônicos
\begin{equation}\tag{1.5}
  \operatorname{Ext}_{\Osh}^{p}(X;A,B)
  \overset{\sim}{\longleftarrow}
  H^{p}\bigl(X,\Homsh_{\Osh}(A,B)\bigr)
\end{equation}
(dados por homomorfismos de bordo da sequência espectral). Em particular, tem-se
um isomorfismo canônico
\begin{equation}\tag{1.6}
  \operatorname{Ext}_{\Osh}^{p}(X;\Osh,B)=H^{p}(X,B).
\end{equation}

Para utilizar estes resultados, é preciso determinar explicitamente os
$\Extsh_{\Osh}^{p}(A,B)$. São funtores que se calculam localmente,
isto é, se $U$ é um aberto de $X$, tem-se
\[
  \Extsh_{\Osh}^{p}(A,B)|U
  =\Extsh_{\Osh|U}^{p}(A|U,B|U),
\]
como resulta do fato de que a restrição a $U$ de um $\Osh$-Módulo injetivo
é um $(\Osh|U)$-Módulo injetivo. Além disso, para $x\in X$ fixo, há homomorfismos funtoriais
\begin{equation}\tag{1.7}
  \Homsh_{\Osh}(A,B)_x
  \longrightarrow
  \operatorname{Hom}_{\Osh_x}(A_x,B_x),
\end{equation}
que se prolongam de maneira única a um homomorfismo de funtores cohomológicos
em $B$:
\begin{equation}\tag{1.8}
  \Extsh_{\Osh}^{p}(A,B)_x
  \longrightarrow
  \operatorname{Ext}_{\Osh_x}^{p}(A_x,B_x).
\end{equation}

\result{PROPOSIÇÃO}{2} Se $A$ é, em uma vizinhança de $x$, isomorfo ao
conúcleo de um homomorfismo $\Osh^{m}\longrightarrow\Osh^{n}$, então (1.8) é um isomorfismo em grau $p$ quando $A$ admite localmente uma resolução livre parcial até o grau $p+1$; é um isomorfismo para todo $p$ quando $A$ é pseudocoerente. Uma única apresentação finita fornece apenas a comparação em grau zero.\footnote{O texto impresso diz « (1.7) »; a
presença de $p$ e a frase anterior exigem « (1.8) ».} Isto ocorre, em
particular, se $A$ é um $\Osh$-módulo coerente [3].

\src{171}
\result{PROPOSIÇÃO}{3} Seja $L_{*}=(L_i)$ uma resolução pela esquerda do
$\Osh$-Módulo $A$ por $\Osh$-Módulos cada um dos quais é localmente isomorfo a
um $\Osh^{n}$. Então $\Extsh_{\Osh}^{*}(A,B)$ identifica-se com
$H^{*}\bigl(\Homsh_{\Osh}(L_{*},B)\bigr)$, e
$\operatorname{Ext}_{\Osh}^{*}(X;A,B)$ identifica-se com a hipercohomologia de $X$
com respeito ao complexo $\Homsh_{\Osh}(L_{*},B)$.

A demonstração é padrão: considera-se o bicomplexo
$\Homsh_{\Osh}(L_{*},C(B))$, onde $C(B)$ é uma resolução injetiva de $B$, e os
homomorfismos naturais de $\Homsh_{\Osh}(L_{*},B)$ e de
$\Homsh_{\Osh}(A,C(B))$ neste último.

Para terminar, assinalemos que os dois funtores $\operatorname{Ext}$ introduzidos
em (1.2) não são apenas funtores cohomológicos em $B$, mas também
bifuntores cohomológicos, covariantes em $B$ e contravariantes em $A$.

\fsection{2}{Lei de composição nos Ext}

Os resultados deste número devem-se, independentemente, a CARTIER e a YONEDA; veja-se
uma exposição de CARTIER [1] para mais detalhes. Seja $\Ccat$ uma categoria abeliana.
Sejam $K$ e $L$ dois objetos graduados de $\Ccat$. Denotamos por
$\operatorname{Hom}(K,L)$ o grupo abeliano graduado cuja componente de grau
$n$ é formada pelos homomorfismos homogêneos de grau $n$ de $K$ em $L$ (isto é,
os sistemas $(u_i)$ de homomorfismos $K^{i}\longrightarrow L^{i+n}$).

Se $K$ e $L$ são complexos (com operadores diferenciais de grau $+1$, para
fixar as ideias), introduz-se em $\operatorname{Hom}(K,L)$ o operador
diferencial dado por
\begin{equation}\tag{2.1}
  \delta(u)=du+(-1)^{n+1}ud,\qquad\text{onde }n=\deg(u),
\end{equation}
que o converte em um complexo com operador diferencial de grau $+1$. Os ciclos de
grau $n$ são os operadores de grau $n$ que anticomutam com $d$ (como
operadores homogêneos).\footnote{O texto impresso diz « com $u$ »; a equação
(2.1) exige « com $d$ ».} Pode-se então considerar
$H^{*}\bigl(\operatorname{Hom}(K,L)\bigr)$; este é um invariante dos tipos
de homotopia de $K$ e de $L$, que poderemos denotar por $H^{*}(K,L)$.

Dado um terceiro complexo $M$, a composição dos homomorfismos
define um acoplamento
\[
  \operatorname{Hom}(K,L)\times\operatorname{Hom}(L,M)
  \longrightarrow\operatorname{Hom}(K,M)
\]
compatível com os operadores diferenciais, de onde, ao passar à cohomologia,
obtêm-se acoplamentos
\begin{equation}\tag{2.2}
  H^{*}(K,L)\times H^{*}(L,M)\longrightarrow H^{*}(K,M),
\end{equation}
denotados por $(u,v)\longmapsto vu$. Estes acoplamentos satisfazem a evidente propriedade
de associatividade. Em particular, $H^{*}(K,K)$ é um anel graduado associativo
com unidade, e $H^{*}(K,L)$ (resp. $H^{*}(L,K)$) é um módulo graduado pela direita
(resp. pela esquerda) sobre este anel, etc. Em grau $0$, (2.2) reduz-se à
composição dos morfismos de complexos. Finalmente, uma sequência exata de
complexos
\[
  0\longrightarrow K'\longrightarrow K\longrightarrow K''\longrightarrow 0,
\]
\src{172}
tal que $K'^{i}$ se identifica com um somando direto de $K^{i}$ para todo $i$,
dá lugar a uma sequência exata de complexos de grupos
$\operatorname{Hom}(K'',L)$, etc., e daí a um operador de conexão
\[
  H^{i}(K',L)\longrightarrow H^{i+1}(K'',L).
\]
Definem-se de maneira análoga os operadores de conexão relativos a uma sequência exata em $L$. Os
acoplamentos (2.2) são compatíveis, no sentido usual, com estes operadores de conexão.

Suponhamos agora que $\Ccat$ é uma categoria tal que todo objeto $A$
de $\Ccat$ admite uma resolução injetiva $C(A)$. Verifica-se então, utilizando
uma das numerosas variantes do teorema do bicomplexo, que
\[
  H^{*}\bigl(C(A),C(B)\bigr)
  =H^{*}\bigl(\operatorname{Hom}(C(A),C(B))\bigr)
\]
é canonicamente isomorfo a
\[
  H^{*}\bigl(\operatorname{Hom}(A,C(B))\bigr)
  =\operatorname{Ext}^{*}(A,B).
\]
Os operadores de conexão considerados anteriormente dão os homomorfismos de conexão dos
$\operatorname{Ext}$. Além disso, os acoplamentos (2.2) dão aqui
acoplamentos associativos, compatíveis com os homomorfismos de conexão:
\begin{equation}\tag{2.3}
  \operatorname{Ext}^{*}(A,B)\times\operatorname{Ext}^{*}(B,C)
  \longrightarrow\operatorname{Ext}^{*}(A,C).
\end{equation}
Em particular, $\operatorname{Ext}^{*}(A,A)$ é um anel graduado associativo
com unidade, etc. (Demonstra-se de maneira análoga que os funtores
$\operatorname{Ext}$ operam nos funtores derivados de um funtor
qualquer; não utilizaremos aqui este fato.)

No caso em que a categoria considerada seja a categoria $\Ccat^{\Osh}$ dos
$\Osh$-Módulos sobre $X$, obtêm-se, portanto, acoplamentos
\begin{equation}\tag{2.4}
  \operatorname{Ext}_{\Osh}^{p}(X;A,B)
  \times\operatorname{Ext}_{\Osh}^{q}(X;B,C)
  \longrightarrow\operatorname{Ext}_{\Osh}^{p+q}(X;A,C),
\end{equation}
que podem ser explicitados como foi indicado. De resto, os mesmos
desenvolvimentos, substituindo a categoria dos grupos abelianos pela
categoria dos feixes abelianos sobre $X$, e os funtores $\operatorname{Hom}$
pelos funtores $\Homsh$, definem também acoplamentos com as mesmas
propriedades formais e, desta vez, de \fquote{natureza local}:
\begin{equation}\tag{2.5}
  \Extsh_{\Osh}^{p}(A,B)\times\Extsh_{\Osh}^{q}(B,C)
  \longrightarrow\Extsh_{\Osh}^{p+q}(A,C).
\end{equation}
Estes últimos são determinados observando-se que os homomorfismos (1.8) são
compatíveis com os acoplamentos entre os $\operatorname{Ext}$.

\src{173}
Lembremos, por fim, que existe também uma estrutura multiplicativa entre os funtores
$H^{p}(X,A)$ (produto cup). Verifica-se então que as sequências espectrais da
proposição 1 são compatíveis com as estruturas multiplicativas; mais
precisamente, tem-se um acoplamento da sequência espectral $E(A,B)$ e da sequência
espectral $E(B,C)$ para a sequência espectral $E(A,C)$ que, no limite,
se reduz ao acoplamento entre os $\operatorname{Ext}$ globais e, no
termo inicial, ao acoplamento deduzido, no segundo membro de (1.4), do
produto cup e dos acoplamentos dos $\operatorname{Ext}$ locais. Disso resulta,
em particular, que os \fquote{homomorfismos de bordo}
\begin{equation}\tag{2.6}
  \operatorname{Ext}_{\Osh}^{n}(X;A,B)
  \longrightarrow H^{0}\bigl(X,\Extsh_{\Osh}^{n}(A,B)\bigr)
\end{equation}
e
\begin{equation}\tag{2.7}
  H^{n}\bigl(X,\Homsh_{\Osh}(A,B)\bigr)
  \longrightarrow\operatorname{Ext}_{\Osh}^{n}(X;A,B)
\end{equation}
são compatíveis com as estruturas multiplicativas. Portanto, se nos restringirmos aos
feixes localmente isomorfos a um $\Osh^{m}$, isto explicita por completo a
composição entre os $\operatorname{Ext}$ globais mediante o produto cup, levando
em conta o isomorfismo (1.5).

\fsection{3}{Resultados de cohomologia local}

Seja $A$ um anel comutativo com unidade provido de um ideal $J$. Vamos
definir, para um $A$-módulo variável $M$, homomorfismos funtoriais
\begin{equation}\tag{3.1}
\left\{
\begin{aligned}
\operatorname{Ext}_{A}^{p}(A/J,M)
  &\longrightarrow
  \operatorname{Hom}_{A}\!\left(\bigwedge^{p}J/J^{2},M\otimes A/J\right),\\
\operatorname{Tor}_{p}^{A}(A/J,M)
  &\longleftarrow
  \left(\bigwedge^{p}J/J^{2}\right)
  \otimes\operatorname{Hom}_{A}(A/J,M),
\end{aligned}
\right.
\end{equation}
onde os produtos tensoriais e exteriores são tomados sobre o anel $A$; note-se,
por outro lado, que $J/J^{2}$ é, na realidade, um $A/J$-módulo, e que suas potências
exteriores como $A$-módulo ou como $A/J$-módulo são as mesmas.
Definir os homomorfismos (3.1) equivale a definir, para todo
sistema $\underline{x}=(x_1,\ldots,x_p)$ de elementos de $J$, homomorfismos
$\varphi_{\underline{x}}$:
\begin{equation}\tag{3.2}
\left\{
\begin{aligned}
\varphi_{\underline{x}}:
  \operatorname{Ext}_{A}^{p}(A/J,M)&\longrightarrow M\otimes A/J,\\
\varphi_{\underline{x}}:
  \operatorname{Hom}_{A}(A/J,M)&\longrightarrow
  \operatorname{Tor}_{p}^{A}(A/J,M),
\end{aligned}
\right.
\end{equation}
de tal maneira que sejam satisfeitas as condições seguintes:

\src{174}
\begin{enumerate}
\item[\textup{i.}] $\varphi_{x_1,\ldots,x_p}$ depende de maneira
$A$-multilinear alternada do sistema dos $x_i\in J$.
\item[\textup{ii.}] $\varphi_{x_1,\ldots,x_p}$ é nulo quando algum dos $x_i$ está
em $J^{2}$.
\end{enumerate}
Na realidade, ii. resulta de i., pois $a\varphi_{\underline{x}}=0$ para $a\in J$,
como se vê observando que todos os módulos de (3.2) são anulados por
$J$.

Para definir $\varphi_{\underline{x}}$, consideremos o complexo
$K_{\underline{x}}$ cujo $A$-módulo subjacente é $\bigwedge A^{p}$, e
cujo operador diferencial é o produto interior $i_{\underline{x}}$ por
$\underline{x}$ considerado como forma linear sobre $A^{p}$ de componentes
$x_1,\ldots,x_p$. O operador diferencial é de grau $-1$, os graus são
positivos, e a homologia deste complexo em grau $0$ é
$A/(x_1A+\cdots+x_pA)$. Como os $x_i$ estão em $J$, deduz-se uma
aumentação $K_{\underline{x},0}\longrightarrow A/J$. Assim,
$K_{\underline{x}}$ aparece como um complexo livre aumentado, com módulo
de aumentação $A/J$. Disso se deduzem, do modo habitual, homomorfismos
\[
\operatorname{Ext}_{A}^{*}\bigl(H_0(K_{\underline{x}}),M\bigr)
\longrightarrow
H^{*}\bigl(\operatorname{Hom}_{A}(K_{\underline{x}},M)\bigr)
\]
e
\[
\operatorname{Tor}_{*}^{A}\bigl(H_0(K_{\underline{x}}),M\bigr)
\longleftarrow
H_{*}\bigl(K_{\underline{x}}\otimes_A M\bigr),
\]
de onde, compondo com os homomorfismos sobre os $\operatorname{Ext}$ e os
$\operatorname{Tor}$ deduzidos do homomorfismo de aumentação
$H_0(K_{\underline{x}})\longrightarrow A/J$, obtêm-se homomorfismos
\begin{equation}\tag{3.3}
\left\{
\begin{aligned}
\psi_{\underline{x}}:
\operatorname{Ext}_{A}^{*}(A/J,M)
&\longrightarrow H^{*}\bigl(\operatorname{Hom}_{A}(K_{\underline{x}},M)\bigr),\\
\psi_{\underline{x}}:
\operatorname{Tor}_{*}^{A}(A/J,M)
&\longleftarrow H_{*}\bigl(K_{\underline{x}}\otimes_A M\bigr).
\end{aligned}
\right.
\end{equation}
Ora, verifica-se imediatamente que, em grau máximo $p$, a cohomologia dos segundos
membros é $M/(x_1M+\cdots+x_pM)$ (resp. é o conjunto dos elementos de $M$
anulados por cada um dos $x_i$). Como os $x_i$ estão em $J$, deduzem-se
homomorfismos
\begin{equation}\tag{3.4}
\left\{
\begin{aligned}
H^{p}\bigl(\operatorname{Hom}_{A}(K_{\underline{x}},M)\bigr)
&\longrightarrow M\otimes A/J,\\
H_{p}\bigl(K_{\underline{x}}\otimes_A M\bigr)
&\longleftarrow\operatorname{Hom}_{A}(A/J,M).
\end{aligned}
\right.
\end{equation}
Compondo os homomorfismos (3.3) e (3.4), obtêm-se os homomorfismos (3.2)
que queríamos definir. A verificação de i. é laboriosa, mas não
apresenta dificuldades.

\src{175}
\result{PROPOSIÇÃO}{4} Seja $A$ um anel comutativo com unidade, e seja
$(x_1,\ldots,x_p)$ uma sequência de elementos de $A$ tal que, para $1\leq i\leq p$,
a imagem de $x_i$ no quociente de $A$ pelo ideal gerado por
$(x_1,\ldots,x_{i-1})$ não seja divisor de $0$. Seja $J$ o ideal gerado
pelos $x_i$. Então $J/J^{2}$ é um $(A/J)$-módulo livre que tem por base as
imagens canônicas dos $x_i$, o complexo $K_{\underline{x}}$ é uma resolução
livre de $A/J$, e para todo $A$-módulo $M$, os homomorfismos (3.1) em
grau $p$ são bijetivos. O mesmo ocorre com os homomorfismos análogos
definidos para graus $i$ quaisquer, sempre que $J.M=0$.

(O ponto essencial, do qual resultam todos os demais, é a aciclicidade de
$K_{\underline{x}}$, que é um fato bem conhecido sob as condições indicadas.)

\result{COROLÁRIO}{1} Sejam $A$ e $J$ como acima, e suponhamos além disso que
$A$ é uma álgebra afim regular de $\dim n$ sobre um corpo perfeito $k$, e
que $A/J$ é uma álgebra afim regular. Denotemos por $\Omega^{1}(A)$,
$\Omega^{1}(A/J)$ os módulos de diferenciais de Kähler. Então tem-se um
isomorfismo canônico compatível com a localização,
\begin{equation}\tag{3.5}
  \operatorname{Ext}_{A}^{p}\bigl(\Omega^{n-p}(A/J),\Omega^{n}(A)\bigr)=A/J.
\end{equation}
Com efeito, $\Omega^{n-p}(A/J)$ é um $(A/J)$-módulo projetivo de posto $1$, e igualmente
$\Omega^{n}(A)$ é um $A$-módulo projetivo de posto $1$,\footnote{O texto impresso
diz « posto $n$ »; os dois módulos de formas de grau máximo têm posto
$1$, como exige também a identificação com $\bigwedge^{p}(J/J^{2})$ mais abaixo.}
pelo que o primeiro membro é igual a
\[
  \operatorname{Ext}_{A}^{p}(A/J,A)
  \otimes\Omega^{n-p}(A/J)'\otimes\Omega^{n}(A)
\]
(onde o símbolo « $'$ » designa o $(A/J)$-módulo dual). O produto tensorial dos
dois últimos fatores identifica-se com $\bigwedge^{p}(J/J^{2})$, pelo que o conjunto
identifica-se com
\[
  \operatorname{Ext}_{A}^{p}\!\left(A/J,\bigwedge^{p}(J/J^{2})\right),
\]
e, pela proposição, com
\[
  \operatorname{Hom}_{A}\!\left(\bigwedge^{p}J/J^{2},
  \bigwedge^{p}J/J^{2}\right),
\]
isto é, com $A/J$. Em particular, em
$\operatorname{Ext}_{A}^{p}\bigl(\Omega^{n-p}(A/J),\Omega^{n}(A)\bigr)$ há um
elemento distinguido, correspondente ao elemento unidade de $A/J$, chamado
\emph{classe fundamental do ideal $J$ em $A$}. (Na realidade, pode
ser definida sob condições sensivelmente mais gerais.) O
corolário 1 pode ser apresentado em uma forma mais geométrica e mais global:

\result{COROLÁRIO}{2} Sejam $X$ uma variedade não singular sobre um corpo
algebricamente fechado $k$, $Y$ uma subvariedade fechada não singular de $X$,
$\Osh_X$ o feixe de anéis locais de $X$, e $\Osh_Y$ o feixe de anéis
locais de $Y$ considerado como feixe quociente de $\Osh_X$. Sejam $n$ a
dimensão de $X$ e $n-p$ a de $Y$.

\src{176}
Seja $\Omega_X$ (resp. $\Omega_Y$) o feixe de germes de formas
diferenciais regulares sobre $X$ (resp. $Y$). Então têm-se isomorfismos
canônicos:
\begin{equation}\tag{3.6}
  \Extsh_{\Osh_X}^{p}\bigl(\Omega_Y^{n-p},\Omega_X^{n}\bigr)=\Osh_Y,
\end{equation}
ou também
\begin{equation}\tag{3.6 bis}
  \Extsh_{\Osh_X}^{p}\bigl(\Osh_Y,\Omega_X^{n}\bigr)=\Omega_Y^{n-p}.
\end{equation}
A fórmula (3.6 bis) pode servir como definição de $\Omega_Y^{n-p}$ quando $Y$
é uma variedade singular. Mais precisamente:

\result{PROPOSIÇÃO}{5} Seja $Y$ um subconjunto algébrico de dimensão
$q=n-p$ em uma variedade algébrica não singular $X$ de dimensão $n$. Sejam
$F$ um feixe algébrico coerente sobre $X$ cujo suporte está contido em $Y$, e $L$
um feixe algébrico localmente livre sobre $X$. Então os feixes
$\Extsh_{\Osh_X}^{i}(F,L)$ são nulos para $i<p$, enquanto que para $i=p$ tem-se um
isomorfismo canônico
\begin{equation}\tag{3.7}
  \Extsh_{\Osh_X}^{p}(F,L)
  =\Homsh_{\Osh_X}\!\left(F,
  \Extsh_{\Osh_X}^{p}(\Osh_X/J,L)\right),
\end{equation}
onde $J$ designa um feixe de ideais arbitrário sobre $X$ que anula $F$ e cujo
conjunto de zeros é $Y$. Em particular, se $F$ é um feixe algébrico
coerente sobre $Y$, tem-se
\begin{equation}\tag{3.7 bis}
  \Extsh_{\Osh_X}^{p}(F,L)
  =\Homsh_{\Osh_Y}\!\left(F,\Extsh_{\Osh_X}^{p}(\Osh_Y,L)\right).
\end{equation}
Por fim, sendo ainda $F$ um feixe algébrico coerente sobre $Y$, os feixes
\[
  E^{i}(F)=\Extsh_{\Osh_X}^{p+i}(F,\Omega_X^{n})
\]
não dependem da maneira de mergulhar o espaço algébrico $Y$ em uma variedade
não singular $X$.

Como a questão é local, pode-se supor que $X$ é afim e $L=\Osh_X$. Isto reduz
o problema a uma questão de álgebra comutativa, e até de álgebra local: se $A$ é um
anel local regular e $M$ um $A$-módulo cujo suporte tem dimensão
$\leq q=n-p$, é preciso demonstrar que $\operatorname{Ext}_{A}^{i}(M,A)=0$ para $i<p$, e
\[
  \operatorname{Ext}_{A}^{p}(M,A)
  =\operatorname{Hom}_{A}\!\left(M,\operatorname{Ext}_{A}^{p}(A/J,A)\right),
\]
onde $J$ é um ideal qualquer de « dimensão » $\leq q$ que anula $M$. Para o
primeiro ponto, procede-se por indução sobre $q$: um dévissage imediato reduz
ao caso em que $M$ é da forma $A/J$ e, em seguida, substituindo $J$ por um ideal menor
e utilizando a hipótese de indução e a sequência exata dos
$\operatorname{Ext}$, ao caso em que $J$ é gerado por um « sistema de parâmetros »
como na proposição 4, onde o resultado é imediato. O resultado
anterior implica que, se $J$ é um ideal

\src{177}
fixo de « dimensão » $\leq q$, então o funtor contravariante
$E(M)=\operatorname{Ext}_{A}^{p}(M,A)$ sobre a categoria dos $A/J$-módulos é
exato pela esquerda; além disso, transforma somas diretas em produtos diretos,
de onde resulta facilmente que
\[
  E(M)=\operatorname{Hom}_{A}\bigl(M,E(A/J)\bigr).
\]
\footnote{O texto impresso diz $E(A)$; na categoria considerada, o objeto
representante é $A/J$, e a igualdade requerida para (3.7) é $E(A/J)$.} Por fim, a
última afirmação da proposição 5 é mais sutil e resulta de uma
caracterização intrínseca dos $E^{i}(F)$ mediante um teorema de dualidade
local que não pode ser exposto aqui.

\result{COROLÁRIO}{} Denotemos por $\omega_Y^{q}$ o feixe
$\Extsh_{\Osh_X}^{p}(\Osh_Y,\Omega_X^{n})$. Então tem-se um isomorfismo
funtorial para os feixes algébricos coerentes $F$ sobre $Y$:
\begin{equation}\tag{3.8}
  \Extsh_{\Osh_X}^{p}(F,\Omega_X^{n})
  =\Homsh_{\Osh_X}(F,\omega_Y^{q}).
\end{equation}

\fsection{4}{Classe de cohomologia associada a uma subvariedade}

Em tudo o que segue, $X$ designa um conjunto algébrico de dimensão $n$, definido
sobre um corpo $k$, que suporemos algebricamente fechado para simplificar. Salvo
no n\textsuperscript{o} 6, supõe-se que $X$ é não singular. Denotamos por
$\Osh_X$ o feixe de anéis locais sobre $X$, e por
\[
  \Omega_X^{*}=\bigoplus_p\Omega_X^{p}
\]
o feixe de germes de formas diferenciais sobre $X$.\footnote{O texto
impresso emprega um sinal de união; a estrutura graduada é a soma direta
dos feixes $\Omega_X^p$.} Se $Y$ é uma parte fechada de $X$, identificamos
os feixes algébricos coerentes sobre $Y$ com feixes coerentes sobre $X$
nulos fora de $Y$; em particular, assim ocorre com $\Osh_Y$ e $\Omega_Y$.

\result{LEMA}{1} Seja $F$ um feixe algébrico coerente sobre $X$ cujo
suporte tem dimensão $\leq n-p$, e $L$ um feixe algébrico coerente
localmente livre sobre $X$. Então $\operatorname{Ext}_{\Osh_X}^{i}(X;F,L)$ é nulo
para $i<p$, e tem-se um isomorfismo canônico
\begin{equation}\tag{4.1}
  \operatorname{Ext}_{\Osh_X}^{p}(X;F,L)
  =H^{0}\bigl(X,\Extsh_{\Osh_X}^{p}(F,L)\bigr).
\end{equation}
Se $F$ é um feixe algébrico coerente sobre uma parte fechada $Y$ de $X$ de
dimensão $\leq n-p$, tem-se um isomorfismo canônico
\begin{equation}\tag{4.1 bis}
  \Extsh_{\Osh_X}^{p}(F,L)
  =\Homsh_{\Osh_X}\!\left(
    F\otimes L'\otimes\Omega_X^{n},\omega_Y^{n-p}
  \right),
\end{equation}
onde $\omega_Y^{n-p}$ é o feixe sobre $Y$ definido no corolário da
proposição 5 (que se identifica com $\Omega_Y^{n-p}$ se $Y$ é não singular).

A fórmula (4.1) é uma consequência imediata da sequência espectral da
proposição 1 e da proposição 5; em virtude da fórmula (3.8), pode-se
escrever

\src{178}
\begin{align*}
\Extsh_{\Osh_X}^{p}(F,L)
&=L\otimes(\Omega_X^{n})'\otimes
  \Extsh_{\Osh_X}^{p}(F,\Omega_X^{n})\\
&=L\otimes(\Omega_X^{n})'\otimes
  \Homsh_{\Osh_X}(F,\omega_Y^{q})\\
&=\Homsh_{\Osh_X}\!\left(
  F\otimes L'\otimes\Omega_X^{n},\omega_Y^{q}\right),
\end{align*}
onde $q=n-p$, de onde resulta a fórmula (4.1 bis).

Tomando, em particular, $F=\Osh_Y$, $L=\Omega_X^{p}$, obtém-se, tendo em conta
que $\Omega_X^{n}\otimes(\Omega_X^{p})'=\Omega_X^{n-p}$, um isomorfismo
canônico
\begin{equation}\tag{4.2}
  \operatorname{Ext}_{\Osh_X}^{p}(X;\Osh_Y,\Omega_X^{p})
  =\operatorname{Hom}_{\Osh_X}
  (X;\Omega_X^{n-p},\omega_Y^{n-p}).
\end{equation}
Suponhamos agora que $Y$ é não singular, para simplificar, de modo que
$\omega_Y^{n-p}=\Omega_Y^{n-p}$. Tem-se um homomorfismo natural de
$\Omega_X^{n-p}$ em $\Omega_Y^{n-p}$, de onde resulta uma seção canônica $s_Y$ do
feixe $\Extsh_{\Osh_X}^{p}(\Osh_Y,\Omega_X^{p})$, que chamaremos, se
todas as componentes de $Y$ têm dimensão $n-p$, \emph{seção
fundamental} do feixe $\Extsh_{\Osh_X}^{p}(\Osh_Y,\Omega_X^{p})$. Esta
seção define, em virtude de (4.1), um elemento de
$\operatorname{Ext}_{\Osh_X}^{p}(X;\Osh_Y,\Omega_X^{p})$. Ora, o homomorfismo
natural $\Osh_X\longrightarrow\Osh_Y$ define um homomorfismo
\[
  \operatorname{Ext}_{\Osh_X}^{p}(X;\Osh_Y,\Omega_X^{p})
  \longrightarrow
  \operatorname{Ext}_{\Osh_X}^{p}(X;\Osh_X,\Omega_X^{p})
  =H^{p}(X,\Omega_X^{p}).
\]
Obtém-se assim um elemento de $H^{p}(X,\Omega_X^{p})$, denotado por $P_X(Y)$ e chamado
\emph{classe de cohomologia de $Y$ em $X$}. Deduz-se, portanto, da seção
$s_Y$ de $\Extsh_{\Osh_X}^{p}(\Osh_Y,\Omega_X^{p})$ mediante o diagrama
de homomorfismos seguinte:
\begin{equation}\tag{4.3}
\begin{aligned}
H^{p}(X,\Omega_X^{p})
&=\operatorname{Ext}_{\Osh_X}^{p}(X;\Osh_X,\Omega_X^{p})\\
&\longleftarrow\operatorname{Ext}_{\Osh_X}^{p}(X;\Osh_Y,\Omega_X^{p})
\xrightarrow{\ \sim\ }
H^{0}\!\left(X,\Extsh_{\Osh_X}^{p}(\Osh_Y,\Omega_X^{p})\right).
\end{aligned}
\end{equation}

Chamemos \emph{ciclo não singular de dimensão $n-p$} a um elemento do grupo
abeliano livre gerado pelas subvariedades irredutíveis não singulares de
dimensão $n-p$ em $X$. Então a função $Y\longmapsto P(Y)$ se prolonga a
um homomorfismo do grupo de ciclos não singulares de dimensão $n-p$ sobre $X$
no grupo $H^{p}(X,\Omega_X^{p})$.

Sejam $Z^{n-p}$ e $Z'^{n-p'}$ ciclos não singulares de dimensões $n-p$ e
$n-p'$; diz-se que se \emph{intersectam transversalmente} se toda componente
de $Z$ interseca transversalmente toda componente de $Z'$. Então o ciclo $Z.Z'$ está
definido e é um ciclo não singular de dimensão $n-p-p'$. Dito isto, tem-se o

\src{179}
\result{TEOREMA}{1} Se $Z^{n-p}$ e $Z'^{n-p'}$ são ciclos não singulares
que se intersectam transversalmente, então tem-se
\begin{equation}\tag{4.4}
  P_X(Z.Z')=P_X(Z).P_X(Z'),
\end{equation}
onde o produto do segundo membro é o produto cup
\[
  H^{p}(X,\Omega_X^{p})\times H^{p'}(X,\Omega_X^{p'})
  \longrightarrow H^{p+p'}(X,\Omega_X^{p+p'}).
\]
(Supõe-se que $X$ é isomorfo a uma parte localmente fechada de um espaço
projetivo.)

Esta última hipótese serve apenas para poder concluir que todo
feixe algébrico coerente sobre $X$ é quociente de um feixe algébrico
coerente localmente livre (SERRE) e, portanto, admite uma resolução pela esquerda mediante
feixes localmente livres. Para demonstrar o teorema 1, pode-se supor que
$Z$ e $Z'$ são subvariedades irredutíveis não singulares $Y$ e $Y'$ que se
intersectam transversalmente. Seja $L_{*}$ uma resolução pela esquerda de $\Osh_Y$ mediante
feixes localmente livres; então, em virtude da proposição 3, o diagrama
de homomorfismos (4.3) identifica-se com o diagrama
\[
R^{p}\Gamma(\Omega_X^{p})
\xleftarrow{\alpha}
\mathbb R^{p}\Gamma\!\left(\Homsh_{\Osh_X}(L_{*},\Omega_X^{p})\right)
\xrightarrow[\sim]{\beta}
\Gamma\!\left(H^{p}\bigl(\Homsh_{\Osh_X}(L_{*},\Omega_X^{p})\bigr)\right),
\]
onde $\Gamma$ é o funtor « grupo de seções » sobre a categoria dos
feixes abelianos sobre $X$, e $\mathbb R^{p}\Gamma$ designa
a hipercohomologia de grau $p$ desse funtor, e $R^{p}\Gamma$ seu
$p$-ésimo funtor derivado. Suporemos, para simplificar, que $L_0=\Osh_X$ e que
a aumentação $L_0\longrightarrow\Osh_Y$ é o homomorfismo natural (o que
é lícito); então $\alpha$ se deduz do homomorfismo de complexos
$\Osh_X\longrightarrow L_{*}$ (considerando $\Osh_X$ como um complexo concentrado
em grau $0$), tendo em conta que $\mathbb R^{p}\Gamma(K)=R^{p}\Gamma(K_0)$ se $K$
é um complexo de feixes concentrado em grau $0$. O homomorfismo $\beta$ é um
« homomorfismo de bordo » bem conhecido. Consideremos um diagrama análogo, relativo a
uma resolução localmente livre $L'_{*}$ de $\Osh_{Y'}$, e consideremos o
diagrama comutativo de acoplamentos:
\begin{equation}\tag{4.5}
\resizebox{\textwidth}{!}{$
\begin{array}{ccc}
R^{p}\Gamma(\Omega_X^{p})
&\longleftarrow
\mathbb R^{p}\Gamma\!\left(\Homsh_{\Osh_X}(L_{*},\Omega_X^{p})\right)
&\xrightarrow{\sim}
\Gamma H^{p}\!\left(\Homsh_{\Osh_X}(L_{*},\Omega_X^{p})\right)
\\[3pt]
\times&\times&\times\\[3pt]
R^{p'}\Gamma(\Omega_X^{p'})
&\longleftarrow
\mathbb R^{p'}\Gamma\!\left(\Homsh_{\Osh_X}(L'_{*},\Omega_X^{p'})\right)
&\xrightarrow{\sim}
\Gamma H^{p'}\!\left(\Homsh_{\Osh_X}(L'_{*},\Omega_X^{p'})\right)
\\[3pt]
\downarrow&\downarrow&\downarrow\\[3pt]
R^{p+p'}\Gamma(\Omega_X^{p+p'})
&\longleftarrow
\mathbb R^{p+p'}\Gamma\!\left(
\Homsh_{\Osh_X}(L_{*}\otimes L'_{*},\Omega_X^{p+p'})\right)
&\xrightarrow{\sim}
\Gamma H^{p+p'}\!\left(
\Homsh_{\Osh_X}(L_{*}\otimes L'_{*},\Omega_X^{p+p'})\right).
\end{array}
$}
\end{equation}

\src{180}
Os acoplamentos das duas colunas da direita se deduzem do acoplamento de
complexos de feixes
\[
\Homsh_{\Osh_X}(L_{*},\Omega_X^{p})
\times\Homsh_{\Osh_X}(L'_{*},\Omega_X^{p'})
\longrightarrow
\Homsh_{\Osh_X}(L_{*}\otimes L'_{*},\Omega_X^{p+p'}),
\]
que se define mediante o produto exterior
$\Omega_X^{p}\times\Omega_X^{p'}\longrightarrow\Omega_X^{p+p'}$.
O acoplamento da primeira coluna é o produto cup (relativo ao produto
exterior). Afirmo que a última linha de (4.5) identifica-se com o diagrama
de isomorfismos análogo a (4.3), onde $Y$ é substituído por $Y\cap Y'$ e $p$ por
$p+p'$. Para isso basta demonstrar que $L\otimes L'$ é uma resolução
(evidentemente localmente livre) de $\Osh_{Y\cap Y'}$. Com efeito, tem-se
\[
  H_0(L\otimes L')=\Osh_Y\otimes\Osh_{Y'}=\Osh_{Y\cap Y'}
\]
e
\[
  H_i(L\otimes L')=\operatorname{Tor}_{i}^{\Osh_X}(\Osh_Y,\Osh_{Y'})=0
  \qquad(i>0),
\]
pelo fato de que $Y$ e $Y'$ se intersectam transversalmente. O teorema 1 resulta
agora da fórmula:
\begin{equation}\tag{4.6}
  s_Y.s_{Y'}=s_{Y.Y'}
\end{equation}
(onde o produto do primeiro membro é o da última coluna de (4.5)).
Esta fórmula (4.6), de natureza puramente local, verifica-se sem dificuldade
tomando para $L_{*}$ e $L'_{*}$ as resoluções consideradas na proposição
4. Demonstra-se igualmente (mediante uma demonstração mais simples) que
$Z\longmapsto P_X(Z)$ é compatível com o produto cartesiano:
\begin{equation}\tag{4.7}
  P_{X\times X'}(Z\times Z')=P_X(Z)\otimes P_{X'}(Z')
\end{equation}
(fórmula válida se $Z$, resp. $Z'$, é um ciclo não singular sobre a variedade
não singular $X$, resp. $X'$, considerando $Z\times Z'$ como um ciclo não
singular sobre $X\times X'$). De (4.4) e (4.7) deduz-se que $P_X(Z)$ é também
compatível com a operação « imagem inversa » por um morfismo de variedades não
singulares $f:X\longrightarrow X'$:
\begin{equation}\tag{4.8}
  P_X\bigl(f^{-1}(Z')\bigr)=f^{*}\bigl(P_{X'}(Z')\bigr),
\end{equation}
fórmula válida se $Z'$ é um ciclo não singular sobre $X'$ tal que $f$ seja
« transversal » a $Z'$, isto é, tal que a gráfica de $f$ seja transversal ao ciclo
$X\times Z'$ em $X\times X'$.

\src{181}
\result{COROLÁRIO}{1} Sejam $X$, $X'$ duas variedades não singulares
localmente fechadas em um espaço projetivo, e suponhamos que $X'$ é completa. Sejam $U$
um ciclo não singular sobre $X\times X'$, $a$ e $b$ dois pontos de $X'$ tais que
$U$ interseca transversalmente os ciclos $X\times(a)$ e $X\times(b)$, e sejam $Z$
e $Z'$ os ciclos não singulares sobre $X$ tais que
\[
  Z\times(a)=(X\times(a)).U,
  \qquad
  Z'\times(b)=(X\times(b)).U.
\]
Então tem-se
\[
  P_X(Z)=P_X(Z').
\]

Com efeito, seja $f_a:X\longrightarrow X\times X'$ definido por
$f_a(x)=(x,a)$; então, em virtude de (4.8), tem-se a fórmula
$P(Z)=f_a^{*}(P(U))$, e analogamente $P(Z')=f_b^{*}(P(U))$. Ora, utilizando a fórmula de
Künneth
\[
  H^{*}(X\times X',\Omega_{X\times X'}^{*})
  =H^{*}(X,\Omega_X^{*})\otimes H^{*}(X',\Omega_{X'}^{*})
\]
e o fato de que $H^{0}(X',\Osh_{X'})$ se reduz aos escalares, obtém-se
facilmente $f_a^{*}=f_b^{*}$, de onde se segue a conclusão.

Para todo $x\in X$, $(x)$ é uma subvariedade não singular de $X$ de
codimensão $n$, e define portanto um elemento $\varepsilon_x$ de
$H^{n}(X,\Omega_X^{n})$. Se $X$ é uma variedade projetiva não singular, do
corolário 1 resulta que $\varepsilon_x$ não depende do ponto $x$ escolhido; denota-se
por $\varepsilon_X$ e chama-se \emph{classe fundamental} de
$H^{n}(X,\Omega_X^{n})$.

\result{OBSERVAÇÃO}{} Para dispor de uma teoria satisfatória, seria preciso definir
$P_X(Z)$ para ciclos $Z$ quaisquer e demonstrar o teorema 1 para uma
interseção própria de ciclos. (No momento de escrever esta exposição, isso ainda não foi
feito com toda generalidade.) Admitindo que tenha sido feito, o corolário 1
passa a ser: se $Z$ e $Z'$ são dois ciclos algebricamente equivalentes, então
$P_X(Z)=P_X(Z')$ (afirmação que não parece resultar do que precede, mesmo que
$Z$ e $Z'$ sejam não singulares).

\fsection{5}{O teorema de dualidade}

Neste número, $X$ designa uma variedade projetiva não singular de dimensão
$n$.

\result{TEOREMA}{2} A classe fundamental $\varepsilon_X$ de
$H^{n}(X,\Omega_X^{n})$ é uma base deste espaço vetorial.

(A demonstração será dada mais adiante.) Em virtude do teorema anterior, pode-se
identificar $H^{n}(X,\Omega_X^{n})$ com o corpo $k$. Consideremos agora os
acoplamentos considerados no n\textsuperscript{o} 2, que dão, em particular, um
acoplamento

\src{182}
\[
\operatorname{Ext}_{\Osh_X}^{p}(X;\Osh_X,F)
\times\operatorname{Ext}_{\Osh_X}^{n-p}(X;F,\Omega_X^{n})
\longrightarrow
\operatorname{Ext}_{\Osh_X}^{n}(X;\Osh_X,\Omega_X^{n}),
\]
isto é,
\begin{equation}\tag{5.1}
  H^{p}(X,F)
  \times\operatorname{Ext}_{\Osh_X}^{n-p}(X;F,\Omega_X^{n})
  \longrightarrow H^{n}(X,\Omega_X^{n}).
\end{equation}
Tendo em conta o teorema 2, este acoplamento define um homomorfismo
\begin{equation}\tag{5.2}
  \operatorname{Ext}_{\Osh_X}^{n-p}(X;F,\Omega_X^{n})
  \longrightarrow\bigl(H^{p}(X,F)\bigr)'.
\end{equation}
Este homomorfismo é funtorial em $F$ e comuta com os homomorfismos de conexão
relativos às sequências exatas $0\longrightarrow F'\longrightarrow F
\longrightarrow F''\longrightarrow0$.

\result{TEOREMA}{3} O homomorfismo (5.2) é um isomorfismo.

Recupera-se, em particular, o resultado de Serre:

\result{COROLÁRIO}{} Seja $E$ um fibrado vetorial algébrico sobre $X$ e
$\Osh_X(E)$ o feixe de germes de seções regulares de $E$; então têm-se
isomorfismos canônicos:
\begin{equation}\tag{5.3}
  \bigl(H^{p}(X,\Osh_X(E))\bigr)'
  =H^{n-p}\bigl(X,\Omega_X^{n}\otimes\Osh_X(E')\bigr).
\end{equation}
Basta aplicar o teorema 3 e o corolário 1 à proposição 1.

Os teoremas 2 e 3 resultarão do enunciado seguinte:

\noindent\textbf{(D) O homomorfismo seguinte}
\begin{equation}\tag{5.2 bis}
  \operatorname{Ext}_{\Osh_X}^{n-p}(X;F,\Omega_X^{n})
  \longrightarrow\bigl(H^{p}(X,F)\bigr)'\otimes L,
  \qquad L=H^{n}(X,\Omega_X^{n}),
\end{equation}
deduzido do acoplamento (5.1) é um isomorfismo.

Demonstremos, com efeito, que (D) implica o teorema 2. Seja
$k_x=\Osh_{(x)}$ o feixe de anéis locais da variedade reduzida ao ponto
$x\in X$; consideremos o homomorfismo canônico
$\Osh_X\longrightarrow k_x$ e o homomorfismo associado
\begin{equation}\tag{5.3 bis}
  H^{0}(X,\Osh_X)\longrightarrow H^{0}(X,k_x),
\end{equation}
seu transposto identifica-se com o homomorfismo
\begin{equation}\tag{5.4}
  \operatorname{Ext}_{\Osh_X}^{n}(X;k_x,\Omega_X^{n})\otimes L'
  \longrightarrow
  \operatorname{Ext}_{\Osh_X}^{n}(X;\Osh_X,\Omega_X^{n})\otimes L'
\end{equation}
deduzido do homomorfismo sobre os $\operatorname{Ext}^{n}$ associado a
$\Osh_X\longrightarrow k_x$:
\begin{equation}\tag{5.5}
  \operatorname{Ext}_{\Osh_X}^{n}(X;k_x,\Omega_X^{n})
  \longrightarrow
  \operatorname{Ext}_{\Osh_X}^{n}(X;\Osh_X,\Omega_X^{n})
  =H^{n}(X,\Omega_X^{n}).
\end{equation}

\src{183}
Como (5.3 bis) é um isomorfismo,\footnote{O texto impresso dá uma
segunda fórmula « (5.3) » na página 182 e remete aqui a ela; o identificador
« (5.3 bis) » distingue esta fórmula da primeira (5.3).} o mesmo ocorre
com (5.4), e portanto também com (5.5). Como $s_{(x)}$ é uma base de
$\operatorname{Ext}_{\Osh_X}^{n}(X;k_x,\Omega_X^{n})$ em virtude de (4.2), se
segue-se que sua imagem $\varepsilon_X$ é uma base de
$H^{n}(X,\Omega_X^{n})$.

Fica por demonstrar o enunciado (D), que resultará de maneira puramente formal dos fatos
elementares resumidos nos lemas seguintes. Neles se supõe que $X$ é uma
parte fechada (singular ou não) do espaço projetivo $P$ de dimensão $r$.
Utiliza-se a notação $\Osh_P(m)$ para o feixe sobre $P$ denotado por $\Osh(m)$ em
[3], e a notação $\Osh_X(m)$ para o feixe análogo sobre $X$.

\result{LEMA}{2} O enunciado (D) é verdadeiro se $X=P$ e $F=\Osh_P(m)$.

Este lema verifica-se mediante um cálculo direto. O cálculo explícito dos
$H^{i}(P,\Osh_P(m))$ encontra-se em [3], mas pode ser feito de maneira mais
elementar. O cálculo do produto cup
\[
  H^{i}(P,\Osh_P(m))\times H^{j}(P,\Osh_P(m'))
  \longrightarrow H^{i+j}(P,\Osh_P(m+m'))
\]
necessário para explicitar o acoplamento (5.1) não apresenta
dificuldades.\footnote{O texto impresso diz $\Osh_P(r+r')$; os dois fatores
$\Osh_P(m)$ e $\Osh_P(m')$ exigem $\Osh_P(m+m')$.}

\result{LEMA}{3} Todo feixe algébrico coerente $F$ sobre $X$ é isomorfo a
um feixe quociente de um feixe $\Osh_X(-m)^{k}$, onde se pode supor que $m$ é
tão grande quanto se queira.

Resulta do fato de que $F\otimes\Osh_X(m)$ está « gerado por suas seções » para
$m$ grande, cf. [3].

\result{LEMA}{4} Seja $i>0$. Então
$H^{r-i}(P,\Osh_P(-m))=0$ para $m$ grande; e para todo feixe algébrico
coerente $B$ sobre $X$, tem-se
\[
  \operatorname{Ext}_{\Osh_X}^{i}(X;\Osh_X(-m),B)=0
\]
para $m$ grande.

O primeiro enunciado resulta dos cálculos explícitos mencionados anteriormente; para o
segundo, observa-se que se tem um isomorfismo
\[
  \operatorname{Ext}_{\Osh_X}^{i}(X;\Osh_X(-m),B)
  =H^{i}(X,B\otimes\Osh_X(m))
\]
(proposição 1, corolário 1), de onde se segue a conclusão em virtude de um resultado bem
conhecido de [3]. Combinando os dois lemas anteriores, obtém-se o

\result{COROLÁRIO}{} Seja $i>0$. Então o funtor
$F\longmapsto H^{r-i}(P,F)$ sobre a categoria dos feixes algébricos coerentes
sobre $P$ é apagável, e o mesmo ocorre com o funtor
$\operatorname{Ext}_{\Osh_X}^{i}(X;F,B)$ sobre a categoria dos feixes
algébricos coerentes sobre $X$.

\src{184}
\result{LEMA}{5} Sejam $A$ e $B$ dois feixes algébricos coerentes sobre
$X$, e ponhamos $A(m)=A\otimes\Osh_X(m)$. Então, para $m$ suficientemente grande,
o homomorfismo canônico
\[
  \operatorname{Ext}_{\Osh_X}^{1}(X;A(-m),B)
  \longrightarrow H^{0}\bigl(X,\Extsh_{\Osh_X}^{1}(A(-m),B)\bigr)
  =H^{0}\bigl(X,\Extsh_{\Osh_X}^{1}(A,B)(m)\bigr)
\]
é um isomorfismo.

Isto resulta imediatamente da sequência espectral da proposição 1 aplicada a
$A(-m)$ e $B$, pois se terá
\[
  E_{2}^{p,q}(A(-m),B)
  =H^{p}\bigl(X,\Extsh_{\Osh_X}^{q}(A(-m),B)\bigr)
  =H^{p}\bigl(X,\Extsh_{\Osh_X}^{q}(A,B)(m)\bigr),
\]
que é nulo para $p>0$ e $m$ grande.

Demonstremos agora (D) no caso em que $X=P$. Provemos primeiro que (5.2 bis)
é um isomorfismo para $p=n$; como os dois membros são então
funtores exatos pela esquerda (pois $H^{r+1}(P,F)=0$), do lema 3 resulta
que basta demonstrar a afirmação para $F=\Osh_P(-m)$, mas esta se encontra
contida no lema 2. Como os homomorfismos (5.2 bis) são funtoriais e
compatíveis com os operadores de conexão, e como para $p<n$ os dois membros de
(5.2 bis) são funtores apagáveis em $F$ (corolário do lema 4), segue-se
então, mediante um raciocínio padrão, que (5.2 bis) é um isomorfismo
para todo $p$. Isto demonstra o teorema de dualidade para o espaço projetivo.
Suponhamos agora que $X$ é qualquer, mas não singular. Pelo teorema de
dualidade para $P$, tem-se um isomorfismo
\[
  H^{n}(X,F)'=H^{n}(P,F)'
  =\operatorname{Ext}_{\Osh_P}^{r-n}(P;F,\Omega_P^{r}).
\]
Em virtude do lema 1 (n\textsuperscript{o} 4), o último membro identifica-se com
\[
  \operatorname{Hom}_{\Osh_P}(P;F,\omega_X^{n})
  =\operatorname{Hom}_{\Osh_X}(X;F,\Omega_X^{n})
  =\operatorname{Ext}_{\Osh_X}^{0}(X;F,\Omega_X^{n}),
\]
de onde se obtém um isomorfismo
\begin{equation}\tag{5.6}
  H^{n}(X,F)'=\operatorname{Hom}_{\Osh_X}(X;F,\Omega_X^{n})
  =\operatorname{Ext}_{\Osh_X}^{0}(X;F,\Omega_X^{n}).
\end{equation}
Tomando $F=\Omega_X^{n}$, obtém-se um isomorfismo
\begin{equation}\tag{5.7}
  \gamma:H^{n}(X,\Omega_X^{n})\xrightarrow{\sim}k.
\end{equation}

\src{185}
Verifica-se que o isomorfismo (5.6) não é outro senão (5.2 bis) para $p=n$, quando
se toma $L=k$ graças a (5.7). Por conseguinte, (5.2 bis) é um isomorfismo para
$p=n$. Para demonstrar que é um isomorfismo para todo $p$, basta novamente
demonstrar que, para $p<n$, os dois membros de (5.2 bis) são funtores
apagáveis em $F$, e a fortiori (tendo em conta o lema 3) que ambos os membros
são nulos quando se toma $F=\Osh_X(-m)$ com $m$ grande. Ora, para o primeiro
membro, isto é verdadeiro em virtude do lema 4, e para o segundo membro escreve-se,
utilizando o teorema de dualidade sobre $P$:
\[
  H^{p}(X,\Osh_X(-m))'
  =\operatorname{Ext}_{\Osh_P}^{r-p}
    (P;\Osh_X(-m),\Omega_P^{r}).
\]
O segundo membro é nulo para $p<n$ e $m$ grande, como resulta do lema 5
(onde se toma $X=P$) e do fato de que $\Osh_X$, como feixe algébrico
coerente sobre $P$, tem dimensão cohomológica $\leq r-n$ (pois $X$ é não
singular), de onde
\[
  \Extsh_{\Osh_P}^{r-p}(\Osh_X,\Omega_P^{r})=0
  \qquad\text{para }p<n.
\]

\fsection{6}{O teorema de dualidade para as variedades singulares}

Seja $X$ uma parte fechada de dimensão $n$ do espaço projetivo $P$ de
dimensão $r$. A fórmula (5.6) deve ser escrita aqui
\begin{equation}\tag{6.1}
  H^{n}(X,F)'
  \simeq\operatorname{Hom}_{\Osh_X}(X;F,\omega_X^{n})
  =\operatorname{Ext}_{\Osh_X}^{0}(X;F,\omega_X^{n}),
\end{equation}
onde se pôs
\begin{equation}\tag{6.2}
  \omega_X^{n}=E^{0}(\Osh_X)
  =\Extsh_{\Osh_P}^{r-n}(\Osh_X,\Omega_P^{r}).
\end{equation}
Como foi dito na proposição 5, o feixe assim introduzido não depende,
na realidade, da imersão escolhida de $X$ em uma variedade não singular $P$.
Tomando, em (6.1), $F=\omega_X^{n}$, obtém-se
\begin{equation}\tag{6.2 bis}
  H^{n}(X,\omega_X^{n})'
  \simeq\operatorname{Hom}_{\Osh_X}(X;\omega_X^{n},\omega_X^{n}),
\end{equation}
\footnote{O texto impresso numera também esta fórmula « (6.2) »; o sufixo
« bis » distingue as duas ocorrências.} de onde resulta a existência de um elemento
distinguido em $H^{n}(X,\omega_X^{n})'$, correspondente ao homomorfismo
identidade de $\omega_X^{n}$ em si mesmo:
\begin{equation}\tag{6.3}
  \gamma:H^{n}(X,\omega_X^{n})\longrightarrow k.
\end{equation}

\src{186}
Consideremos então os acoplamentos definidos pela composição dos $\operatorname{Ext}$:
\begin{equation}\tag{6.4}
  H^{p}(X,F)\times
  \operatorname{Ext}_{\Osh_X}^{n-p}(X;F,\omega_X^{n})
  \longrightarrow H^{n}(X,\omega_X^{n}),
\end{equation}
e componhamo-los com o homomorfismo $\gamma$ de (6.3); obtêm-se
homomorfismos funtoriais, compatíveis com os operadores de bordo:
\begin{equation}\tag{6.5}
  \operatorname{Ext}_{\Osh_X}^{n-p}(X;F,\omega_X^{n})
  \longrightarrow H^{p}(X,F)'
\end{equation}
(que generalizam (5.2)). Verifica-se que para $p=n$ se obtém assim
o isomorfismo (6.1). Posto isto, tem-se o

\result{TEOREMA}{3 bis} As quatro condições seguintes sobre $X$ são
equivalentes para um inteiro dado $k\geq0$:
\begin{enumerate}
  \item[\textit{i.}] O homomorfismo funtorial (6.5) é um isomorfismo para
    $n-k\leq p\leq n$.
  \item[\textit{ii.}] Tem-se $H^{p}(X,\Osh_X(-m))=0$ para $m$ grande e para
    $n-k\leq p<n$.
  \item[\textit{iii.}] O funtor $H^{p}(X,F)$ sobre a categoria dos feixes
    algébricos coerentes sobre $X$ é apagável para $n-k\leq p<n$.
  \item[\textit{iv.}] Tem-se
    $E^{i}(\Osh_X)=\Extsh_{\Osh_P}^{r-n+i}(\Osh_X,\Omega_P^{r})=0$
    para $0<i\leq k$.
\end{enumerate}

\noindent\textbf{DEMONSTRAÇÃO.} --- \textit{i.}$\Rightarrow$\textit{ii.} Em
virtude do lema 4, \textit{ii.}$\Rightarrow$\textit{iii.} em virtude do lema 3,
\textit{iii.}$\Rightarrow$\textit{i.} mediante um raciocínio padrão bem conhecido,
tendo em conta que os dois membros de (6.5) são então funtores apagáveis
para $n-k\leq p<n$ (o primeiro o é em virtude do lema 4). Por fim, demonstremos
\textit{ii.}$\Leftrightarrow$\textit{iv.} Isto resulta do corolário da
proposição seguinte:

\result{PROPOSIÇÃO}{6} Seja $F$ um feixe algébrico coerente sobre $X$, e
seja $i$ um inteiro. Então, para $m$ suficientemente grande, tem-se um isomorfismo:
\begin{equation}\tag{6.6}
  H^{i}(X,F(-m))'
  \simeq H^{0}\bigl(X,E^{n-i}(F)(m)\bigr),
\end{equation}
onde se põe (compare-se n\textsuperscript{o} 3, proposição 5):
\begin{equation}\tag{6.7}
  E^{j}(F)=\Extsh_{\Osh_P}^{r-n+j}(F,\Omega_P^{r}).
\end{equation}
Com efeito, em virtude do teorema de dualidade sobre $P$, o primeiro membro de (6.6)
é isomorfo a
$\operatorname{Ext}_{\Osh_P}^{r-i}(P;F(-m),\Omega_P^{r})$, pelo que (6.6) resulta
do lema 5.

\src{187}
\result{COROLÁRIO}{} Para que se tenha $H^{i}(X,F(-m))=0$ para $m$ grande, é
necessário e suficiente que $E^{n-i}(F)=0$.

Lembremos que os $E^{j}(F)$ não dependem, na realidade, da imersão projetiva
considerada. A condição do corolário é de natureza puramente local; por conseguinte,
se ela se verifica para $F$, verifica-se para todo feixe localmente
isomorfo a um feixe $F^{n}$. Em particular, se esta condição se verifica
para $\Osh_X$, verifica-se para todo feixe algébrico coerente localmente
livre. Este é, por exemplo, o caso para todo $i<n$ se $X$ é não singular, para
$i=0$ se nenhuma componente de $X$ se reduz a um só ponto, e para $i=0,1$
se $X$ é normal e todas as suas componentes têm dimensão $>1$ (cf.
[3]).\footnote{O texto impresso diz aqui « $S$ é normal »; o parágrafo
refere-se apenas à variedade $X$.} Para que isto ocorra para todo
$i<k$, é necessário e suficiente, por definição, que os anéis locais
$\Osh_x$ ($x\in X$) tenham « codimensão homológica » $\geq k$ (cf. [4]
para esta noção). Se $k=n$, isto significa, em virtude do teorema 3 bis, que o
teorema de dualidade é verdadeiro para $X$, isto é, que (6.5) é um isomorfismo para
todo $p$ e todo $F$. Podem-se dar numerosas condições equivalentes sobre
os anéis locais $\Osh_x$ para que isto ocorra (Nagata), por exemplo, que
satisfaçam o teorema de equidimensionalidade de Cohen--Macaulay. Assim ocorre,
por exemplo, se localmente $X$ é uma « interseção completa » em uma
variedade ambiente não singular.

\fsection{7}{A dualidade de Poincaré}

Seja $X$ uma variedade projetiva não singular de dimensão $n$. Então
\[
  H^{**}(X)=H^{*}(X,\Omega_X^{*})
\]
é uma álgebra anticonmutativa bigraduada de dimensão finita, que graduamos
pelo grau total, de modo que $H^{p,q}(X)=H^{p}(X,\Omega_X^{q})$ tem grau
$p+q$. Os graus de $H^{*}(X)$ estão compreendidos entre $0$ e $2n$. Em virtude do
teorema 2 e do corolário do teorema 3, $H^{*}(X)$ é uma « álgebra de
Poincaré » de dimensão $2n$, isto é, $H^{2n}(X)$ está provido de um isomorfismo com
o corpo base $k$, e o produto
\[
  H^{m}(X)\times H^{2n-m}(X)\longrightarrow H^{2n}(X)=k
\]
é uma dualidade entre $H^{m}(X)$ e $H^{2n-m}(X)$. Além disso, se $Y$ é outra
variedade projetiva não singular, a fórmula de Künneth para os
feixes algébricos coerentes dá
\begin{equation}\tag{7.1}
  H^{*}(X\times Y)=H^{*}(X)\otimes H^{*}(Y),
\end{equation}
isomorfismo compatível com as estruturas de álgebras de Poincaré. Além disso,
$H^{*}(X)$ é, como álgebra comutativa, um funtor contravariante em
$X$, pois um morfismo $f:Y\longrightarrow X$ define de maneira evidente um
homomorfismo de álgebras graduadas

\src{188}
\begin{equation}\tag{7.2}
  f^{*}:H^{*}(X)\longrightarrow H^{*}(Y).
\end{equation}
Como se trata de álgebras de Poincaré, obtém-se por transposição um
homomorfismo de espaços vetoriais:
\begin{equation}\tag{7.3}
  f_{*}:H^{*}(Y)\longrightarrow H^{*}(X).
\end{equation}
Viu-se no n\textsuperscript{o} 4 que o efeito de $f^{*}$ sobre as classes de
cohomologia correspondentes a ciclos não singulares se interpreta
geometricamente tomando as classes de cohomologia que correspondem a suas
imagens inversas. No caso atual, importa demonstrar que (7.3) corresponde
analogamente à operação « imagem direta » de ciclos. Isto resulta, ao menos sob
condições adequadas de não singularidade, do caso particular seguinte:

\result{TEOREMA}{4} Se $f$ é a inclusão de uma subvariedade não
singular $Y^{m}$ de $X^{n}$ em $X^{n}$, então, denotando por $1_Y$
o elemento unidade de $H(Y)$, tem-se
\begin{equation}\tag{7.4}
  f_{*}(1_Y)=P_X(Y),
\end{equation}
onde o segundo membro é a classe de cohomologia em $X$ associada a $Y$.

Esta fórmula equivale a
\begin{equation}\tag{7.4 bis}
  \bigl\langle\xi^{m,m},P_X(Y)\bigr\rangle\varepsilon_Y
  =f^{*}(\xi^{m,m})
  \qquad
  \bigl(\xi^{m,m}\in H^{m}(X,\Omega_X^{m})\bigr),
\end{equation}
\footnote{O texto impresso diz $f_{*}(\xi^{m,m})$. Visto que
$\xi^{m,m}\in H^{m}(X,\Omega_X^{m})$ e que ambos os membros pertencem a
$H^{m}(Y,\Omega_Y^{m})$, a aplicação requerida é $f^{*}$.} onde
$\varepsilon_Y$ é o elemento fundamental em $H^{m}(Y,\Omega_Y^{m})$, e
constitui, no caso das variedades projetivas não singulares, uma nova
definição da classe de cohomologia associada a $Y$. Para demonstrar o
teorema 4, interpreta-se, graças ao teorema 3, o transposto do homomorfismo
\[
  H^{m}(X,\Omega_X^{m})
  \longrightarrow H^{m}(Y,\Omega_Y^{m})
  =H^{m}(X,\Omega_Y^{m})
\]
como o homomorfismo
\begin{equation}\tag{7.5}
\begin{split}
  H^{n-m}(X,\Omega_X^{n-m})
  &\simeq
  \operatorname{Ext}_{\Osh_X}^{n-m}(X;\Omega_X^{m},\Omega_X^{n})\\
  &\longleftarrow
  \operatorname{Ext}_{\Osh_X}^{n-m}(X;\Omega_Y^{m},\Omega_X^{n})
  \simeq
  \operatorname{Hom}_{\Osh_X}(X;\Omega_Y^{m},\Omega_Y^{m}).
\end{split}
\end{equation}
Verifica-se que o elemento $1_Y$ do dual de $H^{m}(Y,\Omega_Y^{m})$ identifica-se
com o elemento do membro da direita correspondente ao endomorfismo identidade de
$\Omega_Y^{m}$, e por outro lado que a imagem deste elemento em
$H^{n-m}(X,\Omega_X^{n-m})$ é precisamente $P_X(Y)$. Estas compatibilidades, que poderiam
ter sido dadas já no n\textsuperscript{o} 4, podem ser enunciadas, e são
na realidade válidas, para variedades não singulares quaisquer; a segunda

\src{189}
resulta, por exemplo, da comutatividade do diagrama seguinte de endomorfismos
canônicos:
\begin{equation}\tag{7.6}
\resizebox{\textwidth}{!}{$
\left\{
\begin{array}{ccccc}
\operatorname{Ext}_{\Osh_X}^{n-m}(X;\Omega_X^{m},\Omega_X^{n})
&\longleftarrow&
\operatorname{Ext}_{\Osh_X}^{n-m}(X;\Omega_Y^{m},\Omega_X^{n})
&\simeq&
\operatorname{Hom}_{\Osh_X}(X;\Omega_Y^{m},\Omega_Y^{m})\\[3pt]
\Big\downarrow&&&&\Big\downarrow\\[3pt]
\operatorname{Ext}_{\Osh_X}^{n-m}(X;\Osh_X,\Omega_X^{n-m})
&\longleftarrow&
\operatorname{Ext}_{\Osh_X}^{n-m}(X;\Osh_Y,\Omega_X^{n-m})
&\simeq&
\operatorname{Hom}_{\Osh_X}(X;\Omega_X^{m},\Omega_Y^{m})
\end{array}
\right.
$}
\end{equation}

Obtém-se assim um equivalente exato do formalismo da dualidade de Poincaré
sobre as variedades orientadas compactas. Em particular, o teorema 4 permite
determinar a classe de cohomologia associada à diagonal de $X\times X$. Mediante
um raciocínio bem conhecido, deduz-se, por exemplo, uma fórmula de Lefschetz:

\result{TEOREMA}{5} Seja $f$ um endomorfismo de uma variedade projetiva não
singular $X$, tal que os pontos fixos de $f$ tenham multiplicidade 1.
Então o número desses pontos fixos é congruente, módulo a característica de
$k$, com a soma alternada dos traços dos endomorfismos dos $H^{i}(X)$ definidos
por $f$.

A restrição sobre $f$ que foi necessário impor deve-se às dificuldades mencionadas na
observação do n\textsuperscript{o} 4. Não obstante, observar-se-á que a fórmula de
Lefschetz continua verdadeira se $f$ é « homotópico » a um endomorfismo cujos
pontos fixos têm todos multiplicidade 1.

\fsection{8}{Generalização do teorema de dualidade}

Seja $X$ uma variedade algébrica não singular tal que todo feixe
algébrico coerente $F$ sobre $X$ seja isomorfo a um quociente de um feixe
algébrico coerente localmente livre (o que ocorre se $X$ é localmente
fechado em um espaço projetivo). Então todo feixe algébrico coerente $F$ sobre
$X$ admite uma resolução finita $L$ mediante feixes localmente livres, e para
duas resoluções desta classe sempre se pode encontrar uma terceira e um
homomorfismo desta nas duas primeiras compatível com as
aumentações. Analogamente, se $L$ é uma resolução finita localmente livre de
$F$, e se é dado um homomorfismo $F'\longrightarrow F$, então existe
uma resolução finita localmente livre $L'$ de $F'$ e um homomorfismo
$L'\longrightarrow L$ compatível com $F'\longrightarrow F$, que se pode
supor sobrejetivo se $F'\longrightarrow F$ o é. Isto permite definir, dados
dois inteiros $r,s\geq0$, dois multifuntores cohomológicos nos
argumentos $A_1,\ldots,A_r$; $B_1,\ldots,B_s$,

\src{190}
que variam na categoria dos feixes algébricos coerentes sobre $X$, e cujos
valores estão na categoria dos feixes algébricos coerentes sobre $X$,
resp. na categoria dos módulos sobre $H^{0}(X,\Osh_X)$, mediante as fórmulas:
\begin{equation}\tag{8.1}
\resizebox{\textwidth}{!}{$
\left\{
\begin{aligned}
\underline{T}_{r}^{s*}(A_1,\ldots,A_r;B_1,\ldots,B_s)
  &={}H^{*}\!\left(
    \Homsh_{\Osh_X}\!\left(
      L(A_1)\otimes\cdots\otimes L(A_r),
      L(B_1)\otimes\cdots\otimes L(B_s)
    \right)\right),\\[3pt]
T_{r}^{s*}(A_1,\ldots,A_r;B_1,\ldots,B_s)
  &={}\mathbf R^{*}\Gamma\!\left(
    \Homsh_{\Osh_X}\!\left(
      L(A_1)\otimes\cdots\otimes L(A_r),
      L(B_1)\otimes\cdots\otimes L(B_s)
    \right)\right).
\end{aligned}
\right.
$}
\end{equation}
Nesta fórmula, $L(F)$ designa uma resolução finita localmente livre do
feixe algébrico coerente $F$, e $\mathbf R^{*}\Gamma(K)$ designa
a hipercohomologia do espaço $X$ com respeito ao complexo de feixes $K$. Se
$r$ ou $s$ é nulo, substitui-se o produto tensorial dos $L(A_i)$, resp. dos
$L(B_j)$, por $\Osh_X$. Em particular, $\underline{T}_{0}^{0}$ e $T_{0}^{0}$
são funtores graduados de 0 argumentos, $\underline{T}_{0}^{0}$ se reduz ao
grau 0, onde é o feixe $\Osh_X$, enquanto que $T_{0}^{0}$ é idêntico a
$H^{*}(X,\Osh_X)$. O fato de que os segundos membros de (8.1) não dependam
das resoluções escolhidas é, em qualquer caso, evidente para a primeira linha
(pois a questão é então local), e para a segunda resulta das observações
gerais precedentes, tendo em conta a sequência espectral de
hipercohomologia do complexo de feixes
\[
  K=\Homsh_{\Osh_X}\bigl(L(A_1)\otimes\cdots,
      L(B_1)\otimes\cdots\bigr),
\]
que converge à hipercohomologia de $X$ com respeito a $K$ e cujo termo
inicial é $H^{p}(X,H^{q}(K))$, isto é,
\begin{equation}\tag{8.2}
  E_{2}^{p,q}
  =H^{p}\!\left(X,
    (\underline{T}_{r}^{s})^{(q)}
      (A_1,\ldots,A_r;B_1,\ldots,B_s)\right).
\end{equation}
Vê-se então que esta sequência espectral tampouco depende das resoluções
escolhidas; seu limite é o primeiro membro de (8.1). Definem-se
facilmente os operadores de conexão relativos aos diversos argumentos $A_i$, $B_j$,
observando que toda sequência exata
$0\longrightarrow F'\longrightarrow F\longrightarrow F''\longrightarrow0$
pode ser resolvida mediante uma sequência exata de complexos finitos localmente livres.

No sistema dos funtores $\underline{T}_{r}^{s*}$, resp.
$T_{r}^{s*}$, definem-se operações análogas às do cálculo tensorial, cuja
definição é imediata a partir das fórmulas definidoras (8.1). Assim, tem-se
uma composição (que generaliza a considerada no n\textsuperscript{o} 2):
\begin{equation}\tag{8.3}
\resizebox{\textwidth}{!}{$
\begin{aligned}
&T_{r}^{s*}(A_1,\ldots,A_r;B_1,\ldots,B_s)
\times T_{r'}^{s'*}(A'_1,\ldots,A'_{r'};B'_1,\ldots,B'_{s'})\\
&\qquad\longrightarrow
T_{r+r'}^{s+s'*}
(A_1,\ldots,A_r,A'_1,\ldots,A'_{r'};
 B_1,\ldots,B_s,B'_1,\ldots,B'_{s'}),
\end{aligned}
$}
\end{equation}
que satisfaz as propriedades evidentes de associatividade e de compatibilidade com
os homomorfismos funtoriais, os homomorfismos de conexão e as sequências
espectrais. Do mesmo modo, têm-se operações de simetria cuja explicitação se deixa ao
leitor. Tem-se, além disso, uma operação de contração cada vez que
um dos argumentos $A_i$

\src{191}
é igual a um dos argumentos $B_j$:
\begin{equation}\tag{8.4}
\resizebox{\textwidth}{!}{$
\begin{aligned}
&T_{r}^{s*}(A_1,\ldots,A_{i-1},C,A_{i+1},\ldots,A_r;
 B_1,\ldots,B_{j-1},C,B_{j+1},\ldots,B_s)\\
&\qquad\longrightarrow
T_{r-1}^{s-1*}(A_1,\ldots,\widehat{A_i},\ldots,A_r;
 B_1,\ldots,\widehat{B_j},\ldots,B_s).
\end{aligned}
$}
\end{equation}
Além disso, se um argumento $A_i$ é um feixe localmente livre, pode
ser suprimido, com a condição de substituir um dos $A_j$ ($j\neq i$) por
$A_j\otimes A_i$, ou um dos $B_k$ por $B_k\otimes A_i'$ (onde
$A_i'=\Homsh_{\Osh_X}(A_i,\Osh_X)$), e tem-se uma regra de cálculo análoga no
caso em que um dos argumentos $B_i$ é localmente livre. Em particular, sempre
se pode suprimir um argumento que seja idêntico a $\Osh_X$. Se todos os
argumentos são localmente livres, salvo no máximo um dos argumentos $B_i$, a regra
que acaba de ser enunciada dá um isomorfismo funtorial
\begin{equation}\tag{8.5}
  T_{r}^{s*}(A_1,\ldots,A_r;B_1,\ldots,B_s)
  =H^{*}\bigl(X,A_1'\otimes\cdots\otimes A_r'
    \otimes B_1\otimes\cdots\otimes B_s\bigr)
\end{equation}
(pois se reduz ao caso em que $r=0$, $s=1$, onde isto é imediato; também se
pode utilizar diretamente a sequência espectral de termo inicial (8.2)). Para os
$\underline{T}_{r}^{s}$ se definem operações correspondentes às anteriores.
As relações entre as diversas operações assim
introduzidas são as mesmas que para as operações análogas do cálculo
tensorial.

Seja $n$ a dimensão de $X$. Aplicando sucessivamente uma composição
tensorial (8.3) e contrações (8.4) sobre os argumentos repetidos, obtém-se
um acoplamento
\begin{equation}\tag{8.6}
\resizebox{\textwidth}{!}{$
  (T_{r}^{s})^{p}(A_1,\ldots;B_1,\ldots)
  \times
  (T_{s}^{r})^{n-p}(B_1,\ldots;A_1,\ldots,A_r\otimes\Omega_X^{n})
  \longrightarrow H^{n}(X,\Omega_X^{n}).
$}
\end{equation}

\result{TEOREMA}{6} Se $X$ é uma variedade projetiva não singular, então
os acoplamentos (8.6) são dualidades.

Isto resulta de maneira puramente formal do corolário do teorema 3. Com efeito,
deste corolário resulta facilmente que, se $K$ é um complexo de feixes
algébricos coerentes localmente livres, então a hipercohomologia de $X$ com
respeito a $K$ está em dualidade com a hipercohomologia de $X$ com respeito a
$K'\otimes\Omega_X^{n}$ mediante os acoplamentos naturais
\begin{equation}\tag{8.7}
  \mathbf R^{p}\Gamma(K)
  \times\mathbf R^{n-p}\Gamma(K'\otimes\Omega_X^{n})
  \longrightarrow\mathbf R^{n}\Gamma(\Omega_X^{n})
  =H^{n}(X,\Omega_X^{n}).
\end{equation}
Isto se vê utilizando a sequência espectral de termo inicial
$H^{p}(X,H^{q}(K))$\footnote{O texto impresso diz
$H^{p}(H^{q}(X,K))$; o termo inicial é o da sequência espectral (8.2),
isto é, $H^{p}(X,H^{q}(K))$.} e a sequência espectral análoga para
$K'\otimes\Omega_X^{n}$. Do resultado anterior se deduz o teorema 6
utilizando a definição (8.1).

\src{192}
\subsection*{OBSERVAÇÕES}

\noindent 1\textsuperscript{o} Para as definições anteriores ao teorema 6,
não era necessário que $X$ fosse não singular, pois não era
indispensável trabalhar apenas com resoluções finitas. Mas se $X$
é singular, já não se poderá afirmar, a priori, que os
$(\underline{T}_{r}^{s})^{p}(A_1,\ldots;B_1,\ldots)$ sejam feixes
coerentes, pois no complexo de feixes
\[
  \Homsh_{\Osh_X}\bigl(L(A_1)\otimes\cdots,L(B_1)\otimes\cdots\bigr)
\]
haverá uma infinidade de componentes de um grau total dado.

\medskip
\noindent 2\textsuperscript{o} Verifica-se facilmente que, nas fórmulas
(8.1), pode-se substituir um dos $L(B_i)$ por $B_i$. Tendo em conta a
proposição 3, isto demonstra que se tem
\begin{equation}\tag{8.8}
\left\{
\begin{aligned}
  \underline{T}_{1}^{1*}(A;B)&=\Extsh_{\Osh_X}^{*}(A,B),\\
  T_{1}^{1*}(A;B)&=\operatorname{Ext}_{\Osh_X}^{*}(X;A,B).
\end{aligned}
\right.
\end{equation}
Em particular, tomando em (8.6) $r=s=1$ e $A_1=\Osh_X$, recupera-se o
teorema 3. A fórmula (8.8) implica também
\[
  T_{0}^{1*}(B)=H^{*}(X,B),
  \qquad
  T_{1}^{0*}(A)=\operatorname{Ext}_{\Osh_X}^{*}(X;A,\Osh_X).
\]

\medskip
\noindent 3\textsuperscript{o} Vê-se na fórmula (8.1) que, em geral, os
funtores $\underline{T}_{r}^{s*}$ e $T_{r}^{s*}$ têm componentes de
graus positivos e negativos. Utilizando a observação 2, vê-se que, se a
dimensão de $X$ é $n$, então as componentes não nulas de
$\underline{T}_{r}^{s*}$ estão compreendidas entre $-(s-1)n$ e $rn$, se $s>0$,
entre $0$ e $rn$, se $s=0$, e as componentes não nulas de $T_{r}^{s*}$
estão compreendidas entre $-(s-1)n$ e $(r+1)n$, se $s>0$, entre $0$ e $(r+1)n$,
se $s=0$ (e inclusive, salvo erro, se $r>0$, entre $-(s-1)n$ e $rn$, resp. entre
$0$ e $rn$).

\section*{BIBLIOGRAFIA}

\noindent[1] Cartier (Pierre). --- Les groupes
$\operatorname{Ext}^{s}(A,B)$, Séminaire A. Grothendieck : Algèbre homologique,
t. 1, 1957, n\textsuperscript{o} 3.

\noindent[2] Grothendieck (Alexandre). --- Sur quelques points d'algèbre
homologique, \emph{Tohoku Math. J.}, t. 9, 1957, p. 119--183.

\noindent[3] Serre (Jean-Pierre). --- Faisceaux algébriques cohérents,
\emph{Annals of Math.}, t. 61, 1955, p. 197--278.

\noindent[4] Serre (Jean-Pierre). --- Sur la dimension homologique des anneaux
et des modules noethériens, \emph{Proc. Intern. Symp. on Alg. Number Theory}
[1955, Tokyo et Nikko]. --- Tokyo, Science Council of Japan, 1956 ;
p. 175--189.

\clearpage
\src{193}
\section*{ADENDA}

As dificuldades assinaladas na observação da página 13 estão atualmente
completamente resolvidas, graças a uma extensão dos teoremas de dualidade a
variedades quaisquer (com singularidades arbitrárias). Para esta teoria, que pode
ser formulada no quadro geral da « teoria dos esquemas », remetemos
aos « Éléments de Géométrie algébrique » de J. Dieudonné e A. Grothendieck,
atualmente em preparação. Encontrar-se-ão algumas indicações em:

\begin{quote}
Grothendieck (A.). --- Cohomology theory of algebraic varieties,
\emph{Congrès int. des Mathématiciens}, 1958, Edinburgh (de próxima publicação).
\end{quote}

\begin{flushright}
[Abril de 1959]
\end{flushright}

\begin{center}\rule{36mm}{0.5pt}\end{center}

\par

\def\FgaUnitKey{e149}\clearpage
\chapter*{Errata da exposição n\textsuperscript{o} 149}
\addcontentsline{toc}{chapter}{Errata da exposição n\textsuperscript{o} 149}
\setcounter{footnote}{0}
\markboth{Errata da exposição 149}{Errata da exposição 149}
\sid{e149}
\src{298}

\noindent\textbf{1. Teorema de dualidade para os feixes algébricos
coerentes.}

\begin{center}
[Seminário Bourbaki, t. 9, 1956/57, n\textsuperscript{o} 149, 25 p.]
\end{center}

\noindent\textbf{Observação sobre a página 14.} --- Como assinalei em minha
conferência no Congresso Internacional de Matemáticos de 1958,\footnote{
Grothendieck (Alexander). --- The cohomology theory of abstract algebraic
varieties, \emph{Proceedings of the International Congress of Mathematicians}
[1958, Edinburgh], p. 103--118. --- Cambridge, at the University Press, 1960.}
as questões levantadas aqui estão agora completamente resolvidas.

O leitor encontrará mais informações sobre a dualidade dos feixes
coerentes em \emph{loco citato}\footnotemark[1], p. 112--115, em EGA III,
2\textsuperscript{a} parte (em preparação) e, por enquanto, em SGA 1962.
Um tratamento mais sistemático encontrar-se-á em um capítulo posterior de EGA
(capítulo IX no plano previsto).
% Explicit module paragraph boundary; no prose change.
\par

\par

\def\FgaUnitKey{182}\unit{Exposição 182 — Geometria formal e geometria algébrica}
\sid{182}
\markboth{Exposição 182}{Exposição 182}
\src{193}
\seminar{(Maio de 1959)}
\paperhead{GEOMETRIA FORMAL E GEOMETRIA ALGÉBRICA}%
  {por Alexander GROTHENDIECK}

\fsection{1}{Esquemas}

Sabe-se que um espaço algébrico afim definido sobre um corpo $k$ é
essencialmente determinado por sua álgebra afim $A$ (anel das funções
regulares definidas sobre $k$), pois os morfismos de espaços algébricos
$X\longrightarrow Y$ correspondem biunívocamente aos homomorfismos de
$k$-álgebras $A(Y)\longrightarrow A(X)$. O álgebra afim correspondente a um
espaço algébrico é uma $k$-álgebra de tipo finito e, do ponto de vista
«clássico», não tem elementos nilpotentes; reciprocamente, toda álgebra desse
tipo obtém-se como álgebra afim de um espaço algébrico definido sobre $k$.
Dispõe-se assim de um conhecido dicionário que permite interpretar, em termos
de álgebra comutativa, as situações relativas a espaços algébricos afins.
Há muito se observou que, desse modo, obtêm-se enunciados mais gerais, pois em
geral já não é necessário supor que os anéis em questão sejam do tipo recém
considerado, e a hipótese noetheriana costuma bastar. Em particular, haja ou
não um corpo de base dado, não há razão para excluir o caso em que esses anéis
contenham elementos nilpotentes. Até agora, os geômetras recusaram-se a levar
em conta essas indicações e obstinaram-se em limitar-se a considerar álgebras
afins sem elementos nilpotentes, isto é, espaços algébricos cujos feixes
estruturais não têm elementos nilpotentes (e, na maioria das vezes, até mesmo
espaços algébricos «absolutamente irredutíveis»). O conferencista considera que
essa atitude constituiu um sério obstáculo ao desenvolvimento de métodos
verdadeiramente naturais em geometria algébrica.

Seja $A$ um anel comutativo. É bem sabido que o conjunto
$X=\operatorname{Spec}(A)$ dos ideais primos de $A$ está provido de uma
topologia natural, a «topologia de Zariski» ou topologia espectral. Por outra
parte, há um feixe de anéis comutativos $\Osh_X$ sobre $X$, cujo talo em
$\mathfrak p\in X$ é o anel localizado $A_{\mathfrak p}$ e cujo anel de
seções identifica-se com $A$. Assim, $X$ torna-se um \emph{espaço
anelado}, chamado \emph{espectro primo} de $A$. Um homomorfismo de anéis
$f:A\longrightarrow B$ define um morfismo de espaços anelados
$f':\operatorname{Spec}(B)\longrightarrow\operatorname{Spec}(A)$; a
aplicação subjacente de conjuntos não é outra senão
$\mathfrak p\longmapsto f^{-1}(\mathfrak p)$. Os morfismos

\src{194}
de espaços anelados de $\operatorname{Spec}(B)$ em
$\operatorname{Spec}(A)$ obtidos desse modo são exatamente aqueles para
os quais os homomorfismos $\Osh_x\longrightarrow\Osh_y$ ($x=f'(y)$) são
locais (isto é, a imagem inversa do ideal maximal é o ideal maximal).

Chama-se \emph{esquema afim} um espaço anelado isomorfo a algum
$\operatorname{Spec}(A)$, e \emph{preesquema} um espaço anelado localmente
afim, isto é, tal que cada ponto tem uma vizinhança aberta que, com a
estrutura induzida, é um esquema afim. Os \emph{morfismos} de preesquemas se
definem de maneira evidente; localmente correspondem a homomorfismos de anéis.

Quando se fixa um preesquema $S$ e se consideram morfismos de preesquemas
$X\longrightarrow S$, o espaço $S$ desempenha o papel de um corpo ou de um
anel de base (ou, melhor, de um espaço de base numa fibração). Diz-se
então que $X$ é um $S$-preesquema; se $S=\operatorname{Spec}(A)$, isto
significa também que $\Osh_X$ é um feixe de $A$-álgebras. Assim, todo preesquema
pode ser considerado de maneira única como um $\mathbf Z$-preesquema. Naturalmente,
os $S$-preesquemas formam uma categoria; além disso, demonstra-se que nessa
categoria sempre existe o produto de dois objetos $X,Y$, denotado por
$X\times_S Y$. Essa noção de produto permite definir a mudança de base de um
$S$-preesquema correspondente a um morfismo $S'\longrightarrow S$: com efeito,
$X\times_S S'$ pode ser considerado como um $S'$-preesquema.

Diz-se que $X$ é \emph{separado} sobre $S$ se a diagonal de
$X\times_S X$ é fechada. Chama-se \emph{esquema} um preesquema separado
sobre $\mathbf Z$; ele é então separado sobre qualquer base. Para simplificar,
doravante falaremos apenas de esquemas e, além disso, os suporemos
\emph{noetherianos}, isto é, uniões finitas de abertos afins que são
espectros de anéis noetherianos. Diz-se que $X$ é \emph{de tipo finito
sobre} $S$ se, para todo aberto afim $U$ de $S$, sua imagem inversa em $X$ é
uma união finita de abertos afins cujos anéis são álgebras de tipo finito
sobre o anel de $U$. São esses $S$-esquemas os que se prestam a um estudo
propriamente geométrico. Em particular, para todo $s\in S$, a fibra
$f^{-1}(s)$ de $X$ sobre $s$ é um esquema algébrico sobre o corpo residual
$\kappa(s)$ do anel local $\Osh_s$ de $s$ em $S$. Assim, em certa medida, $X$
pode ser considerado como uma família de «espaços algébricos» $f^{-1}(s)$, onde
o parâmetro $s$ percorre $S$ (isto é, do ponto de vista local, o
conjunto dos ideais primos de um anel dado). Naturalmente, os
$\kappa(s)$ podem ter características distintas. Se
$S=\operatorname{Spec}(k)$, onde $k$ é um corpo, recupera-se essencialmente a
noção usual de «espaço algébrico», com a única diferença de que agora o
feixe estrutural pode ter elementos nilpotentes.

\src{195}
Inspirando-se em noções bem conhecidas, define-se a noção de
\emph{morfismo projetivo} e, mais geralmente, a de \emph{morfismo próprio}.
Tal morfismo é de tipo finito; além disso, transforma partes fechadas em partes
fechadas e conserva essa propriedade sob qualquer mudança de base.

Sendo $X$ um esquema (noetheriano, como sempre), o feixe $\Osh_X$ é um
\emph{feixe coerente de anéis} no sentido de [2]. Portanto, os feixes
coerentes de módulos sobre $X$ são também os feixes que localmente são
isomorfos ao conúcleo de um homomorfismo
$\Osh_X^{m}\longrightarrow\Osh_X^{n}$.

\fsection{2}{Esquemas formais}

Sejam $X$ um esquema e $X'$ uma parte fechada de $X$. Então existe um
subfeixe coerente $J$ de $\Osh_X$ tal que
$X'=\operatorname{Supp}(\Osh_X/J)$ (e inclusive existe um máximo). Provido de
$\Osh_X/J$, o espaço $X'$ torna-se um esquema, denotado por $X_0$;
tal esquema se chama «subesquema fechado de $X$». Para todo $n$ também se
pode considerar $X'$ provido de $\Osh_X/J^{n+1}$, denotado por $X_n$;
é um subpreesquema fechado de $X$ cujo conjunto subjacente continua sendo $X'$,
mas tem outro feixe estrutural, a saber,
$\Osh_{X_n}=\Osh_X/J^{n+1}$. Evidentemente, os $\Osh_{X_n}$ formam um sistema
projetivo de feixes de anéis sobre $X'$, cujo limite projetivo
$\overline{\Osh}_X$ chama-se \emph{completamento formal de $\Osh_X$ ao longo
de $X'$}. Provido desse feixe de anéis, $X'$ chama-se \emph{completamento
formal de $X$ ao longo de $X'$}; é, pois, um espaço anelado, mas em
geral não um esquema. Para todo feixe coerente $F$ sobre $X$, pode-se considerar
analogamente o completamento formal $\overline F=\varprojlim_n F_n$ de $F$ ao
longo de $X'$ (onde $F_n=F\otimes_{\Osh_X}\Osh_X/J^{n+1}$); é um feixe de
módulos sobre $\overline X$. Suas seções chamam-se «seções formais de $F$
ao longo de $X'$» e identificam-se com os elementos de
$\varprojlim_n\Gamma(X',F_n)$. Para $F=\Osh_X$ obtêm-se as «funções
holomorfas» de $X$ ao longo de $X'$ no sentido de Zariski (terminologia que
não seguiremos por suas interferências com a terminologia clássica).

Chama-se \emph{esquema formal} (subentende-se: noetheriano) um espaço
topológico $\mathfrak X$, provido de um feixe de anéis topológicos
$\Osh_{\mathfrak X}$ que satisfaz a condição seguinte: há um isomorfismo
de feixes de anéis topológicos
$\Osh_{\mathfrak X}=\varprojlim_n\Osh_n$, onde os $\Osh_n$ formam um sistema
projetivo de feixes de anéis sobre $\mathfrak X$, cada um dos quais faz
de $\mathfrak X$ um esquema $\mathfrak X_n$, e tal que, para $m\geq n$, o
homomorfismo $\Osh_m\longrightarrow\Osh_n$ é sobrejetivo e tem por núcleo
$J_m^{n+1}$, onde $J_m$ é

\src{196}
o núcleo de $\Osh_m\longrightarrow\Osh_0$. Demonstra-se então que
$\Osh_{\mathfrak X}$ é um \emph{feixe coerente de anéis locais
noetherianos}.

Em virtude das definições, um completamento formal $\overline X$ como o
anterior é um esquema formal e, reciprocamente, todo esquema formal é
\emph{localmente} desse tipo. De fato, dar um esquema formal \emph{afim} (isto
é, tal que $\mathfrak X_0$ seja afim, o que implica que todos os
$\mathfrak X_n$ também o sejam) equivale a dar um anel topológico noetheriano,
separado e completo para a topologia $J$-ádica.

As definições usuais de morfismo, morfismo de tipo finito, morfismo
próprio, etc., para os esquemas ordinários estendem-se sem dificuldade aos
esquemas formais.

\fsection{3}{Os três teoremas fundamentais}

Seja $f:X\longrightarrow Y$ um morfismo próprio de esquemas (noetherianos, como
sempre), seja $Y'$ uma parte fechada de $Y$ e seja $X'$ sua imagem inversa em
$X$; consideremos os completamentos formais correspondentes $\overline Y$ e
$\overline X$. Então $f$ induz um morfismo
$\overline f:\overline X\longrightarrow\overline Y$ de esquemas formais, que
além disso é próprio. Se $F$ é um feixe coerente sobre $X$, então $\overline F$
é um feixe coerente sobre $\overline X$. No teorema 1, esquecem-se $X,Y,F$ e
consideram-se apenas o morfismo próprio $\overline f$ de esquemas formais e o
feixe coerente $\overline F$ sobre $\overline X$. (Não obstante, o conferencista
só escreveu uma demonstração completa para o caso em que se parte de
$X,Y,f,F$.)

\result{TEOREMA}{1 (teorema de finitude)}
\begin{enumerate}
  \item[\textit{i.}] Os $R^{q}\overline f_{*}(\overline F)$ são feixes
    coerentes sobre $\overline Y$.
  \item[\textit{ii.}] Os homomorfismos naturais
    \[
      R^{q}\overline f_{*}(\overline F)
      \longrightarrow \varprojlim_n R^{q}f_{n*}(F_n)
    \]
    são isomorfismos.
\end{enumerate}

Neste enunciado supõe-se escolhido um subfeixe coerente $J$ de $\Osh_Y$ que
define $Y'$; por imagem inversa obtém-se um subfeixe coerente de $\Osh_X$ que
define $X'$ e, por conseguinte, se definem $F_n,X_n,Y_n$ e
$f_n:X_n\longrightarrow Y_n$ como no n\textsuperscript{o} 2. Se partirmos de
um morfismo próprio qualquer entre dois esquemas formais, são evidentes as
pequenas mudanças que devem ser feitas ao explicitar essas notações.

O teorema 1 referia-se apenas à «cohomologia formal». O teorema seguinte

\src{197}
a relaciona com a «cohomologia algébrica» e é análogo a um conhecido teorema
de Serre [4] sobre a comparação entre cohomologia algébrica e cohomologia
analítica.

\result{TEOREMA}{2 (Primeiro teorema de comparação)} Os
$R^{q}f_{*}(F)$ são feixes coerentes sobre $Y$ (o que é um caso particular
do teorema 1), e os homomorfismos naturais
\[
  \overline{R^{q}f_{*}(F)}
  \longrightarrow \varprojlim_n R^{q}f_{n*}(F_n)
\]
são isomorfismos.

\result{COROLÁRIO}{1} Têm-se isomorfismos canônicos:
\[
  \overline{R^{q}f_{*}(F)}=R^{q}\overline f_{*}(\overline F).
\]

Para $q=0$, esse corolário é uma generalização do «teorema fundamental das
funções holomorfas» de Zariski, do qual deduziremos uma generalização do
«teorema de conexão» de Zariski. Observe-se além disso que, enquanto o teorema 1
(\textit{ii}) é trivial para $q=0$, de modo algum ocorre o mesmo com o
teorema 2 nem com sua formulação equivalente (corolário 1). Com efeito, a
demonstração procede por indução descendente sobre $q$ (é trivial para $q$
grande, pois então ambos os membros são nulos), de modo que o caso $q=0$
aparece como a última etapa da indução e, por assim dizer, como o caso «mais
difícil».

\result{COROLÁRIO}{2} Suponhamos $Y=\operatorname{Spec}(A)$ e que $Y'$ está
definido pelo ideal $J$ de $A$. Então, para todo feixe coerente $F$ sobre
$X$, os $H^{q}(X,F)$ são $A$-módulos de tipo finito cujos completamentos
$J$-ádicos são os $H^{q}(\overline X,\overline F)$.

Aplicando, por fim, este corolário a
$H=\Homsh_{\Osh_X}(F,G)$, obtém-se:

\result{COROLÁRIO}{3} Suponhamos $Y=\operatorname{Spec}(A)$ e que $Y'$ está
definido pelo ideal $J$ de $A$. Sejam $F,G$ dois feixes coerentes sobre $X$;
então $\operatorname{Hom}(F,G)$ é um módulo de tipo finito sobre $A$, cujo
completamento $J$-ádico identifica-se com
$\operatorname{Hom}(\overline F,\overline G)$.

Naturalmente, a aplicação natural
$\operatorname{Hom}(F,G)\longrightarrow
\operatorname{Hom}(\overline F,\overline G)$ é a que associa a um
homomorfismo $u:F\longrightarrow G$ sua extensão «por continuidade»
$\overline u:\overline F\longrightarrow\overline G$ (com o que
$\overline F$ se torna um funtor de $F$).

Suponhamos agora que $A$ é separado e completo para sua topologia $J$-ádica;
então os corolários 2 e 3 precedentes dão:
\[
  H^{q}(X,F)=H^{q}(\overline X,\overline F),
  \qquad
  \operatorname{Hom}(F,G)
  =\operatorname{Hom}(\overline F,\overline G).
\]

\src{198}
Esta última identidade mostra que a categoria dos feixes coerentes sobre
$X$ se identifica com uma \emph{subcategoria} (cujos morfismos são os
morfismos induzidos) da categoria dos feixes coerentes sobre
$\overline X$. De fato, tem-se inclusive:

\result{TEOREMA}{3} Para que um feixe de módulos sobre $\overline X$ seja
coerente, é necessário e suficiente que seja isomorfo a um feixe da forma
$\overline F$, onde $F$ é um feixe coerente sobre $X$ (determinado salvo
isomorfismo canônico em virtude do teorema 2, corolário 3). [Lembremos que
agora $Y=\operatorname{Spec}(A)$, sendo $A$ um anel topológico noetheriano,
separado e completo para a topologia $J$-ádica.]

\result{COROLÁRIO}{1} Os subesquemas fechados de $X$ correspondem
biunívocamente aos subesquemas formais fechados de $\overline X$.

Com efeito, correspondem aos subfeixes coerentes de $\Osh_X$ e,
respectivamente, de $\overline{\Osh}_X$. Considerando os grafos de morfismos
como subesquemas fechados, deduz-se do corolário 1:

\result{COROLÁRIO}{2} Sejam $X,Z$ dois esquemas próprios sobre $A$ (anel
noetheriano, separado e completo para a topologia $J$-ádica). Então a
aplicação $g\longmapsto\overline g$ define uma correspondência biunívoca entre
os $Y$-morfismos de $X$ em $Z$ e os $\overline Y$-morfismos de $\overline X$
em $\overline Z$.

Em outros termos, os esquemas algébricos próprios sobre $A$ aparecem como uma
\emph{subcategoria} (cujos morfismos são os morfismos induzidos) da
categoria dos esquemas formais próprios sobre $\overline Y$. No entanto,
deve-se advertir que existem esquemas formais próprios sobre $\overline Y$ que
não são «algebrizáveis», isto é, isomorfos a algum $\overline X$ com $X$
próprio sobre $A$ (do mesmo modo que existem variedades analíticas complexas
compactas que não provêm de variedades algébricas definidas sobre o corpo
dos complexos). Tais esquemas formais aparecem de maneira natural na
«teoria dos módulos». Assinalemos, não obstante, um caso particular interessante em
que um esquema formal é algebrizável:

\result{TEOREMA}{4} Seja $A$ um anel noetheriano local completo de corpo
residual $k$, e seja $\mathfrak X$ um esquema formal próprio sobre $A$ (provido
de sua topologia $\mathfrak m(A)$-ádica). Supõe-se:
\begin{enumerate}
  \item[\textit{i.}] que os anéis locais de $\Osh_{\mathfrak X}$ são
    módulos planos sobre $A$ ou, o que equivale ao mesmo, se se dota a
    $\Osh_{\mathfrak X}$ e a $A$ da filtração definida pelas potências do
    ideal maximal de $A$, para os graduados associados tem-se a relação
    \[
      \operatorname{gr}(\Osh_{\mathfrak X})
      \simeq
      \operatorname{gr}^{0}(\Osh_{\mathfrak X})
      \otimes_k\operatorname{gr}(A).
    \]
  \item[\textit{ii.}] Considerando
    $\mathfrak X_0=\mathfrak X\otimes_A k$ como um esquema algébrico sobre
    $k$, tem-se
    \[
      H^2(\mathfrak X_0,\Osh_{\mathfrak X_0})=0.
    \]
  \item[\textit{iii.}] $\mathfrak X_0$ é projetivo.
\end{enumerate}
Sob essas condições, $\mathfrak X$ é algebrizável e, mais precisamente, é
isomorfo a $\overline X$, onde $X$ é um $A$-esquema projetivo.

As condições \textit{(ii)} e \textit{(iii)} são satisfeitas, em particular,
se $\mathfrak X_0$ é uma curva lisa sobre $k$, e o teorema 4 pode ser aplicado
em particular na «teoria de módulos» para as curvas de gênero dado\ldots{}
Indiquemos como se demonstra o teorema 4: demonstra-se (cf. a proposição 3
mais adiante) que \textit{(i)} e \textit{(ii)} implicam que todo feixe coerente
sobre $\mathfrak X_0$, localmente isomorfo ao feixe fundamental, obtém-se por
redução a partir de um feixe da mesma natureza sobre $\mathfrak X$. Parte-se
então de um feixe «amplo» sobre $\mathfrak X_0$ (existe um por
\textit{(iii)}), levanta-se esse feixe a um feixe invertível sobre $\mathfrak X$ e,
utilizando o teorema 1, se demonstra que um múltiplo deste último define uma
imersão de $\mathfrak X$ no completamento formal de um esquema
$\mathbf P_A^r$ («projetivo tipo» de dimensão $r$ sobre $A$).

Para a demonstração dos teoremas 1 a 3, veja-se [1].

\src{199}
\fsection{4}{Aplicação ao teorema de conexão e ao «Main theorem» de
Zariski}

Seja $f:X\longrightarrow Y$ um morfismo próprio de esquemas. Pelo teorema de
finitude, $f_*(\Osh_X)=\mathcal A$ é um feixe coerente sobre $Y$; além disso, é um
feixe de álgebras comutativas, pelo que corresponde a um $Y$-esquema
$g:Y'\longrightarrow Y$ finito sobre $Y$ (definido pela condição de ser afim
sobre $Y$, isto é, que a imagem inversa de um aberto afim seja afim, e por
$g_*(\Osh_{Y'})=\mathcal A$). É imediato que $f$ se fatora então
canonicamente como $f=gf'$, onde $f':X\longrightarrow Y'$ é um morfismo de
$X$ em $Y'$ que desta vez satisfaz $f'_*(\Osh_X)=\Osh_{Y'}$. Essa fatoração
de $f$ chama-se \emph{fatoração de Stein} de $f$. Aplicando o primeiro
teorema de comparação, corolário 1, a $f'$ e à parte de $Y'$ reduzida a um
ponto $y'$, obtém-se que $f'^{-1}(y')=X'$ é conexo (de outro modo, as
seções formais de $X$ ao longo de $X'$ não formariam um anel local,
enquanto o completamento de
$f'_*(\Osh_X)_{y'}=\Osh_{Y',y'}$ é local!). Demonstramos:

\result{TEOREMA}{5 («Teorema de conexão» de Zariski)} Seja
$f:X\longrightarrow Y$ um morfismo próprio. Então $f$ se fatora de maneira
única (salvo isomorfismo) como $f=gf'$, onde $g:Y'\longrightarrow Y$ é finito
e $f':X\longrightarrow Y'$ satisfaz $f'_*(\Osh_X)=\Osh_{Y'}$ (de onde
$g_*(\Osh_{Y'})=f_*(\Osh_X)$). As fibras de $f'$ são conexas; isto é, o
conjunto das componentes conexas de uma fibra $f^{-1}(y)$ de $f$ está em
correspondência biunívoca com o conjunto dos pontos de $Y'$ sobre $y$, isto
é, com o conjunto dos ideais maximais de $f_*(\Osh_X)_y$.

Daqui se deduzem imediatamente as variantes usuais do teorema de
conexão. Assinalemos aqui apenas o

\src{200}
\result{COROLÁRIO}{1} Para que um ponto $x$ de $X$ esteja isolado em sua fibra
$f^{-1}(y)$, é necessário e suficiente que a fibra $f'^{-1}(y')$ (onde
$y'=f'(x)$) se reduza a $x$, ou que $f'$ induza um isomorfismo de uma
vizinhança de $x$ sobre uma vizinhança de $y'$. O conjunto desses pontos é um
aberto $U$, e $f'$ induz um isomorfismo de $U$ sobre um aberto de $Y'$.

Para demonstrar que $f'$ é um isomorfismo local em $x$, observa-se que $f'$
induz um isomorfismo $\Osh_{y'}\longrightarrow\Osh_x$, como se vê graças a
$f'_*(\Osh_X)=\Osh_{Y'}$ e ao fato de que (sendo $f'$ uma aplicação fechada
cuja fibra em $y'$ se reduz a $x$) os $f'^{-1}(V)$ percorrem um sistema
fundamental de vizinhanças de $x$ quando $V$ recorre um sistema fundamental de
vizinhanças de $y'$. Daqui se deduz imediatamente o resultado seguinte, devido a
Chevalley no caso «geométrico».

\result{COROLÁRIO}{2} Para que $f$ seja um morfismo finito, é necessário e
suficiente que seja próprio e que suas fibras sejam finitas.

Se assim ocorre, $f'$ é efetivamente um isomorfismo pelo que precede.

Seja $f:X\longrightarrow Y$ um morfismo não necessariamente próprio, mas
suponhamos que $X$ está contido como aberto em um $Y$-esquema próprio
$\widetilde f:\widetilde X\longrightarrow Y$ (o que ocorre, em particular,
se $f$ é quase-projetivo). Aplicando o corolário 1, vê-se que $\widetilde f'$
induz um isomorfismo do conjunto $U$ dos pontos de $X$ isolados em sua fibra
sobre um aberto de $Y'$ (e que $U$ é, com efeito, aberto). Deduz-se a
seguinte variante global do «Main Theorem» de Zariski.

\result{TEOREMA}{6 («Main Theorem» de Zariski)} Seja
$f:X\longrightarrow Y$ um morfismo de tipo finito. Então o conjunto $U$ dos
pontos de $X$ que estão isolados em sua fibra é aberto e, se $f$ é
separado,\footnote{A errata de 1962 substitui aqui a hipótese impressa
«$f$ quase-projetivo» pela hipótese mais fraca «$f$ separado», em virtude de
SGA VIII, 6.2: todo morfismo quase-finito e separado é quase-projetivo.} $U$ é
$Y$-isomorfo a uma parte aberta de um esquema $Y'$ finito sobre $Y$.

Como um morfismo de tipo finito é localmente afim e, \textit{a fortiori},
localmente quase-projetivo, do teorema 6 deduzem-se imediatamente as variantes
locais usuais do Main Theorem.

\fsection{5}{Aplicação ao estudo cohomológico dos morfismos próprios e
planos}

Seja $f:X\longrightarrow Y$ um morfismo próprio, e seja $F$ um feixe coerente sobre
$X$; supõe-se que $F$ é $Y$-plano, isto é, os $F_x$ são módulos planos sobre os
anéis $\Osh_y$ ($y=f(x)$). Isto significa também que, para todo $y\in Y$, se
filtrarmos $F$ ao longo da fibra $f^{-1}(y)$ mediante os
$\mathfrak m_y^nF$ (sendo $\mathfrak m_y$ o ideal maximal de $\Osh_y$), o
graduado associado é isomorfo a
$(F/\mathfrak m_yF)\otimes_{\kappa(y)}\operatorname{gr}(\Osh_y)$; em outros
termos, tem-se
\[
  \mathfrak m_y^nF/\mathfrak m_y^{n+1}F
  =F_y\otimes_{\kappa(y)}
  (\mathfrak m_y^n/\mathfrak m_y^{n+1})
\]
para todo inteiro $n$, onde $X_y$ denota a fibra $f^{-1}(y)$ considerada como
esquema próprio sobre o corpo residual $\kappa(y)$ de $y$, e $F_y$ o feixe
$F/\mathfrak m_yF$ induzido por $F$ sobre $X_y$. Tendo em conta esse
isomorfismo e o teorema 2, obtêm-se cotas e, às vezes, cálculos dos
$R^qf_*(F)$ em uma vizinhança de $y$, conhecendo a cohomologia de $X_y$ com
coeficientes em $F_y$. O teorema 2 assume aqui a forma
\[
  \widehat{R^qf_*(F)_y}
  =\varprojlim_n H^q(X_y,F/\mathfrak m_y^nF).
\]
Assinalaremos apenas a consequência seguinte:

\src{201}
\result{PROPOSIÇÃO}{1} Sejam $f:X\longrightarrow Y$ um morfismo próprio e $F$
um feixe sobre $X$, coerente e $Y$-plano. Seja $y\in Y$, seja $q$ um inteiro, e
suponhamos que $H^q(X_y,F_y)=0$. Então $R^qf_*(F)$ é nulo em uma vizinhança de
$y$ e, para todo $n$, o homomorfismo natural
\[
  R^{q-1}f_*(F)_y\longrightarrow
  H^{q-1}(X_y,F/\mathfrak m_y^nF)
\]
é sobrejetivo.

Em particular, se $f$ é um morfismo plano (isto é, $\Osh_X$ é $Y$-plano),
então todo feixe coerente localmente livre $F$ sobre $X$ é $Y$-plano. Sejam
$F,G$ dois feixes desse tipo; aplicando a proposição 1 a
$\Homsh_{\Osh_X}(F,G)$ e a $q=1$, obtém-se:

\result{TEOREMA}{7} Seja $f$ um morfismo próprio e plano, sejam $F,G$ dois feixes
coerentes localmente livres sobre $X$, seja $y\in Y$ e suponhamos que
\[
  H^1\bigl(X_y,\Homsh_{\Osh_{X_y}}(F_y,G_y)\bigr)=0,
\]
então todo homomorfismo $u_0:F_y\longrightarrow G_y$ é induzido por um
homomorfismo
\[
  u:F|V\longrightarrow G|V,
  \qquad V=f^{-1}(U),
\]
onde $U$ é uma vizinhança de $y$.

\result{COROLÁRIO}{} Se $u_0$ é um isomorfismo (resp. um monomorfismo, um
epimorfismo), então o mesmo ocorre com $u$ para $U$ suficientemente
pequeno.

Em particular:

\src{202}
\result{COROLÁRIO}{2} Seja $E_0$ um feixe coerente localmente livre sobre $X_y$
tal que
\[
  H^1\bigl(X_y;\Homsh_{\Osh_{X_y}}(E_0,E_0)\bigr)=0,
\]
então dois feixes localmente livres cujas restrições a $X_y$ sejam isomorfas
a $E_0$ são isomorfos em uma vizinhança de $X_y$.

Assim:

\result{COROLÁRIO}{3} Suponhamos $H^1(X_y,\Osh_{X_y})=0$. Então dois feixes
invertíveis sobre $X$ (isto é, localmente isomorfos a $\Osh_X$) cujas
restrições a $X_y$ sejam isomorfas são isomorfos em uma vizinhança de $X_y$.\footnote{O texto impresso não especifica a vizinhança. Restabelece-se aqui o alcance local do corolário 2, também utilizado na demonstração da proposição 2.}

Deduz-se:

\result{PROPOSIÇÃO}{2} Sejam $Y$ um esquema conexo e $E$ um feixe coerente localmente livre de posto ao menos $2$ sobre $Y$; consideremos o fibrado de espaços projetivos
$X=\mathbf P(E)$ associado a $E$, provido de seu conhecido feixe invertível
$\Osh_X(1)$. Então todo feixe invertível $L$ sobre $X$ é isomorfo a um feixe da
forma
\[
  f^*(L')\otimes\Osh_X(n),
\]
onde $L'$ é um feixe invertível sobre $Y$ e $n$ um inteiro. Este último está
determinado de maneira única, e $L'$ está determinado salvo isomorfismo.

O corolário 3 precedente demonstra que $L$ é isomorfo a algum $\Osh_X(n)$ em
uma vizinhança de cada fibra. O resto é quase formal.

A proposição 2 permite determinar os $Y$-morfismos de $X=\mathbf P(E)$ em
outro fibrado projetivo. Em particular, obtém-se:

\result{COROLÁRIO}{} Seja $u$ um automorfismo de $X=\mathbf P(E)$. Então
existem um feixe invertível $L'$ sobre $Y$ e um isomorfismo $v$ de $E$ sobre
$E\otimes L'$ tais que $u$ é o isomorfismo correspondente
\[
  \mathbf P(E)\xrightarrow{\sim}\mathbf P(E\otimes L')=\mathbf P(E);
\]
o par $(v,L')$ está determinado salvo isomorfismo.

Seja $\Gamma$ o conjunto das classes de fibrados invertíveis $L'$ sobre $Y$
tais que $E\otimes L'$ seja isomorfo a $E$. Seus elementos são de torção, pois,
se $n$ é o posto de $E$, vê-se (tomando as potências exteriores $n$-ésimas)
que se deve ter $L'^{\otimes n}\simeq\Osh_Y$. O corolário pode ser expresso
então dizendo que se tem uma sequência exata de grupos
\[
  e\longrightarrow
  \operatorname{Aut}(E)/\Gamma(Y,\Osh_Y^*)
  \longrightarrow \operatorname{Aut}_Y(X)
  \longrightarrow \Gamma\longrightarrow e
\]
(deduzida, de resto, da sequência exata de cohomologia obtida a partir
da seguinte sequência exata de feixes de grupos
\[
  e\longrightarrow \Osh_Y^*
  \longrightarrow \underline{\operatorname{Aut}}(E)
  \longrightarrow \underline{\operatorname{Aut}}_Y(X)
  \longrightarrow e,
\]
onde $\Osh_Y^*$ é o feixe das «unidades» de $\Osh_Y$, identificado com o
centro de $\underline{\operatorname{Aut}}(E)$).

\src{203}
\fsection{6}{Aplicação a teoremas de existência e unicidade para feixes e
esquemas sobre um anel completo para a topologia \texorpdfstring{$J$}{J}-ádica}

O teorema 7 dava um resultado de unicidade para feixes coerentes localmente
livres utilizando os teoremas 1 e 2. Utilizando o teorema 3, deduzimos agora
teoremas de existência de feixes, de morfismos de esquemas ou de esquemas.
Doravante, $A$ denota um anel noetheriano local, separado e completo. O método
geral consiste sempre em efetuar construções formais, o que consiste
essencialmente em fazer geometria algébrica sobre um anel artiniano e extrair
delas conclusões de natureza «algébrica» mediante os três teoremas
fundamentais.

\result{PROPOSIÇÃO}{3} Seja $\mathfrak X$ um esquema formal próprio e plano
sobre $A$, e seja $F_0$ um feixe localmente livre sobre $X_0$ tal que
\[
  H^2\bigl(X_0,\Homsh_{\Osh_{X_0}}(F_0,F_0)\bigr)=0.
\]
Então existe um feixe localmente livre $F$ sobre $\mathfrak X$ que induz
sobre $X_0$ um feixe isomorfo a $F_0$. (Além disso, este $F$ é único salvo
isomorfismo se
$H^1(X_0,\Homsh_{\Osh_{X_0}}(F_0,F_0))=0$.)

Constroem-se sucessivamente feixes localmente livres $F_n$ sobre os $X_n$, cada
um dos quais induz o anterior. A construção de $F_n$ encontra uma
obstrução em
\[
  H^2\bigl(X_0,\Homsh_{\Osh_{X_0}}(F_0,F_0)\bigr)
  \otimes_{A/J}(J^n/J^{n+1}),
\]
que é, portanto, nula por hipótese. Utilizando agora o teorema 3, obtém-se:

\result{COROLÁRIO}{1} Seja $X$ um esquema próprio e plano sobre $A$, e seja $F_0$
como acima. Então existe um feixe localmente livre $F$ sobre $X$ que induz
sobre $X_0$ um feixe isomorfo a $F_0$. Além disso, este $F$ é único salvo
isomorfismo se
$H^1(X_0,\Homsh_{\Osh_{X_0}}(F_0,F_0))=0$.

Seja $X_0$ um esquema de tipo finito sobre o corpo $k$; supõe-se que $X_0$ é
liso (pelo que entendemos absolutamente liso) sobre $k$, mas não
necessariamente próprio sobre $k$. Seja $A$ um anel artiniano local de corpo
residual $k$. Interessa-nos encontrar os esquemas $X$ planos sobre $A$ tais que
$X\otimes_A k=X_0$ (este é o ponto de partida da «teoria de módulos» ou
das «variações de estrutura» de $X_0$). É equivalente dar um tal $X$, ou dar
sobre o espaço topológico $X_0$ um feixe $\Osh_X$ provido das estruturas
seguintes:
\begin{enumerate}
  \item[\textit{i.}] É um feixe de $A$-álgebras.
  \item[\textit{ii.}] Está provido de um homomorfismo de aumentação
    $\Osh_X\longrightarrow\Osh_{X_0}$ (compatível com as estruturas de
    $A$-álgebras); esses dados estão sujeitos às condições seguintes: a
    aumentação induz um isomorfismo
    $\Osh_X\otimes_A k\xrightarrow{\sim}\Osh_{X_0}$; $\Osh_X$ é plano sobre
    $A$, isto é, o graduado associado a $\Osh_X$, filtrado pelas potências
    do ideal maximal $\mathfrak m$ de $A$, é isomorfo a
    $\operatorname{gr}^0(\Osh_X)\otimes_k\operatorname{gr}(A)$; isto é, se
    têm isomorfismos
    \[
      \mathfrak m^n\Osh_X/\mathfrak m^{n+1}\Osh_X
      =\Osh_{X_0}\otimes_k(\mathfrak m^n/\mathfrak m^{n+1}).
    \]
\end{enumerate}
O fato fundamental é o seguinte:

\src{204}
\result{TEOREMA}{8} Seja $X_0$ um esquema de tipo finito e liso sobre o
corpo $k$, e suponhamos que $X_0$ é afim. Seja $A$ um anel local artiniano de
corpo residual $k$. Então existe um $A$-esquema $X$ plano sobre $A$ tal que
$X\otimes_A k=X_0$, e dois esquemas desse tipo são necessariamente isomorfos.

Observe-se que o isomorfismo em questão não é canônico, pois em geral $X$
terá $A$-automorfismos não triviais que induzem a identidade sobre $X_0$. Por
outro lado, em geral não há uma escolha «canônica» de um $X$ que satisfaça
as condições dadas, salvo quando $A$ é uma $k$-álgebra (caso de
características iguais), em cujo caso pode tomar-se $X=X_0\otimes_k A$, isto
é, $\Osh_X=\Osh_{X_0}\otimes_k A$ (seja ou não afim $X_0$). No caso de
características distintas, ignoro em geral, quando $X_0$ não é afim, se se
pode «levantar» $X_0$ a um $X$ definido sobre $A$. No entanto, seja $n$ um
inteiro $>0$, ponhamos $A_{n-1}=A/\mathfrak m^n$, e suponhamos que se tenha
levantado $X_0$ a um $A_{n-1}$-esquema plano $X_{n-1}$; pretende-se levantar
$X_{n-1}$ a um $A_n$-esquema plano $X_n$. Pelo teorema 8, já se sabe que isso
é possível localmente; por outro lado, verifica-se facilmente que, se $U_n$
levanta um aberto $U_{n-1}$ de $X_{n-1}$, então o feixe dos grupos de
automorfismos de $U_n$ (que induzem a identidade sobre $U_{n-1}$) é
canonicamente isomorfo a
\[
  \Theta_{X_0/k}\otimes_k(\mathfrak m^n/\mathfrak m^{n+1})
\]
restrito a $U_{n-1}$ e, em particular, é comutativo
($\Theta_{X_0/k}$ denota o feixe de germes de $k$-derivações sobre $X_0$).
Segue-se facilmente que há uma obstrução para construir $X_n$ que levante
$X_{n-1}$, a qual se encontra em
\[
  H^2(X_0,\Theta_{X_0/k})\otimes_k
  \mathfrak m^n/\mathfrak m^{n+1}.
\]
Por conseguinte:

\result{COROLÁRIO}{1} Seja $X_0$ um esquema de tipo finito liso sobre $k$, e
suponhamos $H^2(X_0,\Theta_{X_0/k})=0$. Então, para todo anel local
artiniano $A$ de corpo residual $k$, existe um $A$-esquema plano $X$ tal que
$X\otimes_A k=X_0$.

Além disso, se foi possível encontrar um $X$ plano sobre $A$ que levante $X_0$,
então, em virtude do teorema 8, o conjunto das classes (salvo isomorfismo)
de $A$-esquemas planos que levantam $X_0$ identifica-se com
\[
  H^1\bigl(X_0,\underline{\operatorname{Aut}}(X)\bigr),
\]
onde, naturalmente, $\underline{\operatorname{Aut}}(X)$ denota o feixe de
germes de automorfismos do feixe de $A$-álgebras $\Osh_X$ compatíveis com a
aumentação. A filtração de $\Osh_X$ define uma filtração de
$\underline{\operatorname{Aut}}(X)$, e o quociente deste feixe pelo $n$-ésimo
subgrupo da filtração é $\underline{\operatorname{Aut}}(X_n)$; o graduado
associado a essa filtração é comutativo e identifica-se com
$\Theta_{X_0/k}\otimes_k\operatorname{gr}(A)$. Em particular, se
$\mathfrak m^{n+1}$ é a primeira potência de $\mathfrak m$ que é nula,
então $F^n(\underline{\operatorname{Aut}}(X))$ (o último termo da
filtração) pertence ao centro de $\underline{\operatorname{Aut}}(X)$ e é
isomorfo a $\Theta_{X_0/k}\otimes_k\mathfrak m^n$; é também o feixe de
germes de automorfismos de $X$ que induzem a identidade sobre
$X_{n-1}=X\otimes_A A/\mathfrak m^n$. Utilizando estes resultados obtêm-se
imediatamente os enunciados seguintes:

\src{205}
\result{COROLÁRIO}{2} Seja $X_0$ um esquema de tipo finito liso sobre $k$, seja
$A$ um anel local artiniano de corpo residual $k$ e ideal maximal
$\mathfrak m$, suponhamos $\mathfrak m^{n+1}=0$, seja
$A_{n-1}=A/\mathfrak m^n$ e, por fim, seja $X_{n-1}$ um
$A_{n-1}$-esquema plano tal que $X_{n-1}\otimes_A k=X_0$. Então o conjunto
das classes (salvo um isomorfismo que induza a identidade sobre $X_{n-1}$)
de $A$-esquemas planos $X_n$ tais que
$X_n\otimes_A A_{n-1}=X_{n-1}$ é vazio ou é um espaço homogêneo principal
sob
\[
  H^1(X_0,\Theta_{X_0/k})\otimes_k\mathfrak m^n.
\]

(Observe-se que, em geral, não há um elemento origem privilegiado neste
último espaço, pois não há uma maneira privilegiada de levantar $X_{n-1}$ a
$X_n$).

\result{COROLÁRIO}{3} Seja $X_0$ um esquema de tipo finito liso sobre $k$, e
suponhamos $H^1(X_0,\Theta_{X_0/k})=0$. Então, para todo anel local
artiniano $A$ de corpo residual $k$, existe no máximo (salvo isomorfismo) um
$A$-esquema plano $X$ tal que $X\otimes_A k=X_0$.

Os corolários 1 e 3 implicam imediatamente os enunciados, de aparência mais
geral, que se obtêm supondo apenas que $A$ é um anel local
noetheriano completo de corpo residual $k$, contanto que se introduza aí $X$
como um esquema formal sobre $A$:

\result{TEOREMA}{9} Sejam $k$ um corpo e $X_0$ um esquema de tipo finito liso
sobre $k$. Para todo anel local noetheriano completo $A$ de corpo residual
$k$, seja $F(A)$ o conjunto das classes (salvo um isomorfismo que induza a
identidade sobre $X_0$) de esquemas formais $\mathfrak X$ sobre $A$, de tipo
finito e planos sobre $A$, tais que $\mathfrak X\otimes_A k=X_0$. Com estas
notações:
\begin{enumerate}
  \item[\textit{i.}] Se $H^1(X_0,\Theta_{X_0/k})=0$, então, para todo $A$,
    $F(A)$ tem no máximo um elemento.
  \item[\textit{ii.}] Se $H^2(X_0,\Theta_{X_0/k})=0$, então, para todo $A$,
    $F(A)$ tem ao menos um elemento.
\end{enumerate}

\src{206}
\result{COROLÁRIO}{1} Suponhamos que $X_0$ é próprio sobre $k$. Sob a
condição \textit{(i)}, para todo $A$ existe no máximo (salvo um isomorfismo que
induza a identidade sobre $X_0$) um esquema $X$ próprio e plano sobre $A$ tal
que $X\otimes_A k=X_0$.

Utiliza-se o teorema 3, corolário 2. Por exemplo:

\result{COROLÁRIO}{2} Se $X$ é um esquema próprio e plano sobre $A$ tal que
$X\otimes_A k$ é isomorfo ao esquema projetivo tipo $\mathbf P_k^r$ de
dimensão $r$ sobre $k$, então $X$ é isomorfo a $\mathbf P_A^r$.

(Também se pode deduzir esse resultado da proposição 3, corolário 1).

\result{COROLÁRIO}{3} Seja $X_0$ um esquema liso e projetivo sobre $k$, e
suponhamos que se tenha
\[
  H^2(X_0,\Osh_{X_0})=H^2(X_0,\Theta_{X_0/k})=0.
\]
Então, para todo $A$, existe um esquema projetivo e plano $X$ sobre $A$ tal
que $X\otimes_A k=X_0$.

Combinam-se o teorema 9\textit{(ii)} e o teorema 4. Em particular:

\result{COROLÁRIO}{4} Seja $X_0$ o esquema de uma curva algébrica completa e
lisa sobre $k$. Então, para todo anel local noetheriano completo $A$ de
corpo residual $k$, existe um «esquema de curva lisa» $X$ sobre $A$ tal que
$X\otimes_A k=X_0$.

\result{OBSERVAÇÕES}{}

1) Os corolários 3 e 4 são especialmente interessantes se $k$ é de
característica $p\ne0$, tomando como $A$ um anel de valoração discreta de
característica $0$ e corpo residual $k$; por exemplo, o «menor $A$ possível»,
a saber, aquele cujo ideal maximal é gerado por $p$. (De fato, pelos
teoremas de Cohen, basta conhecer os corolários 3 e 4 para um anel $A$ desse
tipo). Assinalemos a este respeito que, segundo os especialistas, não se sabe se
existem esquemas sobre um corpo $k$ que não sejam a redução módulo $p$ de um
esquema plano definido sobre um tal anel $A$. Ao menos, os resultados deste
número proporcionam um meio de investigar sistematicamente essa questão. Dever-se-ia
começar por examinar se a primeira obstrução que aparece em
$H^2(X_0,\Theta_{X_0/k})$ é necessariamente nula.

\emph{Complemento de 1962.} --- J.-P. Serre construiu uma variedade
projetiva não singular, de dimensão 3, sobre um corpo algebricamente fechado
$k$ de característica $p>0$, que não se obtém por redução de um esquema
próprio sobre um anel local íntegro de corpo residual $k$ e corpo de
frações de característica zero. Mumford teria encontrado um exemplo análogo
com uma superfície projetiva não singular.

2) Observe-se que o teorema 3, e a técnica correspondente, só são válidos para
um anel de base completo (local, para fixar as ideias).

Para passar de resultados conhecidos para o completamento de um anel local aos
resultados correspondentes para esse mesmo anel local, seria necessário um
quarto «teorema fundamental», cujo enunciado definitivo ainda está por encontrar.

\src{207}
3) Comparem-se os resultados deste número (em particular, os corolários 1 e
2 precedentes) e do seguinte com os resultados de Kodaira--Spencer sobre a
variação das estruturas complexas. Utilizando o teorema conjectural ao qual
se acaba de aludir, sob as condições do corolário 1, mas sem supor já
que $A$ seja completo, se deveria poder concluir que existe um anel $A'$ que
contém $A$, finito e não ramificado sobre $A$, tal que $X\otimes_A A'$ e
$X'\otimes_A A'$ são $A'$-isomorfos (sendo $X,X'$ dois $A$-esquemas próprios e
planos dados, tais que $X\otimes_A k=X'\otimes_A k=X_0$). Isto é ao menos o
que pode demonstrar-se quando $X_0=\mathbf P_k^r$, utilizando o corolário da
proposição 2. Em todo caso, quando $H^1(X_0,\Theta_{X_0/k})=0$, pode
demonstrar-se que as fibras de $X,X'$ sobre todo ponto $y$ de
$Y=\operatorname{Spec}(A)$ são isomorfas, ao menos depois de passar ao
fecho algébrico do corpo residual $\kappa(y)$. (Tem-se um resultado
local, aparentemente mais forte, sem supor já necessariamente que $A$ seja
local). Quanto às «variações de estrutura» do espaço projetivo,
assinalemos além disso a questão seguinte, sugerida por um problema correspondente
de Kodaira--Spencer. Seja $X$ um esquema próprio e plano sobre o anel local
íntegro $A$ de corpo de frações $K$ e corpo residual $k$, e suponhamos que
$X\otimes_A K$ é isomorfo a $\mathbf P_K^r$; é verdade que
$X\otimes_A k=X_0$ é isomorfo a $\mathbf P_k^r$ (ao menos sobre o fecho
algébrico de $k$)? Nesse problema pode-se supor que $A$ é um anel de
valoração discreta completo. Coloca-se um problema análogo quando $X_0$ é
uma variedade abeliana.

\emph{Complemento de 1962.} --- O conferencista é agora menos otimista
com respeito aos resultados conjecturados acima. Não obstante, Hironaka resolveu
afirmativamente o problema relativo às variações de estrutura
do espaço projetivo, e Koizumi o problema análogo para as variedades
abelianas.

\fsection{7}{Aplicação à «Teoria dos módulos»}

Como o próprio conferencista apenas começou a abordar essa teoria, seremos
obrigados a nos limitar a algumas indicações. Para simplificar, situamo-nos
sobre um corpo $k$, isto é, trabalhamos em igual característica, embora o
teorema 8 permita tratar também o caso geral sem mudança essencial, ao que
parece. Atualmente não superamos a etapa «formal»; no entanto, o
conferencista espera poder construir a partir dela verdadeiros esquemas de
módulos em certos casos e, em particular, construir para cada inteiro $g$ um
esquema sobre os inteiros que desempenhe o papel de esquema de módulos universal
para as curvas lisas de gênero $g$.

\emph{Complemento de 1962.} --- Mumford construiu os esquemas de módulos
para as curvas de gênero $g$ (cf. SHMI). O teorema 10 demonstra além disso que
os «esquemas de Jacobi de nível $n$» da teoria de módulos são não singulares,
e inclusive lisos sobre $\mathbf Z$.

\src{208}
Retomemos a situação e as notações do teorema 9, supondo agora que $A$
é uma álgebra local de posto finito sobre $k$, que para simplificar se supõe
algebricamente fechado. Então $F(A)$ pode ser considerado como um funtor
covariante de $A$, com valores na categoria de conjuntos: um homomorfismo de
$k$-álgebras $A\longrightarrow B$ define uma aplicação
$F(A)\longrightarrow F(B)$, pois todo $A$-esquema plano $X$ tal que
$X\otimes_A k=X_0$ dá lugar a um $B$-esquema $X\otimes_A B$ com as mesmas
propriedades. Suponhamos que se possa encontrar uma $k$-álgebra local noetheriana
completa $\mathcal O$ e um isomorfismo funtorial
\begin{equation}
  \tag{*}
  \operatorname{Hom}(\mathcal O,A)\xrightarrow{\sim}F(A)
\end{equation}
(onde o primeiro membro denota os homomorfismos de $k$-álgebras). Vê-se
facilmente que um tal $\mathcal O$ está determinado salvo isomorfismo canônico;
chama-se então ao espectro formal $\mathfrak Y$ de $\mathcal O$ (isto é,
ao espaço topológico reduzido a um ponto e provido de um feixe de anéis
topológicos reduzido a $\mathcal O$) o \emph{esquema formal de módulos} para
$X_0$. (Observe-se que ele não existe necessariamente). Seja $\mathfrak r$ o ideal
maximal de $\mathcal O$; para cada $n$, seja
$\mathcal O_n=\mathcal O/\mathfrak r^{n+1}$ (portanto $\mathcal O_0=k$).
Então o homomorfismo canônico $\mathcal O\longrightarrow\mathcal O_n$ é um
elemento de $\operatorname{Hom}(\mathcal O,\mathcal O_n)$, pelo que define um
elemento de $F(\mathcal O_n)$, isto é, um $\mathcal O_n$-esquema plano $X_n$
cuja redução módulo $\mathfrak r$ é $X_0$. Esses $X_n$ se deduzem uns dos
outros por extensão de escalares (isto é, aqui por reduções), de onde
resulta que provêm de um esquema formal $\mathfrak X$ bem determinado sobre
o esquema formal de módulos $\mathfrak Y$; $\mathfrak X$ é plano sobre
$\mathfrak Y$ e $\mathfrak X_0=X_0$. Como se vê imediatamente, o isomorfismo (*)
obtém-se então associando a todo homomorfismo de $k$-álgebras
$\mathcal O\longrightarrow A$ a classe do $A$-esquema
$\mathfrak X\otimes_{\mathcal O}A$ (isto é, a todo morfismo de $k$-esquemas
$Y'=\operatorname{Spec}(A)\longrightarrow\mathfrak Y$ associa-se o
$Y'$-esquema $\mathfrak X\otimes_{\mathfrak Y}Y'$ obtido por mudança de base).
Além disso, vê-se que o isomorfismo (*) e sua descrição precedente continuam
válidos se somente se supõe que $A$ é uma $k$-álgebra local noetheriana e
completa (não necessariamente artiniana). Naturalmente, como de costume,
$\mathcal O$ pode perfeitamente ter \textit{a priori} elementos nilpotentes,
e parece provável que existam casos em que $\mathcal O$ seja ele mesmo artiniano
sem ser idêntico a $k$. Isto mostra até que ponto o ponto de vista de
Kodaira--Spencer (que se limita a tomar como $A$ anéis regulares) é
inadequado \textit{a priori} no caso geral.

\src{209}
Fica por dar condições suficientes para que exista um esquema formal de
módulos para $X_0$, que se supõe próprio sobre $k$. Em geral, é fácil dar
condições necessárias e suficientes simples sobre um funtor
$A\longmapsto F(A)$ (de $k$-álgebras locais de posto finito, com valores nos
conjuntos) para que possa ser escrito na forma
$\operatorname{Hom}(\mathcal O,A)$ para um $\mathcal O$ adequado. Não as
detalharemos aqui. Assinalemos apenas que, no caso que nos ocupa, essas
condições impõem sobre $X_0$ condições cohomológicas não triviais, e parece
pouco provável que sejam sempre satisfeitas, embora o conferencista não tenha
construído nenhum contraexemplo. Em contrapartida, parece plausível que a condição
$H^0(X_0,\Theta_{X_0/k})=0$ seja uma condição suficiente (embora de modo algum
necessária) para a existência de um esquema formal de módulos. Limitamo-nos a
enunciar aqui um teorema num caso particularmente simples (cujo análogo na
teoria de espaços analíticos é bem conhecido, cf. Kodaira--Spencer), que se
estabelece sem dificuldade mediante os resultados do número anterior:

\result{TEOREMA}{10} Seja $X_0$ um esquema próprio e liso sobre o corpo $k$
tal que
\[
  H^0(X_0,\Theta_{X_0/k})=H^2(X_0,\Theta_{X_0/k})=0.
\]
Então existe um esquema formal de módulos para $X_0$, correspondente a um
anel local regular $\mathcal O$ (isto é, uma álgebra de séries formais sobre
$k$).

Como já assinalamos, não é certo em geral que o esquema formal
$\mathfrak X$ sobre $\mathcal O$ seja algebrizável; mas sabe-se que o é
quando $X_0$ é projetivo e $H^2(X_0,\Osh_{X_0})=0$ (teorema 4), por exemplo
quando $X_0$ tem dimensão 1. Isto dá certo fundamento à esperança de
construir um esquema de módulos sobre os inteiros para as curvas de gênero
dado\ldots{}

Observe-se também que métodos como os expostos neste número podem ser aplicados
à construção e ao estudo das variedades de Picard, assim como a muitas
outras construções. Voltaremos a isso em breve.

\fsection{8}{Aplicação ao grupo fundamental}

As técnicas expostas permitem abordar o estudo sistemático do grupo
fundamental seguindo o modelo da teoria topológica. Os dois primeiros
teoremas enunciados neste número são generalizações de resultados de um
trabalho recente de Lang--Serre.

Seja $X$ um esquema. Um $X$-esquema $X'$ se chama recobrimento não ramificado de
$X$ se
\begin{enumerate}
  \item[\textit{i.}] $X'$ é finito sobre $X$, isto é, está definido por um
    feixe coerente de álgebras $\mathcal A=\mathcal A(X')$ sobre $X$;
  \item[\textit{ii.}] $\mathcal A$ é um feixe localmente livre sobre $X$;
  \item[\textit{iii.}] Para todo $x\in X$, o quociente
    \[
      \mathcal A_x/\mathfrak m_x\mathcal A_x
      =\mathcal A_x\otimes_{\Osh_x}\kappa(x)
    \]
    é um álgebra separável sobre $\kappa(x)$.
\end{enumerate}

\src{210}
Essa noção de recobrimento não ramificado (devida a Serre e ao conferencista)
possui todas as propriedades elementares que se pode razoavelmente esperar, cuja
lista não daremos. Limitemo-nos a dizer que dá lugar a uma teoria de Galois
modelada sobre a teoria de Galois clássica (e que a contém; as
demonstrações são bem mais simples que as comumente aceitas para
esta última) e sobre a teoria de Galois dos recobrimentos topológicos. Mais
precisamente, chamemos \emph{ponto geométrico} de um esquema $X$ um morfismo
$a$ do espectro $\xi$ de um corpo algebricamente fechado $\Omega$ em $X$,
isto é, o dado de uma extensão algebricamente fechada do corpo residual
$\kappa(x)$ de um ponto $x=|a|$ de $X$ (chamado \emph{localização} do ponto
geométrico $a$). Se $X'$ é então um recobrimento não ramificado de $X$, se
lhe pode associar o conjunto $E_a(X')$ dos «pontos geométricos de $X'$ sobre
$a$», isto é, o conjunto dos pares formados por um $x'\in X'$ sobre $x$
e um $\kappa(x)$-homomorfismo $\kappa(x')\longrightarrow\Omega$. Obtém-se assim
(para $(X,a)$ fixos) um funtor $F(X,a)$ da categoria $R(X)$ dos
recobrimentos não ramificados $X'$ de $X$ na categoria dos conjuntos
finitos. Se $X$ é conexo, o par formado por $R(X)$ e $F(X,a)$ possui as
propriedades formais necessárias para ser isomorfo ao par análogo definido por
um grupo topológico compacto totalmente descontínuo $\pi$ adequado (isto é,
um limite projetivo de grupos finitos): toma-se a categoria $C(\pi)$ dos
conjuntos finitos $E$ sobre os quais $\pi$ atua continuamente, e o funtor
identidade $F(\pi)(E)=E$ desta categoria na categoria dos conjuntos
finitos. Além disso, o grupo $\pi$ está determinado salvo isomorfismo canônico pela
condição de que $(C(\pi),F(\pi))$ seja isomorfo a um par dado. Nesse caso,
$\pi$ chama-se \emph{grupo fundamental} do esquema conexo $X$ no ponto
geométrico $a$, e denota-se por $\pi_1(X,a)$. Se $X$ não é conexo, substitui-se
o esquema pela componente conexa de $x=|a|$. Se, pelo contrário, $X$ é
conexo, os grupos $\pi_1(X,a)$ correspondentes a dois pontos geométricos
$a',a''$ de $X$ são isomorfos (o isomorfismo está determinado salvo
automorfismos interiores), de modo que, como de costume, pode-se escolher $a$
da maneira mais conveniente, por exemplo no ponto genérico de $X$, suposto
irredutível. Naturalmente, $\pi_1(X,a)$ é um funtor covariante do esquema
pontuado $(X,a)$. Todo enunciado relativo à classificação dos
recobrimentos não ramificados\footnote{O texto impresso diz «recobrimentos
inseparáveis»; o parágrafo trata da categoria $R(X)$ dos recobrimentos não
ramificados e de sua equivalência com os conjuntos finitos dotados de uma ação
contínua de $\pi_1(X,a)$.} pode ser traduzido então para a linguagem da teoria de
grupos seguindo o dicionário bem conhecido (salvo que se deve levar em conta
que aqui se trata de grupos topológicos).

\src{211}
Nosso objetivo é desenvolver o análogo da sequência exata de homotopia
dos espaços fibrados para um morfismo próprio $f:X\longrightarrow Y$.
Evidentemente, como não sabemos quais são os grupos de homotopia superiores, só
obteremos resultados necessariamente incompletos. Para poder aplicar os
teoremas fundamentais do n\textsuperscript{o} 3, devemos explicitar primeiro
alguns lemas elementares relativos aos esquemas sobre corpos ou anéis de
Artin (conforme o procedimento geral!).

\result{LEMA}{1} Seja $(X',a')$ um recobrimento não ramificado pontuado associado
a uma representação pontuada de $\pi_1(X,a)$ em um conjunto finito $E$
(provido do ponto marcado $e$). Então o morfismo canônico
$\pi_1(X',a')\longrightarrow\pi_1(X,a)$ identifica o primeiro grupo com o
estabilizador de $e$ em $\pi_1(X,a)$ (e é, portanto, injetivo).

\result{LEMA}{2} Seja $X$ um esquema algébrico sobre o corpo $k$, e seja $k'$
uma extensão radicial de $k$. Então todo recobrimento não ramificado de
$X\otimes_k k'$ provém por imagem inversa (isto é, por extensão de
escalares) de um recobrimento não ramificado de $X$, determinado salvo
isomorfismo canônico.

Desses dois lemas resulta, em particular, que para toda extensão algébrica
$K$ de $k$, e para todo ponto geométrico $a'$ de $X'=X\otimes_k K$ que se
projeta sobre o ponto geométrico $a$ de $X$, o homomorfismo funtorial
$\pi_1(X',a')\longrightarrow\pi_1(X,a)$ é injetivo.

\result{LEMA}{3} Seja $X$ um esquema completo sobre um anel local artiniano
$A$, tal que $H^0(X,\Osh_X)=A$; seja $X'$ um recobrimento não ramificado de $X$,
e seja $A'=H^0(X',\Osh_{X'})$, que é, portanto, um anel finito sobre $A$
(que \textit{a priori} pode estar ramificado sobre $A$). Sejam $X_0,X'_0$ os
subesquemas reduzidos associados a $X,X'$ (obtidos dividindo pelos feixes de
elementos nilpotentes de $\Osh_X,\Osh_{X'}$), e seja $k$ um subcorpo de
$A/\mathfrak r(A)$ sobre o qual $A/\mathfrak r(A)$ seja finito (assim, $X_0$ é
um esquema algébrico completo sobre $k$, e $X'_0$ é um recobrimento não
ramificado seu). Seja, por fim, $\Omega$ uma extensão algebricamente
fechada de $k$, e consideremos o recobrimento não ramificado
$X'_0\otimes_k\Omega$ de $X_0\otimes_k\Omega$. Então as duas condições
seguintes são equivalentes:
\begin{enumerate}
  \item[\textit{i.}] $X'_0\otimes_k\Omega$ está completamente decomposto
    sobre $X_0\otimes_k\Omega$.
  \item[\textit{ii.}] O morfismo natural
    $X'\longrightarrow X\otimes_A A'$ é um isomorfismo.
\end{enumerate}
Sob essas condições, $A'$ é uma extensão não ramificada de $A$. Por fim,
se $X'$ é conexo, a condição \textit{(i)} equivale à condição seguinte,
de aparência mais fraca:
\begin{enumerate}
  \item[\textit{i bis.}] $X'_0\otimes_k\Omega$ admite uma seção regular
    sobre $X_0\otimes_k\Omega$.
\end{enumerate}
Quando se satisfaz a condição \textit{(ii)}, diremos que o recobrimento não
ramificado $X'$ de $X$ é \emph{geometricamente trivial}.

\src{212}
\result{LEMA}{4} Seja $f:X\longrightarrow Y$ um morfismo próprio tal que
$f_*(\Osh_X)=\Osh_Y$. Sejam $a$ um ponto geométrico de $X$ e $b$ sua projeção
sobre $Y$. Então $\pi_1(X,a)\longrightarrow\pi_1(Y,b)$ é sobrejetivo.

Com efeito, é preciso demonstrar o seguinte: se um recobrimento não ramificado
$Y'$ de $Y$ (correspondente a um feixe de álgebras localmente livre
$\mathcal A$) é tal que $X\otimes_Y Y'$ é desconexo, então $Y'$ também o
é. Com efeito, $f^*(\mathcal A)$ será soma direta de dois feixes de anéis não
nulos e, portanto, também o será sua imagem direta; ora, esta última
não é outra que $\mathcal A\otimes f_*(\Osh_X)=\mathcal A$.

\result{LEMA}{5} Seja $X$ um esquema completo sobre um corpo $k$; supõe-se que
$H^0(X,\Osh_X)$ é um anel local $A$ e que $A/\mathfrak r(A)$ é radicial
sobre $k$. Seja $\Omega$ um fecho algébrico de $k$, e seja
$\overline X=X\otimes_k\Omega$ (é conexo). Escolhamos um ponto geométrico
$\overline a$ de $\overline X$ que se projete sobre o ponto geométrico $a$ de
$X$; então tem-se uma sequência exata
\[
  e\longrightarrow\pi_1(\overline X,\overline a)
  \longrightarrow\pi_1(X,a)
  \longrightarrow\pi_1(k,b)\longrightarrow e
\]
(onde $\pi_1(k,b)$ é o grupo de Galois de $\Omega$ sobre $k$).

A injetividade do primeiro homomorfismo foi vista com os lemas 1 e 2, a
exatidão no meio resulta do lema 3 e, por fim, a sobrejetividade do
último homomorfismo (que é a única que utiliza o fato de que
$A/\mathfrak r(A)$ seja radicial) resulta do lema 4.

\result{PROPOSIÇÃO}{4} Seja $f:X\longrightarrow Y$ um morfismo próprio e plano
tal que, para todo $y\in Y$, $H^0(f^{-1}(y),\Osh_{f^{-1}(y)})$ seja uma álgebra
separável sobre o corpo residual $\kappa(y)$ (o que ocorre, por exemplo, se
$f^{-1}(y)$ é um esquema separável sobre $\kappa(y)$, isto é, reduzido e tal
que os corpos correspondentes a suas componentes irredutíveis sejam
extensões separáveis de $\kappa(y)$). Então o recobrimento $Y'$ de $Y$
associado a $f_*(\Osh_X)$ é não ramificado.

A demonstração é fácil graças ao teorema 2.

\src{213}
Essa proposição, junto com o lema 1, reduz praticamente o estudo
homotópico dos morfismos próprios e planos (com fibras separáveis) ao caso em
que $f_*(\Osh_X)=\Osh_Y$ (pois, utilizando a fatoração de Stein, se
substitui $Y$ por $Y'$).

\result{OBSERVAÇÃO}{} Um morfismo plano de tipo finito cujas fibras são
esquemas separáveis (resp. lisos) chama-se \emph{separável} (resp.
\emph{liso}). Demonstra-se que, se $f$ é plano e $f^{-1}(y)$ é separável
(resp. liso), existe uma vizinhança de $f^{-1}(y)$ na qual $f$ é separável
(resp. liso). O mesmo resultado é válido para «absolutamente normal»
(«teorema de Bertini»).

Seja $f:X\longrightarrow Y$ um morfismo próprio tal que
\begin{equation}
  \tag{i}
  f_*(\Osh_X)=\Osh_Y,
\end{equation}
e seja $X'$ um esquema finito sobre $X$. Seja $Y'$ o recobrimento de $Y$
correspondente à fatoração de Stein de $X'\longrightarrow Y$
(cf. teorema 5). Seja $y\in Y$; então o conjunto das componentes conexas
da fibra $F'$ de $X'$ sobre $y$ identifica-se com o conjunto dos pontos
$y'$ de $Y'$ sobre $y$ (teorema 5). Consideremos o morfismo evidente
\begin{equation}
  \tag{*}
  X'\longrightarrow X\times_Y Y'
\end{equation}
deduzido dos morfismos naturais $X'\longrightarrow X$ e
$X'\longrightarrow Y'$\footnote{O texto impresso omite a prima no segundo
$X'$; o morfismo para $Y'$ é o da fatoração de Stein de
$X'\to Y$.}; será um isomorfismo sempre que $X'$ seja da forma
$X\times_Y Y''$, onde $Y''$ é um recobrimento não ramificado de $Y$, e
então $Y'$ não será outro senão $Y''$ e (*) será a identidade. Queremos dar
precisamente condições sob as quais $X'$ seja da forma que acabamos de
indicar, isto é, que $Y'$ seja não ramificado e (*) um isomorfismo. Para isso,
introduzamos a fibra $F$ de $X$ em $y$; é um esquema próprio sobre
$\kappa(y)$ do qual $F'$ é um recobrimento (não ramificado se $X'$ o for).
Seja $F'_1$ uma componente conexa de $F'$ correspondente a um ponto $y'_1$ de
$Y'$ sobre $y$. Suponhamos além disso
\begin{enumerate}
  \item[\textit{(ii)}] $X'$ é não ramificado sobre $X$ nos pontos de $F'_1$
    (portanto, $F'_1$ é um recobrimento não ramificado de $F$), e
  \item[\textit{(iii)}] $F'_1$ é um recobrimento geometricamente trivial de
    $F$ (cf. lema 3).
\end{enumerate}

\result{TEOREMA}{11} Nessas condições, existe uma vizinhança aberta $U'$ de
$y'_1$ em $Y'$ tal que (*) é um isomorfismo sobre $U'$. Além disso, $Y'$ é não
ramificado em $y'_1$ sobre $Y$ (mas pode ser ramificado nos demais pontos
$y'$ de $Y'$ sobre $y$).

Naturalmente, as condições \textit{(ii)} e \textit{(iii)} são também
necessárias para que seja válida a conclusão do teorema. A demonstração do
teorema é fácil, com auxílio do lema 3 e do teorema 2.

\src{214}
\result{COROLÁRIO}{1} Suponhamos sempre satisfeita \textit{(i)}. Para que um
recobrimento não ramificado sobre $X$ seja isomorfo à imagem inversa de um
recobrimento não ramificado $Y'$ de $Y$, é necessário e suficiente que $X'$
induza sobre cada fibra $f^{-1}(y)$ um recobrimento geometricamente trivial.

Como, segundo o teorema 11, o conjunto dos pontos de $Y$ para os quais se
cumpre essa condição é aberto, basta verificá-la nos pontos $y$ que são
fechados\ldots{} Assinalemos o enunciado seguinte, equivalente ao corolário 1:
o núcleo do homomorfismo (sobrejetivo segundo o lema 4)
$\pi_1(X)\longrightarrow\pi_1(Y)$ é o subgrupo invariante fechado gerado
pelas imagens em $\pi_1(X)$ dos $\pi_1(\overline{f^{-1}(y)})$, onde
\[
  \overline{f^{-1}(y)}
  =f^{-1}(y)\otimes_{\kappa(y)}\overline{\kappa(y)},
\]
$\overline{\kappa(y)}$ designa um fecho algébrico de $\kappa(y)$. Observar-se-á
que, como não se pode escolher o mesmo ponto base para todas as fibras,
os homomorfismos $\pi_1(\overline{f^{-1}(y)})\longrightarrow\pi_1(X)$ só
ficam determinados em qualquer caso (uma vez escolhido um ponto base para $X$
e, portanto, para $Y$) módulo composição por um automorfismo interior em
$\pi_1(X)$.

\result{COROLÁRIO}{2} Nas condições gerais do teorema 11, suponhamos
além disso que $Y,X,X'$ são íntegros, e sejam $K,L,L'$ seus corpos. Então existe
uma subextensão separável $K'$ de $K$ em $L'$, linearmente disjunta de $L$, tal
que $L'=LK'$ (de onde $L'=L\otimes_K K'$).

(Aplica-se a última parte do lema 3 à fibra geométrica de $X$). O caso de
aplicação mais interessante do teorema 11 obtém-se quando $f$ é um morfismo
separável. Então $X'$ é também separável sobre $Y$; portanto, em virtude
da proposição 4, $Y'$ é não ramificado sobre $Y$, e o segundo membro
$X\times_Y Y'$ de (*) é não ramificado sobre $X$. Conclui-se facilmente:

\result{COROLÁRIO}{3} Suponhamos, além de \textit{(i)}, que $f$ é separável.
Seja $X'$ um recobrimento não ramificado e conexo de $X$; para que
$X'$\footnote{O texto impresso diz aqui «$X$»; o objeto do critério é o
recobrimento $X'$ introduzido na frase anterior.} seja imagem inversa de um
recobrimento não ramificado $Y'$ de $Y$, é necessário e suficiente que o
recobrimento $\overline F'$ induzido sobre uma fibra geométrica
$\overline F=f^{-1}(y)$ admita uma seção regular.

Observar-se-á que não era necessário supor que $\overline F'$ fosse
geometricamente trivial sobre $\overline F$ (o que será verdadeiro \textit{a
posteriori}, embora \textit{a priori} esta condição seja muito mais forte). O
corolário 3 é, de fato, equivalente ao enunciado seguinte:

\src{215}
\result{COROLÁRIO}{4} Seja $f:X\longrightarrow Y$ um morfismo próprio e separável
tal que $f_*(\Osh_X)=\Osh_Y$, seja $\overline F$ a fibra geométrica de um ponto
$y\in Y$, e escolhamos um ponto geométrico em $\overline F$, do qual se obtêm,
mediante os morfismos $\overline F\longrightarrow X\longrightarrow Y$, pontos
geométricos em $X,Y$, que serão tomados como pontos base para os grupos
fundamentais de $\overline F,X,Y$. Nessas condições tem-se a sequência
exata
\[
  \boxed{
    \pi_1(\overline F)\longrightarrow\pi_1(X)
    \longrightarrow\pi_1(Y)\longrightarrow e
  }.
\]

Daí se deduzem facilmente os dois enunciados seguintes de Serre--Lang,
liberados aqui de toda hipótese de normalidade:

\result{COROLÁRIO}{5} Sejam $X,Y$ esquemas conexos sobre um corpo $k$;
suponhamos que o esquema reduzido $X_{\mathrm{réd}}$ é separável sobre $k$
(o que é automaticamente verdadeiro se $k$ é perfeito), e que $X$ ou $Y$ é
próprio sobre $k$.\footnote{Condição acrescentada na errata de 1962.}
Escolhamos um ponto geométrico $a$ (resp. $b$) em $X$ (resp. $Y$), de onde se
obtêm um ponto geométrico $c=(a,b)$ em $X\times_kY$ e um morfismo natural
\[
  \pi_1(X\times_kY,c)\longrightarrow
  \pi_1(X,a)\times\pi_1(Y,b)
\]
(deduzido dos morfismos funtoriais de $\pi_1(X\times Y,c)$ em $\pi_1(X,a)$
e $\pi_1(Y,b)$). Este morfismo é injetivo, e inclusive bijetivo se $k$ é
algebricamente fechado.

(A sobrejetividade neste último caso é quase trivial). Daí se deduz,
com Serre--Lang:

\result{COROLÁRIO}{6} Seja $X$ um esquema algébrico conexo e próprio sobre um corpo
algebricamente fechado $k$, e seja $K$ uma extensão algebricamente fechada
de $k$; então os grupos fundamentais de $X$ e $X\otimes_kK$ são os
mesmos, isto é, todo recobrimento não ramificado deste último esquema
provém por extensão de escalares de um recobrimento não ramificado (único
salvo isomorfismo) de $X$.

\result{OBSERVAÇÕES}{}

1) Utilizando a proposição 4, vê-se que a hipótese
$f_*(\Osh_X)=\Osh_Y$ do corolário 4 não é essencial. No caso geral, em vez
de colocar o grupo unidade $e$ depois de $\pi_1(Y)$, é preciso continuar com
\[
  \pi_0(\overline F)\longrightarrow\pi_0(X)
  \longrightarrow\pi_0(Y)\longrightarrow e.
\]
como em topologia algébrica.

\src{216}
2) Em geral, por ora não se pode dizer nada sobre o núcleo de
$\pi_1(\overline F)\longrightarrow\pi_1(X)$, no qual deveria intervir um
$\pi_2(Y)$. Parece, no entanto, que se deveria poder demonstrar que
$\pi_1(\overline F)\longrightarrow\pi_1(X)$ é injetivo se $Y$ é o espectro
de um anel local $A$, apoiando-se no teorema 12 mais adiante (que afirma que
isso ocorre se $A$ é completo).

O teorema 11 utilizava somente os teoremas 1 e 2. Vamos aplicar agora o
teorema 3, apoiando-nos no lema elementar seguinte:

\result{LEMA}{6} Seja $X$ um esquema, e $X_0$ o esquema reduzido correspondente
(isto é, no qual se anularam os elementos nilpotentes). Então todo
recobrimento não ramificado $X'_0$ de $X_0$ é induzido por um recobrimento
não ramificado $X'$ de $X$, determinado salvo isomorfismo canônico.

Este lema, de natureza puramente local, desempenha aqui um papel análogo ao do
teorema 8 em teoria de módulos. Combinando-o com o teorema de existência
(teorema 3), deduz-se:

\result{TEOREMA}{12} Seja $A$ um anel local noetheriano e completo de corpo
residual $k$. Seja $X$ um esquema próprio sobre $A$. Então todo recobrimento
não ramificado $X'_0$ de $X_0=X\otimes_A k$ é induzido por um recobrimento não
ramificado $X'$ de $X$, único salvo isomorfismo.

Em outros termos:

\result{COROLÁRIO}{1} Escolhamos um ponto geométrico em $X_0$ como ponto base para
os grupos fundamentais de $X_0$ e $X$. Então o homomorfismo canônico
$\pi_1(X_0)\longrightarrow\pi_1(X)$ é um isomorfismo.

Apliquemos agora a $X_0$ o lema 5 (supondo, para simplificar, que
$H^0(X_0,\Osh_{X_0})=k$), e observemos que, por ser $A$ completo, suas
extensões não ramificadas correspondem às extensões não ramificadas de seu
corpo residual, isto é, tem-se $\pi_1(Y)=\pi_1(k)$ (onde
$Y=\operatorname{Spec}(A)$). Obtém-se a sequência exata:
\[
  e\longrightarrow\pi_1(\overline X_0)
  \longrightarrow\pi_1(X)\longrightarrow\pi_1(Y)\longrightarrow e.
\]

\result{COROLÁRIO}{2} Seja $f:X\longrightarrow Y$ um morfismo próprio e plano,
sejam $y_1$ um ponto de $Y$ e $y_0$ uma especialização de $y_1$, e consideremos
as fibras «geométricas» correspondentes $\overline X_1$ e $\overline X_0$;
supõe-se que $\overline X_0$ é separável e conexa (o que implica que
$\overline X_1$ satisfaz as mesmas condições). Então pode-se encontrar um
homomorfismo de grupos
$\pi_1(\overline X_1)\longrightarrow\pi_1(\overline X_0)$, definido salvo um
automorfismo interior, e este homomorfismo é sobrejetivo.

\src{217}
Poder-se-ia esperar que este homomorfismo fosse sempre bijetivo; infelizmente,
em geral não é assim se $\kappa(y_0)$ tem característica $>0$. Obteremos,
no entanto, mais adiante uma cota para o núcleo deste homomorfismo (ao menos se
$\overline X_0$ é liso), da qual se segue que, se $\kappa(y_0)$ tem
característica $0$, o homomorfismo anterior é bijetivo (resultado que também
poderia ser demonstrado por via transcendente). Em todo caso, obtemos já uma cota
para o $\pi_1$ de uma fibra especial por meio daquele de uma fibra genérica.
Utilizando, por exemplo, o fato de que uma curva algébrica em característica
$p$ se pode levantar a uma curva em característica $0$ (teorema 9, corolário
4), obtém-se por via transcendente:

\result{COROLÁRIO}{3} Seja $X_0$ o esquema de uma curva completa lisa sobre
um corpo algebricamente fechado de característica arbitrária, e seja $g$ o
gênero de $X_0$; então $\pi_1(X_0)$ admite $2g$ geradores topológicos,
ligados pela relação bem conhecida.

Daí se deduz, mediante uma técnica bem conhecida de seções
hiperplanas:

\result{COROLÁRIO}{4} Seja $X$ um esquema projetivo liso sobre um corpo
algebricamente fechado de característica arbitrária; então $\pi_1(X)$
admite um número finito de geradores topológicos.

Determinemos o núcleo do homomorfismo
$\pi_1(\overline X_1)\longrightarrow\pi_1(\overline X_0)$. Para isso, pode-se
supor que $Y$ é o espectro de um anel de valoração discreta completo
$V=\Osh_y$ ($y=y_0$). A questão equivale à seguinte: dado um recobrimento
não ramificado $\overline X'_1$ de $\overline X_1$ (que, se se quiser, pode
ser suposto galoisiano), determinar em que condições provém de um
recobrimento não ramificado de $X_0$. \textit{A priori}, o recobrimento dado
provém, por extensão de escalares, de um recobrimento não ramificado $X'_1$
de $X_1\otimes_K K'$, onde $K'$ é uma subextensão finita do fecho
algébrico $\overline K$ do corpo de frações $K$ de $V$; se
$\overline X'_1$ fosse galoisiano de grupo $G$, pode-se escolher $X'_1$
galoisiano de grupo $G$. Tem-se então: para que
$X'_1\otimes_{K'}\overline K=\overline X'_1$ provenha de um recobrimento não
ramificado de $X_0$, é necessário e suficiente que se possa encontrar uma extensão
finita $K''$ de $K'$ em $\overline K$ tal que
$X''_1=X'_1\otimes_{K'}K''$ seja da forma $X''\otimes_{V''}K''$, onde $V''$
é o fecho normal de $V$ em $K''$ e $X''$ é um recobrimento não ramificado
de $X\otimes_VV''$. Suponhamos, por exemplo, que $X_0$ é absolutamente normal;
daí resulta que $X\otimes_VV''$ é normal (pois é plano sobre $V''$ e tem
fibra especial normal), e seu corpo de funções é idêntico ao $K''(X_1)$ de
\[
  X_1\otimes_KK''
  =(X\otimes_VK)\otimes_KK''
  =X\otimes_VK''
  =(X\otimes_VV'')\otimes_{V''}K''.
\]
Seja $L''=K''(X'_1)$ o corpo de funções de $X'_1\otimes_{K'}K''$; é uma
extensão finita separável de $K''(X_1)$, e a condição anterior significa
também que $L''$ é uma extensão não ramificada do corpo de funções de
$X\otimes_VV''$ (isto é, a normalização de $X\otimes_VV''$ em $L''$ é não
ramificada sobre $X\otimes_VV''$). Basta, de resto, verificar que $L''$ é
não ramificado nos pontos da fibra especial de $X\otimes_VV''$ (pois é não
ramificado sobre a fibra genérica $X_1\otimes_KK''$). Se agora $X_0$ é liso,
do «teorema de pureza» de Nagata--Zariski resulta inclusive que basta verificar
que $L''$ é não ramificado sobre o anel local $O''$ do ponto genérico da
fibra especial de $X\otimes_VV''$, que é um anel de valoração discreta,
normalização em $K''(X_1)$ do anel local $O\subset K(X_1)$ do ponto
genérico da fibra especial de $X$. Portanto, obteve-se:

\src{218}
\result{COROLÁRIO}{5} Nas condições e com as notações anteriores, para
que o recobrimento não ramificado
$X'_1\otimes_{K'}\overline K$ de
$\overline X_1=X_1\otimes_K\overline K$ provenha de um recobrimento não
ramificado de $\overline X_0$, é necessário e suficiente que exista uma
subextensão finita $K''$ de $\overline K/K'$ tal que $K''(X'_1)$ seja não
ramificado sobre o anel de valoração discreta $O''\subset K''(X_1)$.

Observemos agora que $O''$ é a normalização em $K''(X_1)$ do anel de
valoração discreta $O'\subset K'(X_1)$ (normalização de $O$ em $K'(X_1)$), e
que $O'$ contém a normalização $V'$ de $V$ em $K'$; um uniformizante $u$ de
$V'$ é também um uniformizante de $O'$. Suponhamos agora que $X'_1$ é
galoisiano de grupo de Galois $G$, de ordem $n$ primo com a característica $p$
de $\kappa(y_0)$ (que é também a característica do corpo residual de
$O'$). Portanto, $K'(X'_1)$ está «moderadamente ramificado» sobre $O'$, de onde
resulta facilmente («lema de Abhyankar») que, se se lhe adjunta uma raiz
$n$-ésima $v$ de um uniformizante de $O'$, torna-se não ramificado sobre a
normalização de $O'$ em $K'(X_1)(v)$. Ora, pode-se tomar como $v$ uma
raiz $n$-ésima de um uniformizante de $V'$, o que prova que se satisfaz a
condição do corolário 5. (Essa possibilidade de utilizar o lema de Abhyankar e
o teorema de pureza foi-me vendida por Serre). Para exprimir o resultado
obtido, introduzamos, para todo grupo compacto totalmente descontínuo $\pi$,
o grupo quociente $\widetilde\pi$ de $\pi$ pelo subgrupo fechado gerado
por seus $p$-subgrupos de Sylow, isto é, o limite projetivo dos grupos
quocientes discretos de $\pi$ que são de ordem prima com $p$. Com esta notação,
obtém-se:

\result{TEOREMA}{13} Seja $f:X\longrightarrow Y$ um morfismo próprio e plano,
sejam $y_1$ um ponto de $Y$ e $y_0$ uma especialização de $y_1$, e suponhamos
$\overline X_0$ conexo e liso. Então o homomorfismo
$\widetilde\pi_1(\overline X_1)\longrightarrow
\widetilde\pi_1(\overline X_0)$ deduzido do homomorfismo sobrejetivo do
teorema 12, corolário 2, é um isomorfismo.

Em outros termos:

\src{219}
\result{COROLÁRIO}{1} A classificação dos recobrimentos galoisianos não
ramificados, de grupo de Galois de ordem prima com a característica $p$ de
$\kappa(y_0)$, é a mesma para $\overline X_0$ e para $\overline X_1$.

Em particular, se $\kappa(y_0)$ tem característica zero, obtém-se por via
algébrica que $\pi_1(\overline X_1)\longrightarrow\pi_1(\overline X_0)$ é
bijetivo.

Assinalemos, por fim, que as técnicas utilizadas proporcionam também o
resultado seguinte, mais geral que o teorema 13:

\result{TEOREMA}{14} Seja $f:X\longrightarrow Y$ um morfismo próprio e liso,
seja $D$ um subesquema fechado de $X$, liso sobre $Y$, de codimensão 1 em
todos os seus pontos. Para uma fibra $Z=f^{-1}(z)$ de $f$, seja $Z'=Z-Z\cap D$ e seja
$\pi_1^t(Z')$ o quociente do grupo fundamental $\pi_1(Z')$ que classifica os
recobrimentos não ramificados de $Z'$ que estão «moderadamente ramificados»
sobre $\overline{Z\cap D}$. Sejam $y_0,y_1$ como no teorema 13. Então se
tem um homomorfismo sobrejetivo (definido salvo um automorfismo interior)
\[
  \pi_1^t(\overline X'_1)\longrightarrow\pi_1^t(\overline X'_0),
\]
e o homomorfismo correspondente
\[
  \widetilde\pi_1^t(\overline X'_1)
  \longrightarrow\widetilde\pi_1^t(\overline X'_0)
\]
é um isomorfismo.

Daí se deduzem as variantes correspondentes dos corolários do
teorema 13 e do corolário 4 do teorema 12. Do mesmo modo, utilizando o
teorema 9, corolário 3, obtém-se por via transcendente:

\result{COROLÁRIO}{} Seja $X_0$ o esquema de uma curva completa lisa sobre um
corpo algebricamente fechado de característica arbitrária, e seja
$S=(s_i)_{1\le i\le n}$ uma parte finita com $n$ elementos de $X_0$. Então
$\pi_1^t(X_0-S)$ admite $2g+n$ geradores topológicos
$x_i,y_i$ ($1\le i\le g$) e $\sigma_j$ ($1\le j\le n$), ligados pela relação
\[
  \left(\prod_i x_i y_i x_i^{-1}y_i^{-1}\right)
  \sigma_1\cdots\sigma_n=1,
\]
onde os $\sigma_j$ são geradores dos grupos de inércia correspondentes
aos $s_j$. Para todo grupo finito $G$ de ordem prima com a característica,
gerado por elementos $x_i,y_i,\sigma_j$ que satisfazem a relação anterior,
existe um recobrimento galoisiano não ramificado de $X_0-S$, de grupo $G$, cujos
grupos de inércia nos pontos $s_j$ estão gerados pelos $\sigma_j$.

Quando $X_0$ é de gênero 0 e $n=3$, tem-se uma solução do «problema dos
três pontos», ao menos para os recobrimentos galoisianos de ordem prima com
a característica. (Aqui, de resto, o teorema 9 é inútil, e parece, por
outra parte, que o corolário anterior pode ser deduzido do caso particular
considerado no problema dos três pontos).

\src{220}
\result{OBSERVAÇÕES}{}

1) Um estudo mais completo, no qual sem dúvida interviriam recobrimentos
galoisianos generalizados de $X,X_0,X_1$ (com grupo de Galois eventualmente
infinitesimal), deveria permitir recuperar o núcleo do teorema 12, corolário
2. Em contrapartida, o estudo dos recobrimentos que admitem ramificação não
«moderada» parece ter de ser muito mais difícil.

2) O lema 6, junto com um resultado de Grauert relativo ao completamento
formal de um esquema projetivo não singular ao longo de uma seção
hiperplana (ou com o teorema, ainda não demonstrado, mencionado na
observação 2 depois do teorema 11), permite também demonstrar em geometria
algébrica «abstrata» o teorema clássico de Lefschetz sobre o grupo
fundamental.

\begin{center}
  \textbf{BIBLIOGRAFIA}
\end{center}

\begin{description}
  \item[{[1]}] Dieudonné (J.) e Grothendieck (A.). --- \emph{Éléments de
    géométrie algébrique}, Publications mathématiques de l'Institut des Hautes
    Études Scientifiques (de próxima publicação).
  \item[{[2]}] Grothendieck (Alexander). --- \emph{The cohomology theory of
    abstract algebraic varieties}. International Congress of Mathematicians
    [1958. Edinburgh] (de próxima publicação).
  \item[{[3]}] Serre (Jean-Pierre). --- \emph{Faisceaux algébriques
    cohérents}, Annals of Math., Series 2, t. 61, 1955, p. 197--278.
  \item[{[4]}] Serre (Jean-Pierre). --- \emph{Géométrie algébrique et
    géométrie analytique}, Ann. Institut Fourier Grenoble, t. 6, 1955--1956,
    p. 1--42.
\end{description}

\par

\def\FgaUnitKey{e182}\clearpage
\chapter*{Errata da exposição n\textsuperscript{o} 182}
\addcontentsline{toc}{chapter}{Errata da exposição n\textsuperscript{o} 182}
\setcounter{footnote}{1}
\markboth{Errata da exposição 182}{Errata da exposição 182}
\sid{e182}
\src{298}

\noindent\textbf{2. Geometria formal e geometria algébrica.}

\begin{center}
[Seminário Bourbaki, t. 11, 1958/59, n\textsuperscript{o} 182, 28 p.]
\end{center}

\noindent\textbf{Página 8, teorema 6.} --- A hipótese «$f$
quase-projetivo» pode ser substituída pela hipótese mais fraca «$f$
separado», em virtude do resultado seguinte (cf. SGA VIII, 6.2): todo morfismo
$f:X\longrightarrow Y$ que seja quase-finito e separado é quase-projetivo.

\src{299}
\noindent\textbf{Página 12, e página 14, observação 1.} --- Assinalemos que J.-P.
Serre construiu uma variedade projetiva não singular, de dimensão 3, sobre
um corpo algebricamente fechado $k$, de característica $p>0$, que não provém
por redução de um esquema próprio sobre um anel local que seja domínio de
integridade, com corpo residual $k$ e corpo de frações de característica
zero.\footnote{Serre (Jean-Pierre). --- Exemples de variétés projectives en
caractéristique $p$ non relevables en caractéristique zéro, \emph{Proc. Nat.
Acad. Sc. U.S.A.}, t. 47, 1961, p. 108--109.} Mumford teria encontrado um exemplo
análogo, com uma superfície projetiva não singular.

\noindent\textbf{Página 15, observações 2 e 3.} --- Atualmente sou menos
otimista quanto aos resultados aqui conjecturados. No entanto, o
problema relativo às variações da estrutura do espaço projetivo,
assinalado ao final da observação 3, foi resolvido afirmativamente por
Hironaka, e o problema análogo para as variedades abelianas foi resolvido
por Koizumi.

\noindent\textbf{Página 16, linha 3.} --- Assinalemos que Mumford construiu
recentemente os esquemas de módulos para as curvas de gênero $g$ (cf. SHMI).
Além disso, o teorema 10 prova que os «esquemas de Jacobi de nível $n$»
da teoria dos módulos são não singulares (e até mesmo lisos sobre $\mathbf Z$).

\noindent\textbf{Página 23, linha 12, primeira linha do corolário
5.} --- Acrescentar: «sendo $X$ ou $Y$ próprio sobre $k$».

A substância dos §§ 1 a 5 está contida na parte publicada de EGA III, e
a dos §§ 6 e 7 está contida em SGA III. Para o estudo do grupo
fundamental, vejam-se SGA V, IX, X, XI, bem como SGA 1962 (exposições X, XII e
XIII), para os teoremas do tipo de Lefschetz e numerosas questões em aberto.
Somente a teoria dos recobrimentos moderadamente ramificados (cf. teorema 14)
ainda não foi objeto de uma redação formal. O corolário do teorema 14,
que fornece a determinação completa dos recobrimentos galoisianos de ordem
coprima à característica de uma curva algébrica sobre um corpo
algebricamente fechado, foi utilizado de maneira essencial em três ocasiões:

1\textsuperscript{o} na demonstração, por Igusa, da desigualdade de Picard
para as superfícies projetivas não singulares em característica arbitrária;

2\textsuperscript{o} no estudo (desenvolvido independentemente por Ogg e
Šafarevič) do grupo dos fibrados principais homogêneos sob uma variedade
abeliana definida sobre um corpo de funções de uma variável, em
característica arbitrária;

3\textsuperscript{o} na demonstração recente, por Artin, de certos
teoremas fundamentais sobre a «cohomologia de Weil» das variedades algébricas.

\par

\def\FgaUnitKey{190}\unit{Exposição 190 -- Técnica de descida I}
\sid{190}
\markboth{Exposição 190}{Exposição 190}
\src{299}
\seminar{12.\textsuperscript{o} ano, 1959/60, n.\textsuperscript{o} 190\\Dezembro de 1959}
\paperhead{TÉCNICA DE DESCIDA E TEOREMAS DE EXISTÊNCIA EM GEOMETRIA ALGÉBRICA\\
I. GENERALIDADES. DESCIDA POR MORFISMOS FIELMENTE PLANOS}%
  {por Alexander GROTHENDIECK}

Do ponto de vista técnico, a presente exposição e as que se seguirão podem
ser consideradas variações sobre o célebre «teorema 90» de Hilbert.
Parece que A. Weil introduziu o método de descida em geometria algébrica sob
o nome de «descida do corpo base». Aliás, Weil restringiu-se ao caso de
extensões finitas separáveis de corpos. P. Cartier tratou depois o caso de
extensões radiciais de altura 1. Por falta da linguagem dos esquemas e,
mais concretamente, por não dispor de elementos nilpotentes nos anéis que
lhe era permitido considerar, Cartier não pudera formular a identidade
essencial entre esses dois casos.

Atualmente, parece que a técnica geral de descida que será exposta (unida,
quando couber, aos teoremas fundamentais da «geometria formal», cf. [3])
está na base da maioria dos teoremas de existência em geometria
algébrica.\footnote{A errata de 1962 considera excessiva essa formulação: ela é
válida em grande medida para as técnicas não projetivas das duas primeiras
exposições da série TDTE I a VI, mas não para as técnicas projetivas (TDTE IV,
V e VI).} Convém assinalar, além disso, que essa técnica pode, sem dúvida, ser
transposta à «geometria analítica», e cabe esperar que, num futuro próximo,
os especialistas saibam demonstrar os análogos «analíticos» dos teoremas de
existência em geometria formal que serão apresentados na exposição II. Em
todo caso, os trabalhos recentes de Kodaira--Spencer, cujos métodos parecem
inadequados para definir e estudar «variedades de módulos» numa vizinhança
de seus pontos singulares, mostram com bastante clareza a necessidade de
métodos mais próximos da teoria dos esquemas (que deverão completar,
naturalmente, as técnicas transcendentes).

Na presente exposição I trataremos o caso de descida mais elementar (indicado
no título). No entanto, as aplicações dos teoremas 1, 2 e 3 que se seguem já
são muito numerosas. Limitamo-nos a apresentar algumas a título de exemplo,
sem procurar conferir-lhes aqui a máxima generalidade possível.

Utilizaremos livremente a linguagem dos esquemas, para a qual remetemos à
exposição já citada, bem como a [2]. Assinalemos, contudo, expressamente que os
pré-esquemas considerados na presente exposição não são necessariamente
noetherianos, e os

\src{300}
morfismos não são necessariamente de tipo finito. Assim, se $A$ é um anel local
noetheriano, com completamento $\widehat A$, será preciso considerar o anel não
noetheriano $\widehat A\otimes_A\widehat A$, e os morfismos de esquemas afins
correspondentes às inclusões entre os anéis considerados.

\fsection{A}{Preliminares sobre as categorias}

\subsection*{1. Categorias fibradas, dados de descida, morfismos de
$\mathcal F$-descida}
\addcontentsline{toc}{subsection}{1. Categorias fibradas e dados de descida}

\result{a. DEFINIÇÃO}{1.1} Uma \emph{categoria fibrada} $\mathcal F$ de base
$\mathcal C$ consta de uma categoria $\mathcal C$; para cada
$X\in\mathcal C$, de uma categoria $\mathcal F_X$; para cada
$\mathcal C$-morfismo $f:X\longrightarrow Y$, de um funtor
$f^*:\mathcal F_Y\longrightarrow\mathcal F_X$, denotado também por
\[
  f^*(\xi)=\xi\times_Y X
\]
para $\xi\in\mathcal F_Y$ (considerando $X$ como um «objeto de
$\mathcal C$ sobre $Y$», isto é, provido do morfismo $f$); e, por fim,
para dois morfismos componíveis
$X\xrightarrow{f}Y\xrightarrow{g}Z$, de um isomorfismo de funtores
\[
  c_{f,g}:(gf)^*\longrightarrow f^*g^*,
\]
estando esses dados sujeitos às condições seguintes:

\begin{enumerate}
\item[(i)] $\operatorname{id}_X^*=\operatorname{id}$;
\item[(ii)] $c_{f,g}$ é o isomorfismo identidade se $f$ ou $g$ é um
isomorfismo identidade;
\item[(iii)] dados três morfismos componíveis
$X\xrightarrow{f}Y\xrightarrow{g}Z\xrightarrow{h}T$, o seguinte diagrama de
isomorfismos, formado mediante isomorfismos da forma $c_{u,v}$, é
comutativo:
\[
\begin{array}{ccc}
  (h(gf))^*=((hg)f)^* & \longrightarrow & f^*(hg)^* \\
  \big\downarrow && \big\downarrow \\
  (gf)^*h^* & \longrightarrow & (f^*g^*)h^*=f^*(g^*h^*).
\end{array}
\]
\end{enumerate}

\result{EXEMPLO}{1} Seja $\mathcal C$ uma categoria na qual existem os
produtos fibrados. Define-se então uma categoria fibrada $\mathcal F$ de
base $\mathcal C$ tomando $\mathcal F_X=$ categoria dos objetos de
$\mathcal C$ sobre $X$; o funtor
$f^*:\mathcal F_Y\longrightarrow\mathcal F_X$ correspondente a um morfismo
$f:X\longrightarrow Y$ define-se por meio do produto fibrado
$Z\longmapsto Z\times_YX$.

\result{EXEMPLO}{2} Seja $\mathcal C$ a categoria dos pré-esquemas e, para
$X\in\mathcal C$, seja $\mathcal F_X$ a categoria dos feixes
quase-coerentes de módulos sobre $X$. Se $f:X\longrightarrow Y$ é um morfismo
de pré-esquemas, $f^*:\mathcal F_Y\longrightarrow\mathcal F_X$ é o funtor
imagem inversa

\src{301}
de feixes de módulos. Obtém-se assim uma categoria fibrada de base
$\mathcal C$.

\result{b. DEFINIÇÃO}{1.2} Um diagrama
\[
  E\xrightarrow{u}E'
  \overset{v_1}{\underset{v_2}{\rightrightarrows}}E''
\]
de aplicações de conjuntos chama-se \emph{exato} se $u$ é uma bijeção de
$E$ sobre a parte de $E'$ formada pelos $x'\in E'$ tais que
$v_1(x')=v_2(x')$.

\result{DEFINIÇÃO}{1.3} Seja $\mathcal F$ uma categoria fibrada de base
$\mathcal C$; consideremos um diagrama
\[
  \tag{+}
  S\xleftarrow{\alpha}S'
  \overset{\beta_1}{\underset{\beta_2}{\rightrightarrows}}S''
\]
de morfismos em $\mathcal C$ tais que
$\alpha\beta_1=\alpha\beta_2$; esse diagrama chama-se
$\mathcal F$-\emph{exato} se, para cada par $\xi,\eta$ de elementos de
$\mathcal F_S$, o seguinte diagrama de aplicações de conjuntos
\[
  \tag{+'}
  \operatorname{Hom}(\xi,\eta)
  \xrightarrow{\alpha^*}
  \operatorname{Hom}(\alpha^*\xi,\alpha^*\eta)
  \overset{\beta_1^*}{\underset{\beta_2^*}{\rightrightarrows}}
  \operatorname{Hom}(\gamma^*\xi,\gamma^*\eta),
\]
onde $\gamma=\alpha\beta_1=\alpha\beta_2$, é exato. Neste último diagrama
temos identificado $\beta_i^*\alpha^*$ com
$(\alpha\beta_i)^*=\gamma^*$ mediante $c_{\beta_i,\alpha}$, para simplificar.

\result{DEFINIÇÃO}{1.4} Seja $\mathcal F$ uma categoria fibrada de base
$\mathcal C$; consideremos dois morfismos $\beta_1,\beta_2$ de $S''$ em $S'$,
em $\mathcal C$. Seja
$\xi'\in\mathcal F_{S'}$. Chama-se \emph{dado de colagem} sobre $\xi'$
(relativo ao par $(\beta_1,\beta_2)$) a um isomorfismo de
$\beta_1^*(\xi')$ sobre $\beta_2^*(\xi')$. Se $\xi',\eta'\in\mathcal F_{S'}$
estão munidos, cada um, de um dado de colagem, um morfismo
$u:\xi'\longrightarrow\eta'$ em $\mathcal F_{S'}$ se diz
\emph{compatível com os dados de colagem} se o diagrama seguinte é
comutativo:
\[
\begin{array}{ccc}
  \beta_1^*(\xi') & \longrightarrow & \beta_2^*(\xi') \\
  \big\downarrow && \big\downarrow \\
  \beta_1^*(\eta') & \longrightarrow & \beta_2^*(\eta').
\end{array}
\]
Desse modo, os objetos de $\mathcal F_{S'}$ providos de um dado de colagem
(relativo a $\beta_1,\beta_2$) formam uma categoria. Se
$\alpha:S'\longrightarrow S$ é um morfismo tal que
$\alpha\beta_1=\alpha\beta_2$, então, para todo $\xi\in\mathcal F_S$, o
objeto $\xi'=\alpha^*(\xi)$

\src{302}
de $\mathcal F_{S'}$ está provido de um dado de colagem canônico, visto que
\[
  \beta_i^*\alpha^*(\xi)\simeq(\alpha\beta_i)^*(\xi)=\gamma^*(\xi),
\]
tomando de novo $\gamma=\alpha\beta_1=\alpha\beta_2$; além disso, se
$u:\xi\longrightarrow\eta$ é um morfismo em $\mathcal F_S$, então
\[
  \alpha^*(u):\alpha^*(\xi)\longrightarrow\alpha^*(\eta)
\]
é um morfismo em $\mathcal F_{S'}$ compatível com os dados de colagem
canônicos. Obtém-se assim um funtor canônico da categoria $\mathcal F_S$ para
a categoria dos objetos de $\mathcal F_{S'}$ providos de um dado de colagem
relativo ao par $(\beta_1,\beta_2)$. Dito isto, a definição 1.3 também
pode ser expressa dizendo que o diagrama (+) dessa definição é
$\mathcal F$-exato se o funtor anterior é «plenamente fiel», isto é, se
define uma equivalência da categoria $\mathcal F_S$ com uma subcategoria da
categoria dos objetos de $\mathcal F_{S'}$ providos de um dado de colagem
relativo a $(\beta_1,\beta_2)$.

\result{DEFINIÇÃO}{1.5} Diz-se que um dado de colagem sobre
$\xi'\in\mathcal F_{S'}$ é \emph{efetivo} (com respeito a $\alpha$) se
$\xi'$, provido deste dado, é isomorfo a um $\alpha^*(\xi)$, com
$\xi\in\mathcal F_S$.

Quando o diagrama (+) é $\mathcal F$-exato, o objeto $\xi$ da definição
anterior fica, portanto, determinado salvo isomorfismo único, e a categoria
$\mathcal F_S$ é equivalente à categoria dos objetos de
$\mathcal F_{S'}$ providos de um dado de colagem efetivo.

\medskip\noindent
c. O caso mais importante é aquele em que
\[
  S''=S'\times_S S',
\]
sendo as $\beta_i$ as duas projeções $p_1$ e $p_2$ de
$S'\times_S S'$ sobre seus dois fatores (supomos agora que em
$\mathcal C$ existem os produtos fibrados). Têm-se então duas
condições necessárias imediatas para que seja efetivo um dado de colagem
$\varphi:p_1^*(\xi')\longrightarrow p_2^*(\xi')$ sobre um
$\xi'\in\mathcal F_{S'}$. Tem-se
\[
  \tag{i}
  \Delta^*(\varphi)=\operatorname{id}_{\xi'},
\]
onde $\Delta:S'\longrightarrow S'\times_S S'$ designa o morfismo diagonal,
e onde se identificou $\Delta^*p_i^*(\xi')$ com
$(p_i\Delta)^*(\xi')=\xi'$. Tem-se
\[
  \tag{ii}
  p_{31}^*(\varphi)=p_{32}^*(\varphi)\,p_{21}^*(\varphi),
\]

\src{303}
onde $p_{ij}$ designa a projeção canônica de
$S'\times_S S'\times_S S'$ sobre o produto parcial de seus fatores
$i$-ésimo e $j$-ésimo.

\result{DEFINIÇÃO}{1.6} Chama-se \emph{dado de descida} sobre
$\xi'\in\mathcal F_{S'}$, relativo ao morfismo
$\alpha:S'\longrightarrow S$, a um dado de colagem sobre $\xi'$ relativo ao par
$(p_1,p_2)$ de projeções canônicas
$S'\times_S S'\longrightarrow S'$ que satisfaz as condições (i) e (ii)
anteriores.

\result{DEFINIÇÃO}{1.7} Um morfismo $\alpha:S'\longrightarrow S$ chama-se
\emph{morfismo de $\mathcal F$-descida} se o diagrama de morfismos
\[
  S\xleftarrow{\alpha}S'
  \overset{p_1}{\underset{p_2}{\rightrightarrows}}S'\times_S S'
\]
é $\mathcal F$-exato (definição 1.3); diz-se que $\alpha$ é um
\emph{morfismo de $\mathcal F$-descida estrito} se, além disso, todo dado de
descida (definição 1.6) sobre um objeto de $\mathcal F_{S'}$ é efetivo.

Esta última condição também pode ser enunciada, de modo mais gráfico, assim:
«Dá no mesmo» dar um objeto de $\mathcal F_S$ ou um objeto de
$\mathcal F_{S'}$ provido de um dado de descida.

Observe-se que se um morfismo $\alpha:S'\longrightarrow S$ admite uma «seção»
$s:S\longrightarrow S'$ (isto é, um morfismo $s$ tal que
$\alpha s=\operatorname{id}_S$), então é um morfismo de
$\mathcal F$-descida estrito.\footnote{A errata de 1962 esclarece que não é
necessário supor previamente que $\alpha$ seja um morfismo de
$\mathcal F$-descida.} Com efeito, se $\xi'\in\mathcal F_{S'}$ está provido de
um dado de descida relativo a $\alpha$, «provém» de
$\xi=s^*(\xi')$.

\medskip\noindent
d. As noções anteriores podem ser apresentadas de maneira mais intuitiva
introduzindo, para um objeto parâmetro $T$ de $\mathcal C$ sobre $S$, o
conjunto
\[
  \operatorname{Hom}_S(T,S')=S'(T),
\]
cujos elementos serão denotados por $t,t'$, etc. Dado um objeto
$\xi'\in\mathcal F_{S'}$, a cada $t\in S'(T)$ corresponde então um objeto
$t^*(\xi')$ de $\mathcal F_T$, que também será denotado por $\xi'_t$. Um dado de
colagem sobre $\xi'$ (relativo a $p_1,p_2$) fica definido então dando, para
todo $T$ sobre $S$ e todo par de pontos $t,t'\in S'(T)$, um isomorfismo
\[
  \varphi_{t',t}:\xi'_t\longrightarrow\xi'_{t'}
\]
(que satisfaz condições evidentes de funtorialidade ao variar $T$). As
condições (i) e (ii) de 1.c escrevem-se então
\[
  \tag{i bis}
  \varphi_{t,t}=\operatorname{id}_{\xi'_t}
  \qquad\text{para todo $T$ e todo $t\in S'(T)$},
\]

\src{304}
\[
  \tag{ii bis}
  \varphi_{t,t''}=\varphi_{t,t'}\varphi_{t',t''}
  \qquad\text{para todo $T$ e todos $t,t',t''\in S'(T)$}.
\]
Vê-se, além disso, que (ii bis), tomando $t=t'=t''$, implica
$\varphi_{t,t}=\varphi_{t,t}^2$; visto que $\varphi_{t,t}$ é um isomorfismo
por hipótese, daí se obtém a relação (i bis), que é de fato uma consequência
de (ii bis) (portanto, (i) é consequência de (ii)). Mas, se já não se supõe
a priori que os $\varphi_{t,t'}$ sejam isomorfismos (isto é, que
$\varphi:p_1^*(\xi')\longrightarrow p_2^*(\xi')$ seja um isomorfismo), então
(ii bis) não implica necessariamente (i bis); no entanto, a conjunção de
(ii bis) e (i bis) implica que os $\varphi_{t,t'}$ são isomorfismos (pois se
terá
$\varphi_{t,t'}\varphi_{t',t}=\varphi_{t,t}=\operatorname{id}_{\xi'_t}$).

\subsection*{2. Diagramas exatos e epimorfismos estritos, morfismos de
descida. Exemplos}
\addcontentsline{toc}{subsection}{2. Diagramas exatos e epimorfismos estritos}

\result{a. DEFINIÇÃO}{2.1} Seja $\mathcal C$ uma categoria. Um diagrama
\[
  T\xrightarrow{\alpha}T'
  \overset{\beta_1}{\underset{\beta_2}{\rightrightarrows}}T''
\]
de morfismos em $\mathcal C$ chama-se \emph{exato} se, para todo
$Z\in\mathcal C$, o diagrama correspondente de aplicações de conjuntos
\[
  \operatorname{Hom}(Z,T)\longrightarrow
  \operatorname{Hom}(Z,T')
  \rightrightarrows\operatorname{Hom}(Z,T'')
\]
é exato (definição 1.2). Diz-se então que $(T,\alpha)$ (ou, por abuso de
linguagem, que $T$) é um \emph{núcleo} do par de morfismos
$(\beta_1,\beta_2)$.

Evidentemente, esse núcleo está determinado salvo isomorfismo único. Se
$\mathcal C$ é a categoria dos conjuntos, a definição anterior é
compatível com a definição 1.2. Define-se de maneira dual a exatidão de um
diagrama
\[
  S\xleftarrow{\alpha}S'
  \overset{\beta_1}{\underset{\beta_2}{\rightrightarrows}}S''
\]
de morfismos em $\mathcal C$; diz-se então que $(S,\alpha)$ é um
\emph{conúcleo} do par de morfismos $(\beta_1,\beta_2)$.

\result{DEFINIÇÃO}{2.2} Um morfismo $\alpha:S'\longrightarrow S$ chama-se
\emph{epimorfismo estrito} se é um epimorfismo e se, para todo morfismo
$u:S'\longrightarrow Z$, a seguinte condição necessária é também
suficiente para que $u$ se fatore como
$S'\longrightarrow S\longrightarrow Z$: para todo $S''\in\mathcal C$ e todo
par de morfismos $\beta_1,\beta_2:S''\longrightarrow S'$ tal que
$\alpha\beta_1=\alpha\beta_2$, tem-se também $u\beta_1=u\beta_2$.

Se existe o produto fibrado $S'\times_S S'$, equivale a dizer que o diagrama
\[
  S\xleftarrow{\alpha}S'
  \overset{p_1}{\underset{p_2}{\rightrightarrows}}S'\times_S S'
\]

\src{305}
é exato, isto é, que $S$ é um conúcleo do par $(p_1,p_2)$. Em todo
caso, um morfismo conúcleo é um epimorfismo estrito. Assinalemos também que um
epimorfismo estrito que seja um monomorfismo é um isomorfismo. Deixamos ao
leitor desenvolver a noção dual de \emph{monomorfismo estrito}.

Para precisar as relações entre a noção de morfismo de
$\mathcal F$-descida e a de epimorfismo estrito, introduzimos além disso as
definições seguintes:

\result{DEFINIÇÃO}{2.3} Um morfismo $\alpha:S'\longrightarrow S$ chama-se
\emph{epimorfismo universal} (resp. \emph{epimorfismo estrito universal}) se,
para todo $T$ sobre $S$, existe o produto fibrado $T'=S'\times_ST$ e a
projeção $T'\longrightarrow T$ é um epimorfismo (resp. um epimorfismo
estrito).

Nas categorias muito boas (como a categoria dos conjuntos, a
categoria dos feixes de conjuntos sobre um espaço topológico, as categorias
abelianas, etc.), as quatro noções de epimorfismo assim introduzidas
coincidem; ao contrário, todas elas são distintas numa categoria como
a dos pré-esquemas, ou a dos pré-esquemas sobre um pré-esquema não vazio
dado $S$, inclusive se nos limitamos a $S$-esquemas finitos sobre $S$.

\result{DEFINIÇÃO}{2.4} Um morfismo $\alpha:S'\longrightarrow S$ chama-se
\emph{morfismo de descida} (resp. \emph{morfismo de descida estrito}) se é
um morfismo de $\mathcal F$-descida (resp. um morfismo de
$\mathcal F$-descida estrito) (cf. definição 1.7), onde $\mathcal F$ denota
a categoria fibrada de base $\mathcal C$ dos objetos de $\mathcal C$ sobre
objetos de $\mathcal C$ (n\textsuperscript{o} 1, exemplo 1).

\result{PROPOSIÇÃO}{2.1} Se em $\mathcal C$ existem os produtos finitos e
os produtos fibrados (finitos), então, em $\mathcal C$, os morfismos de
descida coincidem com os epimorfismos estritos universais.

\result{b. EXEMPLOS}{} Seja $\mathcal C$ a categoria dos pré-esquemas. Seja
$S\in\mathcal C$, e sejam $S'$ e $S''$ dois pré-esquemas finitos sobre $S$,
isto é, correspondentes a feixes de álgebras $\mathcal A'$, $\mathcal A''$
sobre $S$ que, como feixes de módulos, são quase-coerentes e de tipo finito
(isto é, coerentes se $S$ é localmente noetheriano). Seja
$\alpha:S'\longrightarrow S$ o morfismo estrutural de $S'$, e sejam
$\beta_1,\beta_2$ dois $S$-morfismos de $S''$ em $S'$, definidos por
homomorfismos de álgebras $\mathcal A'\longrightarrow\mathcal A''$, denotados
também por $\beta_1,\beta_2$. Utilizando o fato de que um morfismo finito é
fechado (primeiro teorema de Cohen--Seidenberg), demonstra-se facilmente que o
diagrama
\[
  \tag{+}
  S\xleftarrow{\alpha}S'
  \overset{\beta_1}{\underset{\beta_2}{\rightrightarrows}}S''
\]
em $\mathcal C$ é exato se e somente se também o é o diagrama de feixes

\src{306}
\[
  \mathcal O_S=\mathcal A\longrightarrow\mathcal A'
  \overset{\beta_1}{\underset{\beta_2}{\rightrightarrows}}\mathcal A''
\]
sobre $S$. Em particular, se $\alpha:S'\longrightarrow S$ é um morfismo
finito correspondente a um feixe $\mathcal A'$ de álgebras sobre $S$, então
$\alpha$ é um epimorfismo estrito se e somente se o diagrama de feixes
\[
  \mathcal O_S=\mathcal A\longrightarrow\mathcal A'
  \overset{p_1}{\underset{p_2}{\rightrightarrows}}
  \mathcal A'\otimes_{\mathcal A}\mathcal A'
\]
é exato (é um epimorfismo se e somente se
$\mathcal A\longrightarrow\mathcal A'$ é injetivo). Se $S$ é afim de anel
$A$, e portanto $S'$ é afim de anel $A'$ finito sobre $A$, então
$S'\longrightarrow S$ é um epimorfismo estrito se e somente se
$A\longrightarrow A'$ é um isomorfismo de $A$ sobre o subanel de $A'$
formado pelos $x'\in A'$ tais que
\[
  1_{A'}\otimes_Ax'-x'\otimes_A1_{A'}=0
\]
(é um epimorfismo se e somente se $A\longrightarrow A'$ é injetivo). Como já
assinalamos, mesmo se $S$ é o esquema de um anel local artiniano, um
morfismo finito $S'\longrightarrow S$ que seja um epimorfismo não é
necessariamente um epimorfismo estrito. No entanto, pode-se demonstrar que, se
$S$ é um pré-esquema noetheriano, todo morfismo finito
$S'\longrightarrow S$ que seja um epimorfismo é a composição de uma sucessão
finita de epimorfismos estritos (também finitos). Isso mostra, além disso, que
a composição de dois epimorfismos estritos não é necessariamente um
epimorfismo estrito.

\result{c}{} Se $(+)$ é um diagrama exato de morfismos finitos, então,
para todo morfismo plano $T\longrightarrow S$ de pré-esquemas, o diagrama
obtido de $(+)$ mediante a mudança de base $T\longrightarrow S$ continua
exato. Disso se deduz que, se $X,Y$ são dois $S$-pré-esquemas, sendo $X$
plano sobre $S$, então o seguinte diagrama de aplicações de conjuntos
(onde $X',Y'$ são as imagens inversas de $X,Y$ sobre $S'$, e $X'',Y''$ suas
imagens inversas sobre $S''$)
\[
  \operatorname{Hom}_S(X,Y)\longrightarrow
  \operatorname{Hom}_{S'}(X',Y')
  \rightrightarrows
  \operatorname{Hom}_{S''}(X'',Y'')
\]
é exato. Em particular, se $\mathcal F$ denota a categoria fibrada cuja base
é a categoria $\mathcal C$ dos pré-esquemas e tal que, para
$X\in\mathcal C$, $\mathcal F_X$ é a categoria dos $X$-pré-esquemas
planos, então o diagrama $(+)$ é $\mathcal F$-exato. (Esse resultado deixa
de ser verdadeiro se não se impõe a hipótese de planicidade; em particular, um
epimorfismo estrito finito não é necessariamente um morfismo de descida.)
Vê-se igualmente que $(+)$ é $\mathcal F$-exato se $\mathcal F$ denota a
categoria fibrada para a qual $\mathcal F_X$ é a categoria dos feixes
quase-coerentes e planos sobre o pré-esquema $X$ (também aqui a hipótese

\src{307}
de planicidade é essencial). Num e noutro caso, a questão da efetividade de
um dado de colagem (e, mais concretamente, de um dado de descida quando
$S''=S'\times_SS'$) sobre um objeto plano sobre $S'$ é delicada (e sua
resposta em diversos casos particulares é um dos objetos principais das
presentes exposições). O conferenciante ignora se, para todo epimorfismo
estrito finito $S'\longrightarrow S$, todo dado de descida sobre um feixe
quase-coerente plano sobre $S'$ é efetivo (mesmo supondo que $S$ seja o
espectro de um anel local artiniano e limitando-nos aos feixes localmente
livres de posto 1). Mais geralmente, sejam $A$ um anel e $A'$ uma
$A$-álgebra (tudo é comutativo) tal que o diagrama de aplicações
\[
  A\longrightarrow A'\rightrightarrows A'\otimes_AA'
\]
seja exato, o que equivale também a que o morfismo
$S'\longrightarrow S$ correspondente entre os espectros de $A,A'$ seja um
morfismo de $\mathcal F$-descida, onde $\mathcal F$ é a categoria fibrada
dos feixes quase-coerentes planos. Seja $M'$ um $A'$-módulo plano provido de
um dado de descida a $A$, isto é, de um isomorfismo
\[
  \varphi:M'\otimes_AA'\xrightarrow{\sim}A'\otimes_AM'
\]
de $A'\otimes_AA'$-módulos que satisfaz as condições (i) e (ii) de 1.c
(cuja explicitação em termos de módulos deixamos ao leitor). É efetivo
este dado (com respeito à categoria fibrada dos feixes quase-coerentes
planos)? Seja $M$ o subconjunto de $M'$ formado pelos $x'\in M'$ tais que
\[
  \varphi(x'\otimes_A1_{A'})=1_{A'}\otimes_Ax',
\]
que é um sub-$A$-módulo de $M'$. A injeção canônica
$M\longrightarrow M'$ define um homomorfismo de $A'$-módulos
$M\otimes_AA'\longrightarrow M'$. A efetividade de $\varphi$ significa
então o seguinte: $M$ é um $A$-módulo plano e o homomorfismo anterior é
um isomorfismo.

\result{OBSERVAÇÃO}{} Nas considerações anteriores não havíamos imposto
nenhuma hipótese de planicidade aos morfismos do diagrama $(+)$, o que nos
obrigava, para dispor de uma técnica de descida, a impor hipóteses de
planicidade aos objetos sobre $S,S'$ considerados. No parágrafo 2 imporemos
uma hipótese de planicidade a $\alpha:S'\longrightarrow S$, o que nos permitirá
dispor de uma técnica de descida para objetos sobre $S,S'$ que já não estão
submetidos a condição alguma de planicidade. Em todos os casos intervém uma
hipótese de planicidade. Essa é uma das principais razões da importância da
noção de planicidade em geometria algébrica (cujo papel não podia manifestar-se
enquanto nos limitávamos a corpos base, sobre os quais, com efeito, tudo é
plano!).

\src{308}
\subsection*{3. Aplicação aos morfismos étale}
\addcontentsline{toc}{subsection}{3. Aplicação aos morfismos étale}

Sejam $A$ um anel local e $B$ um álgebra local sobre $A$ cujo ideal maximal
induz o de $A$. Diremos que $B$ é \emph{étale} sobre $A$ (em lugar de
«não ramificada», expressão usada em outros lugares) se satisfaz as condições
seguintes:
\begin{enumerate}
\item[(i)] $B$ é plana sobre $A$;
\item[(ii)] $B/\mathfrak mB$ é uma extensão finita separável de
$A/\mathfrak m=k$ (onde $\mathfrak m$ denota o ideal maximal de $A$).
\end{enumerate}
Quando $A$ e $B$ são noetherianos e $k$ é algebricamente fechado, isso
significa que o homomorfismo $\widehat A\longrightarrow\widehat B$ entre os
completamentos que prolonga $A\longrightarrow B$ é um isomorfismo. Um morfismo de
tipo finito $f:T\longrightarrow S$ diz-se \emph{étale em $x\in T$}, ou também
que $T$ é \emph{étale sobre $S$ em $x$}, se $\mathcal O_x$ é étale sobre
$\mathcal O_{f(x)}$; diz-se que $f$ é \emph{étale}, ou chama-se $f$ um
\emph{morfismo étale}, ou diz-se que $T$ é \emph{étale sobre $S$}, se $f$ é
étale em todo $x\in T$. Assinalemos, além disso, que se $S$ é localmente
noetheriano, o conjunto de pontos de $T$ nos quais $f$ é étale é aberto, e
o uso do «main theorem» de Zariski permite precisar a estrutura de $T/S$
numa vizinhança de tal ponto (por meio de uma equação de tipo bem conhecido).

Quando $S$ é um esquema de tipo finito sobre o corpo dos números
complexos, corresponde-lhe um espaço analítico $\overline S$ no sentido de
Serre [5], exceto que $\overline S$ pode conter elementos nilpotentes em seu
feixe estrutural, o que em nada altera o essencial de [5]. Vê-se então facilmente
que $f$ é étale se e somente se
$\overline f:\overline T\longrightarrow\overline S$ também o é, isto é, se todo
ponto de $\overline T$ admite uma vizinhança na qual $\overline f$ induz um
isomorfismo sobre um aberto de $\overline S$. Em particular, a todo
\emph{recobrimento étale} $T$ de $S$ (isto é, um morfismo finito étale
$f:T\longrightarrow S$) corresponde um recobrimento étale $\overline T$ de
$\overline S$, que é conexo se e somente se $T$ também o é [5]. Vê-se, além
disso, facilmente que, se $T,T'$ são dois esquemas étale sobre $S$, então a
aplicação natural
\[
  \operatorname{Hom}_S(T,T')\longrightarrow
  \operatorname{Hom}_{\overline S}(\overline T,\overline{T'})
\]
é bijetiva, isto é, o funtor $T\longmapsto\overline T$ da categoria dos
esquemas étale sobre $S$ à categoria dos espaços analíticos étale
sobre $\overline S$ é «plenamente fiel» e, portanto, define uma equivalência
da primeira categoria com uma subcategoria da segunda. Um teorema de
Grauert--Remmert [2] implica que, se $S$ é normal, obtém-se assim uma
equivalência entre a categoria dos \emph{recobrimentos étale} de $S$ e a
categoria dos recobrimentos étale (finitos) de $\overline S$, isto é, que
todo recobrimento étale de $\overline S$ é $\overline S$-isomorfo a um
$\overline T$, onde $T$ é um recobrimento étale de $S$. Demonstremos que o
\emph{teorema de Grauert--Remmert continua válido sem hipótese de
normalidade sobre $S$}. Seja, com efeito, primeiro $S'\longrightarrow S$ um
epimorfismo estrito finito; suponhamos demonstrado o

\src{309}
teorema para $S'$ e demonstremos que será verdadeiro para $S$. Seja
$\mathfrak T$ um recobrimento étale de $\overline S$ e consideremos sua imagem
inversa $\mathfrak T'$ sobre $\overline{S'}$, que corresponde a um feixe
analítico coerente $\mathcal A'$ de álgebras sobre $\overline{S'}$, imagem
inversa do feixe de álgebras $\mathcal A$ sobre $\overline S$ que define
$\mathfrak T$. Por hipótese, sobre $\overline{S'}$, $\mathfrak T'$ provém de
um recobrimento étale $T'$ de $S'$, isto é, $\mathcal A'$ provém de um feixe
coerente de álgebras $A'$ sobre $S'$. Por outro lado, $\mathcal A'$ está
provido de um dado de descida canônico relativo a
$\overline{S'}\longrightarrow\overline S$, isto é, de um isomorfismo entre
suas duas imagens inversas sobre
\[
  \overline{S'}\times_{\overline S}\overline{S'}
  =\overline{(S'\times_SS')},
\]
(que satisfaz as condições (i), (ii)), e esse isomorfismo provém, conforme
dito, de um isomorfismo entre os feixes algébricos correspondentes, isto é,
de um dado de descida sobre $A'$ relativo a
$S'\longrightarrow S$. Verifica-se facilmente que este último é efetivo
(pois o de $\mathcal A'$ também o é, e a efetividade de um dado de descida,
tal como se explicitou no número anterior, reconhece-se localmente nos
\emph{completamentos} dos módulos que intervêm). Obtém-se assim um feixe
coerente de álgebras $A$ sobre $S$, que define um recobrimento $T$ de $S$,
o recobrimento procurado. O resultado anterior continua então manifestamente
válido se $S'\longrightarrow S$ for apenas uma composição de um número finito
de epimorfismos estritos finitos, isto é, se for um epimorfismo finito
qualquer (pelo resultado de fatoração assinalado no parágrafo 2). Segue-se que
o teorema de Grauert--Remmert continua válido se $S$ é um esquema
\emph{reduzido}, isto é, se $\mathcal O_S$ não tem elementos nilpotentes, como
se vê introduzindo sua normalização $S'$. Daí se passa facilmente ao caso geral.

Uma demonstração inteiramente análoga, que utiliza de novo o resultado de
fatoração para os epimorfismos estritos finitos e a natureza «formal» da
efetividade dos dados de descida, permite demonstrar o resultado seguinte:
seja $S$ um pré-esquema localmente noetheriano e seja
$S'\longrightarrow S$ um morfismo \emph{finito, sobrejetivo e radicial} (ou,
o que dá no mesmo, um morfismo de tipo finito tal que, para todo $T$
sobre $S$, o morfismo $T'=S'\times_ST\longrightarrow T$ seja um homeomorfismo,
o que também se exprime dizendo que $S'\longrightarrow S$ é um
\emph{homeomorfismo universal}). Para todo $T$ étale sobre $S$, consideremos
sua imagem inversa $T'=T\times_SS'$, que é étale sobre $S'$. Então o
funtor $T\longmapsto T'$ é uma \emph{equivalência de categorias entre a
categoria dos pré-esquemas $T$ étale sobre $S$ e a categoria dos
pré-esquemas $T'$ étale sobre $S'$}. (Utiliza-se a bijetividade de
\[
  \operatorname{Hom}_S(T_1,T_2)\longrightarrow
  \operatorname{Hom}_{S'}(T'_1,T'_2)
\]
para dois pré-esquemas $T_1,T_2$ étale sobre $S$, fato cuja verificação direta
é fácil, e o fato de que o teorema enunciado é verdadeiro se
$S'=(S,\mathcal O_S/\mathcal J)$,

\src{310}
onde $\mathcal J$ é um feixe coerente nilpotente de ideais de $\mathcal O_S$
([4], lema 6)). Observe-se, além disso, que não supomos aqui que os $T$
considerados sejam finitos sobre $S$. Este resultado implica, em particular,
que o morfismo $S'\longrightarrow S$ induz um isomorfismo do grupo
fundamental $[\,]$ de $S'$ sobre o de $S$ («\emph{invariância topológica do
grupo fundamental de um pré-esquema}»).

\subsection*{4. Relações com a 1-cohomologia}
\addcontentsline{toc}{subsection}{4. Relações com a 1-cohomologia}

\medskip\noindent
a. Seja $\mathcal C$ uma categoria na qual sempre existe o produto de dois
objetos, e seja $T\in\mathcal C$. Para todo conjunto finito
$I\ne\varnothing$, pode-se considerar $T^I$; ao variar $I$, obtém-se assim um
funtor covariante da categoria dos conjuntos finitos não vazios para
$\mathcal C$, isto é, o que se pode chamar um \emph{objeto simplicial} de
$\mathcal C$, denotado por $K_T$. Este último depende covariantemente de $T$;
além disso, se $u,v$ são dois morfismos $T\longrightarrow T'$, os morfismos
correspondentes $K_T\longrightarrow K_{T'}$ são homotópicos. Diremos que $T$
\emph{domina} a $T'$ se $\operatorname{Hom}(T,T')\ne\varnothing$; esta é uma
relação de pré-ordem filtrante crescente em $\mathcal C$. Do anterior se
deduz que, se $T$ domina $T'$, existe uma classe canônica (salvo homotopia)
de homomorfismos de objetos simpliciais $K_T\longrightarrow K_{T'}$; em
particular, se $K_T$ e $K_{T'}$ são tais que cada um domina ao outro, então
$K_T$ e $K_{T'}$ são homotopicamente equivalentes. Seja agora $F$ um funtor
(contravariante, para fixar ideias) de $\mathcal C$ para uma categoria abeliana
$\mathcal C'$; então
\[
  C^*(T,F)=F(K_T)
\]
é um objeto cosimplicial de $\mathcal C'$ e, portanto, define da maneira
habitual um complexo (de cocadeias) em $\mathcal C'$, do qual se pode tomar a
cohomologia
\[
  H^*(T,F)=H^*(C^*(T,F))=H^*(F(K_T))
\]
(poder-se-á escrever uma $\mathcal C$ como subíndice de $H^*$ se houver
possibilidade de confusão). Esse é um funtor cohomológico em $F$, cuja
variância ao variar $T$ se deduz do que foi dito sobre os $K_T$; mais
precisamente, para $F$ fixo e $T$ variável em $\mathcal C$ (pré-ordenada pela
relação de dominação), os $H^*(T,F)$ formam um sistema indutivo de objetos
graduados de $\mathcal C'$;
em particular, se $T$ e $T'$ são tais que cada um domina ao outro, então
$H^*(T,F)$ e $H^*(T',F)$ são canonicamente isomorfos.

Suponhamos que em $\mathcal C$ existem os produtos fibrados. Para
$S\in\mathcal C$ fixo, pode-se aplicar o anterior à categoria
$\mathcal C_S$ dos objetos de $\mathcal C$ sobre $S$; escreveremos
$C^*(T/S,F)$ e $H^*(T/S,F)$ em lugar de $C^*(T,F)$ e $H^*(T,F)$ quando
quisermos precisar que trabalhamos na categoria $\mathcal C_S$; assim,

\src{311}
$C^*(T/S,F)$ é um complexo de cocadeias em $\mathcal C'$ que, em dimensão
$n$, é igual a
\[
  F(T\times_ST\times_S\cdots\times_ST)
\]
(onde o parêntese contém $n+1$ fatores).

Observe-se que, como de costume, pode-se definir $H^0(T/S,F)$ sem supor
abeliana a categoria $\mathcal C'$: é o núcleo (definição 2.1), se existe,
do par de morfismos $F(p_i)$ ($i=1,2$)
\[
  F(T)\rightrightarrows F(T\times_ST)
\]
correspondentes às duas projeções
$p_1,p_2:T\times_ST\rightrightarrows T$. Em particular, ter-se-á um morfismo
natural (chamado de \emph{aumento})
\[
  F(S)\longrightarrow H^0(T/S,F)
\]
que será um isomorfismo nos casos favoráveis (em particular, se
$T\longrightarrow S$ é um epimorfismo estrito e $F$ é «exato pela
esquerda»). Do mesmo modo, quando $F$ toma valores na categoria dos
grupos numa categoria $\mathcal C''$, também se pode definir $H^1(T/S,F)$;
quando $\mathcal C''$ é a categoria dos conjuntos (isto é, $F$ toma
valores na categoria dos grupos ordinários, não necessariamente
comutativos), $H^1(T/S,F)$ é o quociente do subgrupo $Z^1(T/S,F)$ de
$C^1(T/S,F)=F(T\times_ST)$ formado pelos $g$ tais que
\[
  F(p_{31})(g)=F(p_{32})(g)\,F(p_{21})(g)
\]
pelo grupo de operadores $F(T)$, que atua sobre $C^1(T/S,F)$ e, em
particular, sobre o subconjunto $Z^1(T/S,F)$ mediante
\[
  \rho(g')\cdot g=F(p_2)(g')\,g\,F(p_1)(g')^{-1}.
\]

\medskip\noindent
b. Seja, por exemplo, $\mathcal F$ uma categoria fibrada de base $\mathcal C$.
Sejam $\xi,\eta\in\mathcal F_S$ e, para todo $S'$ sobre $S$, seja
\[
  F_{\xi,\eta}(S')=
  \operatorname{Hom}(\xi\times_SS',\eta\times_SS').
\]
Assim, $F_{\xi,\eta}$ é um funtor contravariante de $\mathcal C_S$ para a
categoria dos conjuntos. Dito isto, \emph{afirmar que o morfismo de
aumento}
\[
  F_{\xi,\eta}(S)\longrightarrow H^0(S'/S,F_{\xi,\eta})
\]
\emph{é um isomorfismo para todo par de elementos
$\xi,\eta\in\mathcal F_S$ significa que
$\alpha:S'\longrightarrow S$ é um morfismo de $\mathcal F$-descida}
(definição 1.7).

\medskip\noindent
c. Ponhamos igualmente, para $\xi\in\mathcal F_S$ e todo objeto $S'$ de
$\mathcal C$ sobre $S$,
\[
  G_\xi(S')=\operatorname{Aut}(\xi\times_SS') ;
\]
definiu-se assim um funtor contravariante $G_\xi$ de $\mathcal C_S$ para a
categoria dos

\src{312}
grupos. Dito isto, constata-se que $Z^1(S'/S,G_\xi)$ \emph{se identifica
canonicamente com o conjunto dos dados de descida sobre
$\xi'=\xi\times_SS'$ relativos a $S'\longrightarrow S$} (definição 1.6), e
$H^1(S'/S,G_\xi)$ \emph{se identifica com o conjunto das classes (salvo
isomorfismo) de objetos de $\mathcal F_{S'}$ providos de um dado de descida
relativo a $\alpha:S'\longrightarrow S$ que, como objetos de
$\mathcal F_{S'}$, são isomorfos a $\xi'=\xi\times_SS'$}. Portanto, se
$\alpha:S'\longrightarrow S$ é um morfismo de $\mathcal F$-descida
(cf. (b)), então $H^1(S'/S,G_\xi)$ \emph{contém como subconjunto o
conjunto das classes (salvo isomorfismo) de objetos $\eta$ de
$\mathcal F_S$ tais que $\eta\times_SS'$ seja isomorfo em
$\mathcal F_{S'}$ a $\xi\times_SS'$; e essa inclusão é uma identidade se e
somente se todo dado de descida sobre $\xi'=\xi\times_SS'$ relativo a
$\alpha:S'\longrightarrow S$ é efetivo}. (Esse será o caso, em particular,
se $\alpha:S'\longrightarrow S$ é um morfismo de $\mathcal F$-descida
estrito.)

\result{OBSERVAÇÃO}{} Os complexos de cocadeias do tipo $C^*(T/S,F)$
contêm como casos particulares a maioria dos complexos padrão
conhecidos (cohomologia de Čech, cohomologia de grupos, etc.) e desempenham um
papel importante em geometria algébrica, em particular na «cohomologia de
Weil» dos pré-esquemas.

\medskip\noindent
d. \result{EXEMPLO}{1} Seja $S'$ um objeto sobre $S\in\mathcal C$, e seja
$\Gamma$ um grupo de automorfismos de $S'$ tal que $S'$ seja «formalmente
principal sobre $S$, de grupo $\Gamma$», isto é, tal que o morfismo natural
\[
  \Gamma\times S'\longrightarrow S'\times_SS',
\]
onde $\Gamma\times S'$ denota a soma direta de $\Gamma$ cópias de $S'$, seja
um isomorfismo. (Supõe-se que em $\mathcal C$ existem as somas diretas que
intervêm aqui.) Seja $F$ um funtor contravariante de $\mathcal C$ para a
categoria dos grupos abelianos. Então $C^*(S'/S,F)$ é canonicamente
isomorfo ao grupo simplicial padrão de cocadeias homogêneas
$C^*(\Gamma,F(S'))$ e, portanto, $H^*(S'/S,F)$ é canonicamente isomorfo a
$H^*(\Gamma,F(S'))$.

\medskip\noindent
e. \result{EXEMPLO}{2} Seja $\mathcal C$ a categoria dos pré-esquemas. Denota-se
por $G_a$ («grupo aditivo») o funtor contravariante de $\mathcal C$ para a
categoria dos grupos abelianos definido por
\[
  G_a(X)=H^0(X,\mathcal O_X).
\]
Define-se de igual modo o funtor $G_m$ («grupo multiplicativo») por meio de
\[
  G_m(X)=H^0(X,\mathcal O_X)^*
\]
(grupo dos elementos invertíveis do anel $H^0(X,\mathcal O_X)$), e, mais
geralmente, o funtor $\operatorname{Gl}(n)$ («grupo linear de ordem $n$»)
mediante

\src{313}
\[
  \operatorname{Gl}(n)(X)=
  \operatorname{Gl}\bigl(n,H^0(X,\mathcal O_X)\bigr),
\]
que é um funtor de $\mathcal C$ à categoria dos grupos (não
necessariamente comutativos se $n>1$; para $n=1$ se recupera $G_m$). Além disso,
pode-se interpretar $\operatorname{Gl}(n)$ como um funtor de automorfismos
(cf. (c)) considerando a categoria fibrada $\mathcal F$ de base $\mathcal C$
tal que, para $X\in\mathcal C$, $\mathcal F_X$ seja a categoria dos feixes
localmente livres sobre $X$: com efeito, tem-se
\[
  \operatorname{Gl}(n)(X)=
  \operatorname{Aut}_{\mathcal F_X}(\mathcal O_X^n).
\]
De (b) segue-se que, se $\alpha:S'\longrightarrow S$ é um morfismo de
$\mathcal F$-descida (cf. parágrafo 2.c),
$H^1(S'/S,\operatorname{Gl}(n))$ contém o conjunto das classes (salvo
isomorfismo) de feixes localmente livres sobre $S$ cuja imagem inversa sobre
$S'$ é isomorfa a $\mathcal O_{S'}^n$, e essa inclusão é uma igualdade se e
somente se todo dado de descida sobre $\mathcal O_{S'}^n$ (relativo a
$\alpha:S'\longrightarrow S$) é efetivo. Quando $S$ é o espectro de um
anel local, isto significa, portanto,
$H^1(S'/S,\operatorname{Gl}(n))=(e)$, pois todo feixe localmente livre sobre $S$
é então trivial.

Assinalemos a equivalência das condições seguintes para um morfismo
$\alpha:S'\longrightarrow S$:
\begin{enumerate}
\item[(i)] o homomorfismo de aumento
\[
  H^0(S,\mathcal O_S)=G_a(S)\longrightarrow H^0(S'/S,G_a)
\]
é um isomorfismo;
\item[(ii)] $S'\longrightarrow S$ é um morfismo de $\mathcal F$-descida
(sendo $\mathcal F$ a categoria fibrada de base $\mathcal C$ considerada
acima).
\end{enumerate}
Se $S'\longrightarrow S$ é finito, essas condições equivalem também a:
\begin{enumerate}
\item[(iii)] $S'\longrightarrow S$ é um epimorfismo estrito
(cf. parágrafo 2.c).
\end{enumerate}

Suponhamos agora que $S=\operatorname{Spec}(A)$,
$S'=\operatorname{Spec}(A')$; então
\[
  C^n(S'/S,G_a)=C^n(A'/A,G_a)=\bigotimes_A^{n+1}A',
\]
sendo o operador coborde
$C^n(A'/A,G_a)\longrightarrow C^{n+1}(A'/A,G_a)$ a soma alternada dos
operadores de face
\[
  \partial_i(x_0\otimes x_1\otimes\cdots\otimes x_n)
  =x_0\otimes\cdots\otimes x_{i-1}\otimes1_{A'}\otimes
   x_i\otimes\cdots\otimes x_n.
\]
Do mesmo modo,
\[
  C^n(S'/S,G_m)=C^n(A'/A,G_m)
  \simeq\left(\bigotimes_A^{n+1}A'\right)^*,
\]
e as operações simpliciais em $C^*(A'/A,G_m)$ são induzidas pelas de
$C^*(S'/S,G_a)$. Explicitam-se analogamente as operações simpliciais em
$C^*(A'/A,\operatorname{Gl}(n))$. \emph{Em todos os casos que o
conferenciante conhece, tem-se}
\[
  H^i(A'/A,G_a)=0\qquad\text{para $i>0$},
\]
\emph{e, se $A$ é local, tem-se}
\[
  H^1(A'/A,G_m)=0
  \quad\text{e, mais geralmente,}\quad
  H^1(A'/A,\operatorname{Gl}(n))=(e)
\]
(quando $S'\longrightarrow S$ é um morfismo de $\mathcal F$-descida,
isto é, quando o diagrama
$A\longrightarrow A'\rightrightarrows A'\otimes_AA'$ é exato; compare-se com
o parágrafo 2.c). Observe-se que \emph{o «teorema 90» de Hilbert nada mais é
que a relação}
\src{314}
\[
  H^1(S'/S,G_m)=0
\]
\emph{quando $A$ é um corpo e $A'$ uma extensão de Galois finita deste
(cf. exemplo 1), e também se pode exprimir dizendo que, no caso
considerado, $S'\longrightarrow S$ é um morfismo de descida estrito para a
categoria fibrada dos feixes localmente livres de posto 1.} É sob esta
última forma que convém generalizar o teorema de Hilbert, variando tanto
as hipóteses sobre o morfismo $S'\longrightarrow S$ como as relativas aos
feixes quase-coerentes considerados.

Assinalemos, por fim, a equivalência das propriedades seguintes quando
$A$ é um anel local artiniano de ideal maximal $\mathfrak m$ e $A'$ uma
$A$-álgebra (denotando, para todo inteiro $k>0$, por $A_k$ (resp. $A'_k$) os
anéis $A/\mathfrak m^{k+1}$ (resp. $A'/\mathfrak m^{k+1}A'$)):
\begin{enumerate}
\item[(i)] $H^1(A'_k/A_k,G_a)=0$ para todo $k$.
\item[(ii)] $H^1(A'_k/A_k,G_m)=0$ para todo $k$.
\item[(iii)] $H^1(A'_k/A_k,\operatorname{Gl}(n))=(e)$ para todo $k$ e todo
$n$.
\end{enumerate}
Se $S'\longrightarrow S$ é um epimorfismo estrito, as condições
anteriores implicam inclusive que é um morfismo de descida estrito para os
módulos livres (de tipo finito ou não) sobre $A'$.

\result{OBSERVAÇÃO}{} A definição dos grupos $H^i(S'/S,G_m)$, quando
$S,S'$ são esquemas de corpos $A,A'$, se deve a Amitsur. O grupo
$H^2(S'/S,G_m)$ é particularmente interessante como variante «global» do
grupo de Brauer; para essa variante pode-se consultar [1], capítulo VII.

\fsection{B}{Descida por morfismos fielmente planos}

\subsection*{1. Enunciado dos teoremas de descida}
\addcontentsline{toc}{subsection}{1. Enunciado dos teoremas de descida}

\result{DEFINIÇÃO}{1.1} Um morfismo
$\alpha:S'\longrightarrow S$ de pré-esquemas diz-se \emph{plano} se, para todo
$x'\in S'$, $\mathcal O_{x'}$ é um módulo plano sobre o anel
$\mathcal O_{\alpha(x')}$ (isto é, se
$\mathcal O_{x'}\otimes_{\mathcal O_{\alpha(x')}}M$ é um funtor exato no
$\mathcal O_{\alpha(x')}$-módulo $M$). Um morfismo se diz
\emph{fielmente plano} se é plano e sobrejetivo.

Por exemplo, se $S=\operatorname{Spec}(A)$ e
$S'=\operatorname{Spec}(A')$, então $S'$ é plano sobre $S$ se e somente se
$A'$ é um $A$-módulo plano, e $S'$ é fielmente plano sobre $S$ se e somente se
$A'$ é um $A$-módulo fielmente plano (isto é, o funtor
$A'\otimes_A M$ no $A$-módulo $M$ é exato e fiel); isto significa também,
na terminologia de Serre [5], que o par $(A,A')$ é plano. Se $S'$ é
fielmente plano sobre $S$, então o funtor imagem inversa de feixes
quase-coerentes sobre $S$ é exato e fiel; em outras palavras, para que uma
sucessão de homomorfismos de feixes

\src{315}
quase-coerentes sobre $S$ seja exata, é necessário e suficiente que também o seja a sua
imagem inversa sobre $S'$ (em particular, para que um homomorfismo de feixes
quase-coerentes sobre $S$ seja um monomorfismo, resp. um epimorfismo, resp. um
isomorfismo, é necessário e suficiente que também o seja a sua imagem inversa sobre
$S'$). Essa propriedade continua válida se nos restringirmos a um aberto
qualquer de $S$ e, sob essa forma, caracteriza os morfismos fielmente
planos.

\result{DEFINIÇÃO}{1.2} Um morfismo $\alpha:S'\longrightarrow S$ diz-se
\emph{quase-compacto} se a imagem inversa de toda parte aberta quase-compacta
$U$ de $S$ é quase-compacta (isto é, união finita de abertos afins).

Evidentemente, basta verificar esta propriedade para os abertos afins de $S$.
Por exemplo, um morfismo afim (isto é, tal que a imagem inversa de um aberto
afim seja afim) é quase-compacto.

A classe dos morfismos planos, resp. fielmente planos, resp. quase-compactos,
é estável por composição e por «mudança de base», e contém, naturalmente,
os isomorfismos.

\result{TEOREMA}{1} Seja $\alpha:S'\longrightarrow S$ um morfismo de
pré-esquemas, fielmente plano e quase-compacto. Então $\alpha$ é um
\emph{morfismo de descida estrito} (A, definição 1.7) para a categoria
fibrada $\mathcal F$ dos feixes quase-coerentes (A, parágrafo 1, exemplo 2).

Esse enunciado significa duas coisas:
\begin{enumerate}
\item[(i)] Se $F$ e $G$ são dois feixes quase-coerentes sobre $S$, e $F'$ e
$G'$ suas imagens inversas sobre $S'$, então o homomorfismo natural
\[
  \operatorname{Hom}(F,G)\longrightarrow\operatorname{Hom}(F',G')
\]
é uma bijeção do primeiro membro sobre o subgrupo do segundo formado pelos
homomorfismos $F'\longrightarrow G'$ compatíveis com os dados de descida
canônicos sobre esses feixes, isto é, aqueles cujas imagens inversas pelas
duas projeções de $S''=S'\times_SS'$ sobre $S'$ dão um mesmo
homomorfismo $F''\longrightarrow G''$.
\item[(ii)] Todo feixe quase-coerente $F'$ sobre $S'$, provido de um dado de
descida relativo ao morfismo $\alpha:S'\longrightarrow S$
(A, definição 1.6), é isomorfo (provido deste dado) à imagem inversa de
um feixe quase-coerente $F$ sobre $S$.
\end{enumerate}
Tomando $F=\mathcal O_S$ em (i), obtém-se:

\result{COROLÁRIO}{1} Seja $G$ um feixe quase-coerente sobre $S$; sejam
$G'$ e $G''$ suas imagens inversas sobre $S'$ e sobre
$S''=S'\times_SS'$, e sejam $p_1,p_2$ as duas projeções de $S''$ sobre
$S'$. Então o seguinte diagrama de aplicações de conjuntos

\src{316}
\[
  \Gamma(G)\xrightarrow{\alpha^*}\Gamma(G')
  \underset{p_2^*}{\overset{p_1^*}{\rightrightarrows}}\Gamma(G'')
\]
é exato (A, definição 1.(a)).

Por outro lado, a conjunção de (i), (ii) e da definição 1.1 dá o

\result{COROLÁRIO}{2} Seja $G$ como no corolário 1. Então existe uma
correspondência biunívoca entre os subfeixes quase-coerentes de $G$ e os
subfeixes quase-coerentes de $G'$ cujas imagens inversas sobre $S''$, pelas
duas projeções $p_1,p_2$, dão o mesmo subfeixe de $G''$.

Naturalmente, tem-se um enunciado equivalente em termos de feixes quociente.
Como se viu (A, parágrafo 4, (e)), o teorema 1 deve ser considerado uma
generalização do «teorema 90» de Hilbert, e implica como casos particulares
diversas formulações em termos de 1-cohomologia. Para a demonstração, reduz-se
facilmente ao caso em que $S=\operatorname{Spec}(A)$,
$S'=\operatorname{Spec}(A')$, e para (i) reduz-se facilmente a demonstrar o
corolário 1, isto é, a exatidão do diagrama
\[
  M=A\otimes_A M\longrightarrow A'\otimes_A M
  \rightrightarrows A'\otimes_AA'\otimes_A M
\]
para todo $A$-módulo $M$, o que se deduz do lema mais geral:

\result{LEMA}{1.1} Seja $A'$ uma $A$-álgebra fielmente plana. Então, para
todo $A$-módulo $M$, o complexo aumentado por $M$
$C^*(A'/A,G_a)\otimes_A M$ (cf. A, parágrafo 4, (e)) é uma
\emph{resolução} de $M$.

Basta demonstrar que o complexo aumentado deduzido do anterior mediante a
mudança de base de $A$ para $A'$ satisfaz as mesmas conclusões. Isso conduz a
verificar o enunciado quando se substituem $A$ por $A'$ e $A'$ por
$A'\otimes_AA'$, o que nos reduz a um caso no qual existe um
homomorfismo de $A$-álgebras $A'\longrightarrow A$ (ou, em termos
geométricos, ao caso em que $S'$ sobre $S$ admite uma seção). Nesse caso,
isso se deduz das generalidades de A, parágrafo 4, (a). Assinalemos de passagem
o seguinte corolário, que generaliza um enunciado bem conhecido em
cohomologia de Galois (compare-se com A, parágrafo 4, (e)):

\result{COROLÁRIO}{} Se $A'$ é fielmente plana sobre $A$, tem-se
$H^0(A'/A,G_a)=A$ e
\[
  H^i(A'/A,G_a)=0\qquad\text{para $i\geq1$}.
\]

Para demonstrar a parte (ii) do teorema 1, procede-se como para (i),
reduzindo-se ao caso em que $S'$ sobre $S$ admite uma seção, onde se deduz
de (i) (cf. A, parágrafo 1, (c)).

Evidentemente, pode-se variar ad libitum o teorema 1 e seus corolários,
introduzindo diversas estruturas adicionais nos feixes (ou sistemas de
feixes) quase-coerentes considerados. Por exemplo, dar sobre $S$ um feixe
\src{317}
quase-coerente de álgebras comutativas «equivale» a dar sobre $S'$ um feixe desse
tipo provido de um dado de descida relativo a
$\alpha:S'\longrightarrow S$. Tendo em conta a correspondência funtorial
entre tais feixes quase-coerentes sobre $S$ e os pré-esquemas afins sobre
$S$, obtém-se a segunda afirmação do teorema seguinte:

\result{TEOREMA}{2} Seja $\alpha:S'\longrightarrow S$ como no teorema 1.
Então $\alpha$ é um \emph{morfismo de descida} (não estrito em geral)
(A, definição 2.4), e é um \emph{morfismo de descida estrito} para a
categoria fibrada dos esquemas afins sobre pré-esquemas
(A, definição 1.7).

A primeira afirmação do teorema significa o seguinte: sejam $X,Y$ dois
pré-esquemas sobre $S$; sejam $X',Y'$ suas imagens inversas sobre $S'$, e
$X'',Y''$ suas imagens inversas sobre $S''=S'\times_SS'$. Então o seguinte
diagrama de aplicações naturais
\[
  \operatorname{Hom}_S(X,Y)\xrightarrow{\alpha^*}
  \operatorname{Hom}_{S'}(X',Y')
  \underset{p_2^*}{\overset{p_1^*}{\rightrightarrows}}
  \operatorname{Hom}_{S''}(X'',Y'')
\]
é exato; isto é, $\alpha^*$ é uma bijeção de
$\operatorname{Hom}_S(X,Y)$ sobre a parte de
$\operatorname{Hom}_{S'}(X',Y')$ formada pelos homomorfismos compatíveis com
os dados de descida canônicos sobre $X',Y'$ (isto é, aqueles cujas
imagens inversas pelas duas projeções de $S''$ sobre $S'$ são iguais).
Isso se deduz facilmente do teorema 1, corolário 1, quando se restringe
a $Y$ afim sobre $S$; no caso geral, é preciso combinar o teorema 1 com o
resultado seguinte:

\result{LEMA}{1.2} Seja $\alpha:S'\longrightarrow S$ um morfismo fielmente plano
e quase-compacto. Então $S$ identifica-se com um \emph{espaço topológico
cociente} de $S'$; isto é, toda parte $U$ de $S$ tal que
$\alpha^{-1}(U)$ seja aberta é aberta.

Para completar o teorema 2, é preciso dar critérios de efetividade para um dado
de descida sobre um $S'$-pré-esquema $X'$ (quando não se supõe que $X'$ seja
afim sobre $S'$). Assinalemos primeiro que \emph{tal dado de descida não é
necessariamente efetivo}, mesmo se $S$ é o espectro de um corpo $k$, $S'$
o espectro de uma extensão quadrática $k'$ deste e $X'$ um esquema
algébrico próprio de dimensão 2 sobre $S'$ (como se pode ver, segundo Serre,
utilizando a superfície não projetiva de Nagata). \emph{Para que um dado de
descida sobre $X'/S'$ relativo a $\alpha:S'\longrightarrow S$ (fielmente
plano e quase-compacto) seja efetivo, é necessário e suficiente que $X'$ seja
união de abertos $X'_i$, afins sobre $S'$, que sejam «estáveis» pelo dado de
descida sobre $X'$.} Isso certamente se cumpre (quaisquer que sejam $X'/S'$
e o dado de descida sobre $X'$) se o morfismo
$\alpha:S'\longrightarrow S$ é \emph{radicial} (isto é, injetivo e com
extensões residuais que

\src{318}
são radiciais). Pode-se demonstrar também que continua a cumprir-se se
$\alpha:S'\longrightarrow S$ é \emph{finito} e toda parte finita de $X'$,
contida numa fibra de $X'$ sobre $S$, está contida num aberto de $X'$
afim sobre $S$ (esse é o \emph{critério de Weil}). Em particular, isso se
cumpre se $X'/S'$ é quase-projetivo, e nesse caso pode-se demonstrar que o
pré-esquema «descido» $X/S$ também é quase-projetivo (e projetivo se
$X'/S'$ também o é). Em resumo:

\result{TEOREMA}{3} Seja $\alpha:S'\longrightarrow S$ um morfismo de
pré-esquemas, fielmente plano e quase-compacto. Se $\alpha$ é \emph{radicial},
é um \emph{morfismo de descida estrito}. Se $\alpha$ é finito, é um
morfismo de descida estrito com respeito à categoria fibrada dos
pré-esquemas quase-projetivos (ou projetivos) sobre pré-esquemas.

\result{OBSERVAÇÕES}{} Ignoro se, na segunda afirmação anterior, a
hipótese de que $\alpha$ seja um morfismo \emph{finito} é realmente necessária;
em todo caso, verifica-se de maneira «formal» que pode ser substituída pela
seguinte hipótese, aparentemente mais fraca: \emph{para todo ponto de $S$
existe uma vizinhança aberta $U$, um $U'$ finito e fielmente plano sobre $U$, e um
$S$-morfismo de $U'$ em $S'$.} Um caso típico que não entra no anterior é
aquele em que $S=\operatorname{Spec}(A)$,
$S'=\operatorname{Spec}(\widehat A)$, onde $A$ é um anel local noetheriano
e $\widehat A$ seu completamento; ou também aquele em que $S'$ é quase-finito sobre
$S$ (isto é, localmente isomorfo a um aberto de um $S$-esquema finito) e não
finito. Nesses dois casos, o conferenciante também ignora a resposta à
pergunta seguinte: seja $X$ um $S$-esquema tal que $X'=X\times_SS'$ seja
projetivo sobre $S'$; é verdade que $X$ é projetivo sobre
$S$?\footnote{A errata de 1962 responde negativamente: um morfismo
$S'\longrightarrow S$ quase-finito, étale e sobrejetivo, ou um morfismo
$\operatorname{Spec}(\widehat A)\longrightarrow\operatorname{Spec}(A)$, não é
sempre um morfismo de descida estrito, mesmo se $A$ é o anel local de
uma curva algébrica sobre um corpo algebricamente fechado $k$ e
$S=\operatorname{Spec}(A)$. Pode-se encontrar assim um morfismo próprio e liso
$f:X\longrightarrow S$, que faz de $X$ um fibrado principal de base $S$ sob
uma curva elíptica $E$, tal que $f':X'\longrightarrow S'$ seja projetivo, mas
$f$ não o seja. Portanto, também é um exemplo de fibrado principal homogêneo
não isotrivial sob um esquema abeliano.}

\subsection*{2. Aplicação à descida de certas propriedades dos morfismos}
\addcontentsline{toc}{subsection}{2. Descida de propriedades dos morfismos}

Seja $P$ uma classe de morfismos de pré-esquemas. Seja
$\alpha:S'\longrightarrow S$ um morfismo de pré-esquemas, seja
$f:X\longrightarrow Y$ um morfismo de $S$-pré-esquemas e seja
$f':X'\longrightarrow Y'$ a imagem inversa de $f$ por $\alpha$. Cabe
perguntar-se então se a relação «$f'\in P$» implica «$f\in P$». A resposta é
afirmativa em muitos casos importantes quando se supõe que $\alpha$ é
\emph{fielmente plano e quase-compacto} (sendo às vezes supérflua essa última
hipótese). Isso se vê diretamente e sem dificuldade se $P$ é a classe dos
morfismos sobrejetivos, resp. radiciais (esses casos se deduzem da
sobrejetividade de $\alpha$), resp. planos, resp. fielmente planos, resp. lisos
(esses casos se deduzem da planicidade fiel de $\alpha$), resp. de tipo finito.
Utilizando os teoremas 1 e 2 e o lema 1.2, vê-se também que ocorre o mesmo se
$P$ é uma das classes seguintes: isomorfismos, imersões abertas, imersões
fechadas, imersões (se $f$ é de tipo finito e $Y$
localmente noetheriano), morfismos afins, morfismos finitos, morfismos
quase-finitos, morfismos abertos, morfismos fechados, homeomorfismos, morfismos
separados,

\src{319}
morfismos próprios. O único caso importante não elucidado é o dos
morfismos projetivos e quase-projetivos, já assinalado na observação do
parágrafo 1.

\subsection*{3. Descida por morfismos finitos fielmente planos}
\addcontentsline{toc}{subsection}{3. Descida por morfismos finitos fielmente planos}

Seja $\alpha:S'\longrightarrow S$ um morfismo finito correspondente a um feixe de
álgebras $\mathcal A'$ sobre $S$ que, como feixe de módulos, é localmente livre
de tipo finito e é $\ne0$ em toda parte; então $\alpha$ é um morfismo
fielmente plano e quase-compacto, ao qual se podem aplicar os resultados
anteriores. Dar um feixe quase-coerente $F'$ sobre $S'$ equivale a dar o feixe
quase-coerente $\alpha_*(F')$ sobre $S$, provido de sua estrutura de
$\mathcal A'$-módulo (observando que
$\mathcal A'=\alpha_*(\mathcal O_{S'})$). Para simplificar, este feixe sobre $S$
também será denotado por $F'$. As duas imagens inversas $p_i^*(F')$ de $F'$
sobre $S'\times_SS'$ correspondem igualmente aos feixes quase-coerentes de
$(\mathcal A'\otimes_{\mathcal O_S}\mathcal A')$-módulos
\[
  F'\otimes_{\mathcal O_S}\mathcal A'
  \quad\text{e}\quad
  \mathcal A'\otimes_{\mathcal O_S}F'.
\]
Dar um $(S'\times_SS')$-homomorfismo do primeiro ao segundo equivale a dar um
homomorfismo de $(\mathcal A'\otimes_{\mathcal O_S}\mathcal A')$-módulos e,
tendo em conta que $\mathcal A'$ é localmente livre, isto equivale a dar
um homomorfismo de $(\mathcal A'\otimes\mathcal A')$-módulos:
\[
  \mathcal U=
  \operatorname{Hom}_{\mathcal O_S}(\mathcal A',\mathcal A')
  =\mathcal A'\otimes\check{\mathcal A}'
  \longrightarrow
  \operatorname{Hom}_{\mathcal O_S}(F',F'),
\]
isto é, a dar, para cada seção $\xi$ de $\mathcal U$ sobre um aberto $V$,
um homomorfismo de $\mathcal O_S$-módulos
$T_\xi:F'|_V\longrightarrow F'|_V$ que satisfaz as condições
\begin{equation}
  T_{f\xi}(x)=fT_\xi(x),
  \qquad T_{\xi f}(x)=T_\xi(fx),
  \tag{3.1}
\end{equation}
onde $f,x$ são, respectivamente, seções de $\mathcal A',F'$ sobre um
aberto de $S$ contido em $V$. As condições (i) e (ii) de um dado de
descida (A, parágrafo 1, (c)) escrevem-se então, respectivamente:
\begin{equation}
  T_{1_{\mathcal U}}x=x,
  \qquad\text{isto é,}\qquad
  T_{1_{\mathcal U}}=\operatorname{id}_{F'},
  \tag{3.2}
\end{equation}
\begin{equation}
  T_{\xi\eta}=T_\xi T_\eta.
  \tag{3.3}
\end{equation}
Em outras palavras, \emph{um dado de descida sobre $F'$ equivale a uma
representação do feixe de $\mathcal O_S$-álgebras
$\mathcal U=\operatorname{Hom}_{\mathcal O_S}(\mathcal A',\mathcal A')$ no
feixe de $\mathcal O_S$-álgebras
$\operatorname{Hom}_{\mathcal O_S}(F',F')$ que satisfaz as duas condições de
linearidade (3.1).} Se se tem um emparelhamento de feixes quase-coerentes sobre
$S'$:
\[
  F'_1\times F'_2\longrightarrow F'_3
\]
(que pode ser interpretado como um emparelhamento de feixes de
$\mathcal A'$-módulos sobre $S$),

\src{320}
e dados de colagem sobre os $F'_i$, definidos por homomorfismos
\[
  T_i:\mathcal U\longrightarrow
  \operatorname{Hom}_{\mathcal O_S}(F'_i,F'_i)
  \qquad (i=1,2,3),
\]
então esses dados são compatíveis com o emparelhamento dado, no sentido
evidente do termo, se e somente se vigora a condição seguinte.

Para toda seção $\xi$ de $\mathcal U$ sobre um aberto, denotando por
\[
  \Delta\xi=\sum_i\xi'_i\otimes_{\mathcal A'}\xi''_i
\]
a seção de $\mathcal U\otimes_{\mathcal A'}\mathcal U$ (considerando
$\mathcal U$ como $\mathcal A'$-módulo mediante sua estrutura esquerda)
definida pela fórmula
\[
  \xi\mathbin{\cdot}(fg)=\sum_i\xi'_i(f)\xi''_i(g)
\]
(onde $f$ e $g$ são duas seções de $\mathcal A'$ sobre um aberto menor),
tem-se a fórmula
\begin{equation}
  T^{(3)}_\xi(x\mathbin{\cdot}y)
  =\sum_iT^{(1)}_{\xi'_i}x\mathbin{\cdot}T^{(2)}_{\xi''_i}y
  \tag{3.4}
\end{equation}
para todo par de seções $x,y$ de $F'_1,F'_2$, respectivamente, sobre um
aberto menor.\footnote{O texto impresso diz «de $\mathcal A'$». Os domínios de
$T^{(1)}$ e $T^{(2)}$, assim como o emparelhamento
$F'_1\times F'_2\longrightarrow F'_3$, obrigam a que $x$ e $y$ sejam seções
de $F'_1$ e $F'_2$, respectivamente.} (Essa propriedade pode ser expressa
dizendo que os homomorfismos $T^{(i)}$ são compatíveis com a aplicação
diagonal de $\mathcal U$ com respeito ao emparelhamento dado.) Em particular, as
fórmulas (3.1) a (3.4) permitem interpretar, em termos de representações de
álgebras com aplicações diagonais, os dados de descida sobre um feixe
quase-coerente de álgebras sobre $S'$ e, portanto, também (restringindo-se às
álgebras comutativas) os dados de descida sobre um $S'$-esquema afim.

Daí se passa a uma interpretação análoga dos dados de descida sobre um
$S'$-pré-esquema $X'$ qualquer: dar tal $X'$ equivale a dar um pré-esquema
$X'$ sobre $S$, provido de um homomorfismo de $\mathcal O_S$-álgebras
\[
  \mathcal A'\longrightarrow\mathcal O_{X'},
\]
e um dado de descida sobre $X'$ equivale a dar um homomorfismo de feixes
(compatível com o morfismo $h:X'\longrightarrow S'$)
\[
  \mathcal U\longrightarrow
  \operatorname{Hom}_{h^{-1}(\mathcal O_S)}
  (\mathcal O_{X'},\mathcal O_{X'}),
\]
que satisfaz condições análogas às condições (3.1) a (3.4) anteriores.

\textsc{Exemplo 1 («Weil»).} -- Suponhamos que $S'/S$ é um recobrimento étale
de Galois de grupo de Galois $\Gamma$ (cf. A, parágrafos 3 e 4.d). Então, um
dado de descida sobre um feixe quase-coerente $F'$ sobre $S'$ (resp. sobre um
$S'$-pré-esquema $X'$) equivale a dar uma representação de $\Gamma$ por
automorfismos de $(S',F')$ (resp. de $(S',X')$) compatível com as operações
de $\Gamma$ sobre $S'$. Este resultado

\src{321}
é «formal», isto é, demonstra-se em termos de categorias, mas, do
ponto de vista deste número, deduz-se também da estrutura explícita de
$\mathcal U$ (provido de sua estrutura de anel, do homomorfismo de anéis
$\mathcal A'\longrightarrow\mathcal U$ e da aplicação diagonal), conhecida
por completo graças ao resultado seguinte: $\mathcal U$ admite, como
$\mathcal A'$-módulo esquerdo, uma base formada pelas seções de
$\mathcal U$ que correspondem aos elementos de $\Gamma$.

\textsc{Exemplo 2 («Cartier»).} -- Seja $p$ um número primo; suponhamos que
$p\mathcal O_S=0$ (isto é, $\mathcal O_S$ é de característica $p$),
$(\mathcal A')^p\subset\mathcal O_S=\mathcal A$ (isto é, $S'/S$ é radicial
de altura 1) e que o feixe de álgebras $\mathcal A'$ sobre $\mathcal A$ admite
localmente uma $p$-base, isto é, uma família $(x_i)$ de seções tal que
$\mathcal A'$ está gerada como álgebra pelos $x_i$, submetidos às únicas
condições $x_i^p=0$. Supomos que o conjunto dos $i$ é finito, de
cardinal $n$. Seja $\mathfrak g$ o feixe das $\mathcal A$-derivações de
$\mathcal A'$; é um feixe de $\mathcal A'$-módulos localmente livre de posto
$n$ e, além disso, um feixe de $p$-álgebras de Lie sobre $\mathcal A$ (mas não sobre
$\mathcal A'$) que satisfaz a condição
\begin{equation}
  [X,fY]=X(f)Y+f[X,Y].
  \tag{3.5}
\end{equation}

\textsc{Lema.} -- $\mathcal U=
\operatorname{Hom}_{\mathcal O_S}(\mathcal A',\mathcal A')$ está gerada,
como $\mathcal O_S$-álgebra provida de um homomorfismo de álgebras
$\mathcal A'\longrightarrow\mathcal U$, pelo sub-$\mathcal A'$-módulo
esquerdo $\mathfrak g$, com as relações adicionais
\begin{equation}
  \left\{
  \begin{aligned}
    Xf-fX&=X(f),\\
    XY-YX&=[X,Y],\\
    X^p&=X^{(p)}.
  \end{aligned}
  \right.
  \tag{3.6}
\end{equation}

Do lema anterior se deduz que um dado de descida sobre o feixe
quase-coerente $F'$ sobre $S'$ equivale a dar, para todo $X\in\mathfrak g$, um
$\mathcal O_S$-endomorfismo $\overline X$ de $F'$ que satisfaz as condições
\begin{align}
  \overline{fX}&=f\overline X, \tag{3.7}\\
  \overline X(fx)&=X(f)x+f\overline X(x), \tag{3.8}\\
  \overline{[X,Y]}&=[\overline X,\overline Y], \tag{3.9}\\
  \overline{X^{(p)}}&=\overline X^p. \tag{3.10}
\end{align}
(Isso é o que se poderia chamar uma conexão linear sobre $F'$, sem curvatura
e compatível com a potência $p$-ésima.) Explicita-se de igual modo a noção
de dado de descida sobre um $S'$-pré-esquema $X'$; a relação (3.4) é
substituída aqui pela condição de que os $\overline X$ sejam derivações de
$\mathcal O_{X'}$. Como o morfismo $S'\longrightarrow S$ é radicial, o
teorema 3 garante que todo dado de descida desse tipo é

\src{322}
efetivo e, portanto, define um $S$-pré-esquema $X$.

Observe-se que não foi necessário impor nenhuma hipótese de planicidade, de
não singularidade nem de finitude de tipo algum sobre $F'$, resp. $X'$.

\subsection*{4. Aplicação a critérios de racionalidade}
\addcontentsline{toc}{subsection}{4. Aplicação a critérios de racionalidade}

Seja $X$ um $S$-pré-esquema tal que a imagem direta de $\mathcal O_X$ sobre
$S$ seja $\mathcal O_S$; essa propriedade continuará válida após toda mudança
plana $S'\longrightarrow S$ da base $S$. Se $F$ é um feixe invertível
(isto é, localmente livre de posto 1) sobre $X$, os automorfismos de $F$, que se
identificam com as seções invertíveis de $\mathcal O_X$, correspondem
biunívocamente às seções invertíveis de $\mathcal O_S$. Seja então $s$
uma seção de $X$ sobre $S$; chamaremos, de maneira gráfica, seção de $F$
sobre $s$ a uma seção do feixe invertível $s^*(F)$ sobre $S$. Do anterior
deduz-se que, se $F_i$ ($i=1,2$) são dois feixes invertíveis sobre $X$, cada um
provido de uma seção sobre $s$, e se $F_1$ e $F_2$ são isomorfos, existe um
único isomorfismo de $F_1$ sobre $F_2$ compatível com as seções em questão
(isto é, que transforma a primeira na segunda). Por outro lado, e
independentemente da seção $s$, convenhamos em considerar equivalentes
dois feixes invertíveis $F_1$ e $F_2$ sobre $X$ tais que todo ponto de $S$ tenha
uma vizinhança aberta $U$ para a qual as restrições de $F_1$ e $F_2$ a
$X|U$ sejam isomorfas. Então todo feixe invertível $F$ sobre $X$ é equivalente
a um feixe invertível $F_1$ provido de uma seção distinguida sobre $s$
(toma-se $F_1=F\,s^*(F)^{-1}$), e $F_1$ está determinado salvo isomorfismo. Em
outras palavras, a classificação dos feixes invertíveis sobre $X$ salvo
equivalência é a mesma que a classificação, salvo isomorfismo, dos feixes
invertíveis providos de uma seção distinguida.

Como essas propriedades continuam válidas após uma mudança plana
$\alpha:S'\longrightarrow S$ da base (substituindo a seção $s$ por sua
imagem inversa $s'$ por $\alpha$), se conclui, tendo em conta o teorema 1:

\emph{Seja o pré-esquema $X/S$ como acima e suponhamos que admita uma seção
$s$; seja $\alpha:S'\longrightarrow S$ um morfismo fielmente plano e
quase-compacto; seja $F'$ um feixe invertível sobre $X'=X\times_SS'$. Para que $F'$
seja equivalente à imagem inversa sobre $X'$ de um feixe invertível $F$ sobre
$X$, é necessário e suficiente que suas imagens inversas $p_1^*(F')$ e
$p_2^*(F')$ sobre $X'\times_XX'=X\times_S(S'\times_SS')$ sejam equivalentes. Se
for assim, $F$ está determinado salvo equivalência.} (Dir-se-á então que $F'$
é racional sobre $S$.)

Inspirando-se nesse princípio quando $\alpha:S'\longrightarrow S$ é como no
exemplo 1 ou no exemplo 2 do número anterior, obtêm-se os critérios de
racionalidade de Weil ou de Cartier. (Observe-se que esses autores se limitam ao
caso em que $S$ e $S'$

\src{323}
são espectros de corpos; a fortiori, $S$ é então o espectro de um anel
local, e a relação de equivalência introduzida acima nada mais é que a relação
de isomorfia.) No primeiro caso, $F'$ é racional sobre $S$ se e somente se seus
transformados por $\Gamma$ são equivalentes a $F'$. Para exprimir o critério
de racionalidade no segundo caso, consideram-se, em geral, o morfismo
diagonal
\[
  X'\longrightarrow X''=X'\times_XX'
\]
de $X'/X$, o feixe de ideais $\mathcal I$ correspondente sobre
$X'\times_XX'$ e o feixe $\mathcal I/\mathcal I^2$, que se identifica com sua
imagem inversa $\Omega^1_{X'/X}$ sobre $X'$. Como as restrições dos
$F''_i=p_i^*(F')$ ($i=1,2$) à diagonal são isomorfas (pois são isomorfas a
$F'$), isto é, $F''_1(F''_2)^{-1}=F''$ tem uma restrição à diagonal
que é trivial, segue-se que a restrição de $F''$ a
$(X'',\mathcal O_{X''}/\mathcal I^2)$ é dada, salvo isomorfismo, por um
elemento bem determinado $\xi$ de
\[
  H^1(X'',\mathcal I/\mathcal I^2)
  =H^1(X',\Omega^1_{X'/X}).
\]
Além disso, nesse caso tem-se
\[
  \Omega^1_{X'/X}
  =\Omega^1_{S'/S}\otimes_{\mathcal O_S}\mathcal O_X,
\]
e, portanto, se $\Omega^1_{S'/S}$ é localmente livre sobre $S$ (como no
caso de Cartier), $\xi$ define uma seção de
\[
  R^1f'_*(\mathcal O_{X'})\otimes\Omega^1_{S'/S}
\]
sobre $S'$ (chamada classe de Atiyah--Cartier do feixe invertível $F'$ sobre
$X'/S$), cuja anulação é necessária e suficiente para que sejam equivalentes
as imagens inversas de $F'$ pelas duas projeções de
\[
  (X'',\mathcal O_{X''}/\mathcal I^2)
  =X\times_S(S'',\mathcal O_{S''}/\mathcal J^2)
\]
sobre $X'$ (onde $\mathcal J$ é o feixe de ideais sobre
$S''=S'\times_SS'$ definido pelo morfismo diagonal
$S'\longrightarrow S'\times_SS'$). Essa anulação é, pois, trivialmente
necessária para que sejam equivalentes as imagens inversas de $F'$ sobre o
próprio $X''=X\times_SS''$ e, portanto, também para que $F'$ seja equivalente
à imagem inversa de um feixe invertível $F$ sobre $X$. Além disso, a classe de
Atiyah--Cartier pode ser interpretada também como a obstrução à existência,
localmente sobre $S'$, de uma conexão de $F'$ relativa às derivações de
$X'/X$; tal conexão fica, além disso, determinada quando se conhecem as
derivações de $F'$ correspondentes às extensões naturais a $X'$ das
derivações de $S'/S$. Disso e dos desenvolvimentos do número anterior
conclui-se facilmente que, no caso do exemplo 2 desse número e quando $X/S$
admite uma seção, a anulação da classe de Atiyah--Cartier é também
suficiente para que $F'$ seja racional sobre $S$.

\subsection*{5. Aplicação à restrição da base em um esquema abeliano}
\addcontentsline{toc}{subsection}{5. Restrição da base em um esquema abeliano}

Seja $S$ um pré-esquema. Chama-se esquema abeliano sobre $S$ um esquema $X$
liso e próprio

\src{324}
sobre $S$ cujas fibras nos pontos $x\in S$ são variedades abelianas sobre os
$\kappa(x)$. Suponhamos que $S$ seja noetheriano e regular
(isto é, que seus anéis locais sejam regulares); então, utilizando o
teorema de conexidade de Murre [4] (ao menos no caso de «características
iguais», no qual o teorema citado está atualmente demonstrado), pode-se
demonstrar que toda seção racional de $X$ sobre $S$ está definida em toda
parte (isto é, é uma seção), o que generaliza um teorema clássico de
Weil. Deduz-se, mais geralmente, que, se $X'$ é um esquema liso sobre $S$,
então toda $S$-aplicação racional de $X'$ em $X$ está definida em toda
parte. Disso se obtém o seguinte, que generaliza um resultado de
Chow--Iang:\footnote{O texto impresso diz «Igusa--Lang»; a errata de 1962
prescreve «Chow--Iang».} Seja $S$ noetheriano e regular, e denote $K$ seu
anel de funções racionais (produto direto de corpos); seja $X$ um
esquema abeliano sobre $K$. Se $X$ é isomorfo a um $K$-esquema da forma
$X_0\times_S\operatorname{Spec}(K)$, onde $X_0$ é um esquema abeliano sobre
$S$, então $X_0$ está determinado salvo isomorfismo único.

Utilizando o resultado de unicidade anterior, vê-se que a questão da
restrição da base em $X$ é local sobre $S$ (e, portanto, que basta saber
efetuar a restrição aos $\operatorname{Spec}(\mathcal O_x)$, com
$x\in S$). Vê-se de igual maneira que, se
$\alpha:S'\longrightarrow S$ é um morfismo liso de tipo finito, se $K'$ é o
anel de funções racionais de $S'$, e se $X\otimes_KK'$ é da forma
$X'_0\times_{S'}\operatorname{Spec}(K')$, então $X'_0$ está provido de um
dado de descida canônico relativo a $\alpha$. Tendo em conta o teorema 3,
conclui-se:

\result{PROPOSIÇÃO}{5.1} Seja $S$ um pré-esquema noetheriano, regular e
irredutível, de corpo de funções racionais $K$; seja $K'$ uma extensão
finita de $K$, não ramificada sobre $S$; seja $S'$ a normalização de $S$ em
$K'$ (que é, portanto, um recobrimento étale de $S$); e seja $X$ um esquema
abeliano sobre $K$ tal que $X\otimes_KK'$ seja da forma
$X'_0\times_{S'}\operatorname{Spec}(K')$, onde $X'_0$ é um esquema abeliano
projetivo sobre $S'$. Então $X$ é da forma
$X_0\times_S\operatorname{Spec}(K)$, onde $X_0$ é um esquema abeliano
projetivo sobre $S$.

\result{OBSERVAÇÕES}{} O conferenciante ignora se pode ser substituída a
hipótese de que $S'\longrightarrow S$ seja um recobrimento étale sobrejetivo
(que permite utilizar o teorema 3) pela hipótese de que seja um morfismo
liso, sobrejetivo e de tipo finito (inclusive se se supõe que é um morfismo
étale), ou se a proposição continua válida sem supor que $X'_0$ seja
projetivo sobre $S'$ (condição que talvez se cumpra automaticamente).

\subsection*{6. Aplicação a critérios de trivialidade local e isotrivialidade}
\addcontentsline{toc}{subsection}{6. Critérios de trivialidade local e isotrivialidade}

Sejam $S$ um pré-esquema, $G$ um «pré-esquema em grupos» sobre $S$, e $P$ um
pré-esquema sobre $S$ sobre o qual «atua $G$» (pela direita). Diz-se que $P$
é formalmente principal homogêneo sob $G$ se o morfismo bem conhecido
\[
  G\times_SP\longrightarrow P\times_SP
\]

\src{325}
induzido pela ação de $G$ sobre $P$ é um isomorfismo. Doravante, suporemos que
$G$ seja plano sobre $S$ (e, portanto, fielmente plano sobre $S$), e reservaremos
o nome de fibrado principal homogêneo sob $G$ para um fibrado formalmente
principal homogêneo $P$ que seja fielmente plano e quase-compacto sobre $S$.
É imediato que isso equivale a dizer que se pode encontrar uma mudança
fielmente plana e quase-compacta
$S'\longrightarrow S$ da base $S$, tal que o fibrado formalmente principal
homogêneo
\[
  P'=P\times_SS'
\]
sob $G'=G\times_SS'$ seja trivial, isto é, isomorfo a $G'$ (isto é, que
admita uma seção); em particular, pode-se tomar $S'=P$. Observe-se, além disso,
que, se $S$ é localmente noetheriano, a hipótese de planicidade fiel sobre $P$
equivale à hipótese de que
\[
  \widehat P_s=P\times_S\operatorname{Spec}(\widehat{\mathcal O}_s)
\]
seja fielmente plano sobre $\widehat{\mathcal O}_s$ para todo $s\in S$ (onde
$\widehat{\mathcal O}_s$ denota o completamento do anel local
$\mathcal O_s$), como se deduz do fato de que $\widehat{\mathcal O}_s$ é
fielmente plano sobre $\mathcal O_s$. Além disso, se $P$ é de tipo finito sobre
$S$ localmente noetheriano, o conjunto dos pontos $s$ que satisfazem a
condição anterior é construtível; portanto, se $S$ é um «pré-esquema de
Jacobson» (por exemplo, um esquema de tipo finito sobre um corpo ou, mais
geralmente, o espectro de um anel de Jacobson), basta verificar a
condição em questão nos pontos fechados de $S$. Isso nos reduz ao caso em
que a base é o espectro de um anel local completo $A$. Quando
$S=\operatorname{Spec}(A)$ ($A$ um anel local noetheriano completo) e $P$ é
de tipo finito sobre $S$, a planicidade fiel de $P/S$ equivale também à
existência de um $S'$ finito e plano sobre $S$ tal que $P'$ seja trivial; e, se
além disso $G$ é liso sobre $S$, pode-se supor que $S'$ seja étale sobre $S$.
Por conseguinte, se além disso o corpo residual de $A$ é algebricamente fechado
(«caso geométrico»), $P$ é fielmente plano sobre $A$ se e somente se é trivial.
Assim, pois, se $S$ é um pré-esquema algébrico sobre um corpo algebricamente
fechado e $G$ é liso e de tipo finito sobre $S$, vê-se que a condição de
planicidade fiel sobre $S$ equivale à condição de trivialidade analítica (SLF)
de Serre ([6], pp. 1--12).

Podem-se introduzir outros tipos, mais fortes, de condições sobre $P$ que
tenham a natureza de uma «trivialidade local». Diremos, em particular, que
$P$ é isotrivial (resp. estritamente isotrivial) se, para todo $s\in S$,
existem uma vizinhança aberta $U$ de $s$ e um morfismo finito e fielmente plano
(resp. um recobrimento étale sobrejetivo) $U'\longrightarrow U$ tais que
$P'=P\times_SU'$ seja trivial. (Nossa terminologia difere da de Serre [6], que
chama localmente isotrivial ao que nós chamamos estritamente isotrivial.) A
isotrivialidade estrita é especialmente útil se $G$ é liso sobre $S$, mas, em
contrapartida, é uma noção inadequada nos demais casos.

Se $G$ é afim sobre $S$, todo fibrado principal homogêneo $P$ sob $G$ é afim
pelo parágrafo 2, de onde, graças ao teorema 2, resulta a possibilidade de
«descer»

\src{326}
tais fibrados por morfismos fielmente planos e quase-compactos. Tomemos, em
particular, $G=\operatorname{Gl}(n)_S$, definido pela condição de que o
funtor de $S$-pré-esquemas
\[
  S'\longmapsto\operatorname{Hom}_S(S',G)
\]
(com valores na categoria dos grupos) se identifique com o funtor
\[
  \operatorname{Gl}(n)(S')
  =\operatorname{Gl}\bigl(n,H^0(S',\mathcal O_{S'})\bigr)
\]
de A, parágrafo 4.e. Utilizando os fatos (i) de que todo fibrado principal
homogêneo sob $G$ (resp. todo feixe localmente livre de posto $n$ sobre $S$) se
torna isomorfo ao objeto «trivial» $G$ (resp. $\mathcal O_S^n$) após uma mudança
fielmente plana e quase-compacta adequada de $S$; (ii) de que podem ser descidos
os objetos do tipo considerado (fibrados principais homogêneos sob $G$,
resp. feixes localmente livres de posto $n$) por tais morfismos; e, por fim,
(iii) de que o grupo de automorfismos do fibrado trivial sobre um $S'/S$ é
funtorialmente isomorfo ao grupo de automorfismos do feixe localmente livre
trivial de posto $n$ sobre $S'$, conclui-se «formalmente» que «dá no mesmo»
dar sobre $S$ (ou sobre um $S'/S$) um fibrado principal homogêneo de grupo $G$
ou dar aí um feixe localmente livre de posto $n$. (Mais precisamente, tem-se
uma equivalência de categorias fibradas.) Conclui-se, em particular:

\result{PROPOSIÇÃO}{6.1} Todo fibrado principal homogêneo de grupo
$\operatorname{Gl}(n)_S$ é localmente trivial.

Mediante argumentos conhecidos obtém-se o mesmo resultado para grupos
estruturais tais como $\operatorname{Sl}(n)_S$, $\operatorname{Sp}(n)_S$ e
produtos de tais grupos. Deduz-se também que, se $F$ é um subgrupo fechado
de $G=\operatorname{Gl}(n)_S$, plano sobre $S$, tal que existe o quociente
$G/F$, e se $G$ é um fibrado principal homogêneo isotrivial (resp.
estritamente isotrivial) sobre $G/F$, de grupo $F\times_S(G/F)$, então todo
fibrado principal homogêneo de grupo $F$ é isotrivial (resp. estritamente
isotrivial). Isso se aplica a todos os «grupos lineares» sobre $S$ utilizados
até agora e, em particular, quando $G=S\times_k\Gamma$, onde $S$ é um
pré-esquema sobre o corpo $k$ e $\Gamma$ um grupo linear no sentido clássico
e, em particular, liso, sobre $k$. Isso resolve, pois, para tais grupos, uma
questão de Serre (loc. cit.).

Assinalemos igualmente que, para a maioria dos grupos lisos sobre $S$ conhecidos,
sejam lineares ou não, e em todo caso para os da forma
$S\times_k\Gamma$ anterior, pode demonstrar-se que todo fibrado principal
homogêneo isotrivial é estritamente isotrivial. Isso resolve, em particular,
outra questão de Serre (loc. cit., pp. 1--14), dado que um fibrado principal
homogêneo obtido por uma descida à Cartier (cf. parágrafo 3, exemplo 2) é
manifestamente isotrivial.

\result{OBSERVAÇÃO}{} Uma das dificuldades essenciais nessas questões
(além da questão da existência dos esquemas quociente $G/F$) é a
falta de critérios de efetividade para um dado de descida por um morfismo
fielmente plano não finito.

\src{327}
\begin{center}
  \textbf{BIBLIOGRAFIA}
\end{center}

\begin{description}
  \item[{[1]}] Dieudonné (J.) e Grothendieck (Alexander). --- \emph{Éléments
    de géométrie algébrique}, a aparecer nas Publications
    mathématiques de l'Institut des Hautes Études Scientifiques.
  \item[{[2]}] Grauert (H.) und Remmert (R.). --- \emph{Komplexe Räume}, Math.
    Annalen, t. 136, 1958, pp. 245--318.
  \item[{[3]}] Grothendieck (Alexander). --- \emph{Géométrie formelle et
    géométrie algébrique}. Séminaire Bourbaki, t. 11, 1958/59, n.º 182.
  \item[{[4]}] Murre (J. P.). --- \emph{On a connectedness theorem for a
    birational transformation at a simple point}, Amer. J. Math., t. 80, 1958,
    pp. 3--15.
  \item[{[5]}] Serre (Jean-Pierre). --- \emph{Géométrie algébrique et géométrie
    analytique}, Ann. Inst. Fourier Grenoble, t. 6, 1955--56, pp. 1--42.
  \item[{[6]}] Serre (Jean-Pierre). --- \emph{Espaces fibrés algébriques},
    Séminaire Chevalley, t. 2, 1958: Anneaux de Chow et applications, n.º 1.
\end{description}

\par

\def\FgaUnitKey{e190}\clearpage
\chapter*{Errata da exposição n\textsuperscript{o} 190}
\addcontentsline{toc}{chapter}{Errata da exposição n\textsuperscript{o} 190}
\setcounter{footnote}{0}
\markboth{Errata da exposição 190}{Errata da exposição 190}
\sid{e190}
\src{300}

\noindent\textbf{3. Técnica de descida e teoremas de existência em geometria
algébrica [TDTE I]. I: Generalidades. Descida por morfismos fielmente planos.}

\begin{center}
[Seminário Bourbaki, t. 12, 1959/60, n\textsuperscript{o} 190, 29 p.]
\end{center}

\noindent\textbf{Página 1, linha 10.} --- Agora parece excessivo dizer que a
técnica de descida está «na base da maioria dos teoremas de
existência em geometria algébrica». Isto é verdadeiro em grande medida para as
técnicas não projetivas tratadas nas duas primeiras exposições
da série TDTE I a VI, mas não para as técnicas projetivas (TDTE IV, V,
VI).

\noindent\textbf{Página 5, linha 16.} --- Não é necessário supor que $\alpha$
seja um morfismo de $\mathcal F$-descida.

\noindent\textbf{Página 20, observação.} --- Um morfismo
$S'\longrightarrow S$ quase-finito, étale e sobrejetivo, ou um morfismo
$\operatorname{Spec}(\widehat A)\longrightarrow\operatorname{Spec}(A)$, não é
sempre um morfismo de descida estrito, mesmo que $A$ seja um anel local de uma
curva algébrica sobre um corpo algebraicamente fechado $k$ e
$S=\operatorname{Spec}(A)$. Assim, pode-se encontrar um morfismo próprio e liso
$f:X\longrightarrow S$, que faz de $X$ um fibrado principal de base $S$ sob
uma curva elíptica $E$, tal que $f':X'\longrightarrow S'$ seja projetivo, mas
$f$ não o seja. É, portanto, também um exemplo de fibrado
principal homogêneo não isotrivial sob um esquema abeliano.

\noindent\textbf{Página 26, linha 9.} --- Leia-se «Chow--Lang» em vez de
«Igusa--Lang».

Para diversos detalhes acerca da teoria da descida, vejam-se SGA VI, VII e
VIII.

\par

\def\FgaUnitKey{195}\unit{Exposição 195 -- Técnica de descida II}
\sid{195}
\markboth{Exposição 195}{Exposição 195}
\src{369}
\seminar{12.\textsuperscript{o} ano, 1959/60, n.\textsuperscript{o} 195\\Fevereiro de 1960}
\paperhead{TÉCNICA DE DESCIDA E TEOREMAS DE EXISTÊNCIA EM GEOMETRIA ALGÉBRICA\\
II. O TEOREMA DE EXISTÊNCIA NA TEORIA FORMAL DOS MÓDULOS}%
  {por Alexander GROTHENDIECK}

\fsection{A}{Funtores representáveis e pro-representáveis}

\subsection*{1. Funtores representáveis}
\addcontentsline{toc}{subsection}{1. Funtores representáveis}

Seja $\mathcal C$ uma categoria. Para todo $X\in\mathcal C$, seja $h_X$ o
funtor contravariante de $\mathcal C$ na categoria $(\operatorname{Ens})$
dos conjuntos,
\[
  h_X:\mathcal C\longrightarrow(\operatorname{Ens}),
\]
definido pela fórmula
\[
  h_X(Y)=\operatorname{Hom}(Y,X).
\]
Um morfismo $X\longrightarrow X'$ em $\mathcal C$ induz de maneira evidente um
homomorfismo de funtores $h_X\longrightarrow h_{X'}$: $h_X$ é um funtor
covariante em $X$; isto é, temos definido um funtor covariante canônico
\[
  h:\mathcal C\longrightarrow
  \operatorname{Hom}\bigl(\mathcal C^o,(\operatorname{Ens})\bigr)
\]
de $\mathcal C$ na categoria dos funtores covariantes da categoria dual
$\mathcal C^o$ de $\mathcal C$ na categoria dos conjuntos. Recordemos
então:

\result{PROPOSIÇÃO}{1.1} Este funtor $h$ é \emph{plenamente fiel}; em outros
termos, para todo par $X,X'$ de objetos de $\mathcal C$, a aplicação
natural
\[
  \operatorname{Hom}(X,X')\longrightarrow
  \operatorname{Hom}(h_X,h_{X'})
\]
é bijetiva.

Em particular, se um funtor
$F\in\operatorname{Hom}(\mathcal C^o,(\operatorname{Ens}))$ é isomorfo a um
funtor da forma $h_X$, então $X$ fica determinado salvo isomorfismo
único. Diz-se nesse caso que o funtor $F$ é \emph{representável}. A
proposição anterior significa que o funtor canônico $h$ define uma
equivalência da categoria $\mathcal C$ com a subcategoria plena de
$\operatorname{Hom}(\mathcal C^o,(\operatorname{Ens}))$ formada pelos
funtores representáveis. Este fato fundamenta a noção de «solução de um
problema universal»: tal problema consiste sempre em examinar se um funtor
dado (contravariante como aqui, ou covariante no caso dual) de $\mathcal C$
em $(\operatorname{Ens})$ é representável.

\src{370}
Observemos além disso que, pela própria definição de produto em uma categoria
[1], o funtor $h:X\longmapsto h_X$ comuta com os produtos sempre que estes
existam (e, mais geralmente, com os limites projetivos finitos ou infinitos;
em particular, com os produtos fibrados, com a formação de «núcleos» [2],
etc., sempre que existam): tem-se um isomorfismo de funtores
\[
  h_{X\times X'}\xrightarrow{\sim}h_X\times h_{X'}
\]
sempre que exista $X\times X'$; isto é, têm-se bijeções funtoriais
em $Y$
\[
  h_{X\times X'}(Y)\xrightarrow{\sim}h_X(Y)\times h_{X'}(Y).
\]
Em particular, dar um morfismo
\[
  X\times X'\longrightarrow X''
\]
em $\mathcal C$ (isto é, uma «lei de composição» em $\mathcal C$ entre
$X,X',X''$) equivale a dar um morfismo
\[
  h_{X\times X'}=h_X\times h_{X'}\longrightarrow h_{X''},
\]
isto é, a dar, para todo $Y\in\mathcal C$, uma lei de composição de
conjuntos
\[
  h_X(Y)\times h_{X'}(Y)\longrightarrow h_{X''}(Y)
\]
de tal modo que, para todo morfismo $Y\longrightarrow Y'$ em $\mathcal C$, o
sistema de aplicações entre conjuntos
\[
  h_{X^{(i)}}(Y)\longleftarrow h_{X^{(i)}}(Y')
  \qquad(\text{para }i=0,1,2)
\]
seja um morfismo com respeito às duas leis de composição correspondentes a
$Y$ e $Y'$. Vê-se assim que a noção de estrutura de «$\mathcal C$-grupo»,
«$\mathcal C$-anel», etc. sobre um objeto $X$ de $\mathcal C$ se exprime do
modo mais cômodo (tanto na teoria como na prática) dizendo que, para
todo $Y\in\mathcal C$, o conjunto $h_X(Y)$ está provido de uma lei de grupo
(resp. de anel, etc.) no sentido usual, e que as aplicações
$h_X(Y)\longleftarrow h_X(Y')$ correspondentes a morfismos
$Y\longrightarrow Y'$ são homomorfismos de grupos (resp. de anéis, etc.).
Esta é, por exemplo, a forma mais intuitiva e cômoda de definir os diversos
grupos clássicos $G_a$, $G_m$, $\operatorname{Gl}(n)$, etc. sobre um pré-esquema
base arbitrário $S$ e de escrever as relações clássicas entre esses grupos;
de dotar o esquema afim $V(\mathcal F)$ sobre $S$, definido por um feixe
quase coerente $\mathcal F$, de uma estrutura de «fibrado vetorial»; e de
definir e estudar as diversas variedades de bandeiras associadas
(grassmannianas, fibrados projetivos), etc. O yoga geral consiste
simplesmente em identificar, mediante o funtor canônico $h$, os objetos de
$\mathcal C$ com certos funtores contravariantes: os funtores representáveis,

\src{371}
de $\mathcal C$ na categoria dos conjuntos.

O procedimento habitual de inverter as setas, necessário, por exemplo, no
caso dos esquemas afins para passar da linguagem geométrica à da álgebra
comutativa, leva a dualizar as considerações anteriores e, em
particular, a introduzir também os funtores covariantes representáveis
$\mathcal C\longrightarrow(\operatorname{Ens})$; isto é, os da forma
\[
  Y\longmapsto\operatorname{Hom}(X,Y)=h'_X(Y).
\]

\subsection*{2. Funtores pro-representáveis, pro-objetos}
\addcontentsline{toc}{subsection}{2. Funtores pro-representáveis, pro-objetos}

Seja $\widetilde X=(X_i)_{i\in I}$ um sistema projetivo de objetos de
$\mathcal C$; a ele corresponde um funtor covariante
\[
  h'_{\widetilde X}=\varinjlim_i h'_{X_i},
\]
mais explicitamente,
\[
  h'_{\widetilde X}(Y)=\varinjlim_i h'_{X_i}(Y)
  =\varinjlim_i\operatorname{Hom}(X_i,Y),
\]
de $\mathcal C$ em $(\operatorname{Ens})$. Um funtor de $\mathcal C$ em
$(\operatorname{Ens})$ que seja isomorfo a um funtor desse tipo com $I$
filtrante chama-se \emph{pro-representável}. Pelo número anterior, estes são
exatamente os funtores isomorfos a limites indutivos filtrantes de funtores
representáveis. Seja $\widetilde X'=(X'_j)_{j\in J}$ um segundo sistema
projetivo filtrante em $\mathcal C$ (construído sobre um segundo conjunto
preordenado filtrante de índices, $J$); verifica-se facilmente que há uma
bijeção canônica (que generaliza a proposição 1.1)
\[
  \operatorname{Hom}(h'_{\widetilde X'},h'_{\widetilde X})
  =\varprojlim_j\varinjlim_i\operatorname{Hom}(X_i,X'_j).
\]
Isso leva a introduzir a categoria $\operatorname{Pro}(\mathcal C)$ dos
pro-objetos de $\mathcal C$; seus objetos são os sistemas projetivos de objetos
de $\mathcal C$ (sobre conjuntos variáveis de índices preordenados filtrantes),
e se $\widetilde X=(X_i)_{i\in I}$ e
$\widetilde X'=(X'_j)_{j\in J}$ são dois objetos dessa classe, põe-se
\[
  \operatorname{ProHom}(\widetilde X,\widetilde X')
  =\varprojlim_j\varinjlim_i\operatorname{Hom}(X_i,X'_j),
\]
sendo evidente a composição dos pro-homomorfismos. Por sua própria
construção, $\widetilde X\longmapsto h'_{\widetilde X}$ pode ser considerado
como um funtor contravariante em $\widetilde X$, que estabelece uma equivalência
entre a categoria dual da categoria $\operatorname{Pro}(\mathcal C)$ dos
pro-objetos de $\mathcal C$ e a categoria dos funtores covariantes
pro-representáveis de $\mathcal C$ em $(\operatorname{Ens})$. Naturalmente, um
objeto $X$ de $\mathcal C$ define canonicamente um pro-objeto, denotado

\src{372}
também por $X$, de modo que $\mathcal C$ é equivalente a uma subcategoria
plena de $\operatorname{Pro}(\mathcal C)$. Se então
$\widetilde X=(X_i)_{i\in I}$ é um pro-objeto qualquer de $\mathcal C$, tem-se
(com a identificação anterior)
\[
  \widetilde X=\varprojlim_i X_i,
\]
tomando-se o limite projetivo em $\operatorname{Pro}(\mathcal C)$
(visto que $h'_{\widetilde X}=\varinjlim_i h'_{X_i}$). Convém observar que,
mesmo quando o limite projetivo dos $X_i$ exista em $\mathcal C$, em
geral não será isomorfo ao limite projetivo $\widetilde X$ em
$\operatorname{Pro}(\mathcal C)$, como já resulta evidente quando $\mathcal C$
é a categoria dos conjuntos. Note-se que, por definição,
$\varprojlim_{\mathcal C}X_i=L$ fica definido pela condição de que o
funtor em $Y\in\mathcal C$
\[
  \varprojlim_i\operatorname{Hom}_{\mathcal C}(Y,X_i)
  =\operatorname{Hom}_{\operatorname{Pro}(\mathcal C)}(Y,\widetilde X)
\]
com valores em $(\operatorname{Ens})$ seja representável mediante $L$, isto é,
seja isomorfo a $\operatorname{Hom}_{\mathcal C}(Y,L)$; por conseguinte,
$\varprojlim_{\mathcal C}X_i$ já está definido em termos do pro-objeto
$\widetilde X$ e, mais precisamente, depende funtorialmente do pro-objeto
$\widetilde X$ sempre que esteja definido; não há, pois, inconveniente em
denotá-lo por $\varprojlim_{\mathcal C}(\widetilde X)$. Se os limites
projetivos existem sempre em $\mathcal C$,
$\varprojlim_{\mathcal C}(\widetilde X)$ é um funtor de
$\operatorname{Pro}(\mathcal C)$ em $\mathcal C$, e há um homomorfismo
canônico de funtores
$\varprojlim_{\mathcal C}(\widetilde X)\longrightarrow\widetilde X$.

Como todo funtor, que suporemos covariante para fixar as ideias,
\[
  F:\mathcal C\longrightarrow\mathcal C'
\]
se prolonga de maneira evidente a um funtor
\[
  \operatorname{Pro}(F):\operatorname{Pro}(\mathcal C)
  \longrightarrow\operatorname{Pro}(\mathcal C'),
\]
segue-se que, se em $\mathcal C'$ sempre existem os limites projetivos, $F$
define também canonicamente um funtor composto
\[
  \overline F:\operatorname{Pro}(\mathcal C)\longrightarrow\mathcal C',
  \qquad
  \overline F=\varprojlim_{\mathcal C'}\circ\operatorname{Pro}(F),
\]
que transforma $\widetilde X=(X_i)_{i\in I}$ em
$\varprojlim_{\mathcal C'}F(X_i)$.

Um pro-objeto $\widetilde X$ chama-se \emph{pro-objeto estrito} se é isomorfo a
um pro-objeto $(X_i)_{i\in I}$ cujos morfismos de transição
$X_i\longrightarrow X_j$ são epimorfismos; um funtor definido por tal objeto
chama-se \emph{estritamente pro-representável}. Pode-se exigir, além disso, que $I$
seja ordenado filtrante e que todo epimorfismo $X_i\longrightarrow X'$ seja
equivalente a um epimorfismo

\src{373}
$X_i\longrightarrow X_j$ para um $j\in I$ adequado (determinado de maneira
única por essa condição). Nessas condições, o sistema projetivo
$(X_i)_{i\in I}$ fica determinado salvo isomorfismo único (no sentido usual
dos isomorfismos de sistemas projetivos). Segue-se que em
$\operatorname{Pro}(\mathcal C)$ existe sempre o limite projetivo de um
sistema projetivo de pro-objetos estritos $\widetilde X^{(\alpha)}$ e que, com
as notações anteriores para $F$, $\overline F$, ter-se-á
\[
  \overline F\left(\varprojlim_\alpha\widetilde X^{(\alpha)}\right)
  =\varprojlim_{\alpha,\mathcal C'}
  \overline F\left(\widetilde X^{(\alpha)}\right).
\]
Em particular, se todo pro-objeto de $\mathcal C$ é estrito (cf. o número
seguinte), o funtor prolongado $\overline F$ comuta com os limites
projetivos.

\subsection*{3. Caracterização dos funtores pro-representáveis}
\addcontentsline{toc}{subsection}{3. Caracterização dos funtores pro-representáveis}

Sejam $\mathcal C$, $\mathcal C'$ duas categorias nas quais existem os limites
projetivos finitos (isto é, relativos a conjuntos preordenados finitos, não
necessariamente filtrantes), ou, o que é equivalente, nas quais existem os
produtos finitos e os produtos fibrados finitos (o que implica, em
particular, a existência de um «objeto final» $e$, tal que
$\operatorname{Hom}(X,e)$ se reduza a um elemento para todo $X$). Seja $F$ um
funtor covariante de $\mathcal C$ em $\mathcal C'$. As condições seguintes
são equivalentes:

\begin{enumerate}
\item[(i)] $F$ comuta com os limites projetivos finitos;
\item[(ii)] $F$ comuta com os produtos finitos e os produtos fibrados finitos;
\item[(iii)] $F$ comuta com os produtos finitos e, para todo diagrama exato
\[
  X\longrightarrow X'\rightrightarrows X''
\]
em $\mathcal C$ ([3], A, definição 2.1), o diagrama transformado por $F$
\[
  F(X)\longrightarrow F(X')\rightrightarrows F(X'')
\]
é exato.
\end{enumerate}

Diz-se então que $F$ é \emph{exato à esquerda}.

De agora em diante, supõe-se que em $\mathcal C$ existem os limites projetivos
finitos. Das definições resulta então imediatamente que um funtor
representável é exato à esquerda e, passando ao limite, que um funtor
pro-representável é exato à esquerda.

Para obter uma recíproca, seja
\[
  F:\mathcal C\longrightarrow(\operatorname{Ens})
\]
um funtor covariante, e sejam $X\in\mathcal C$ e $\xi\in F(X)$. Diz-se que
$\xi$ (ou o par \src{374}$(X,\xi)$) é \emph{minimal} se, para todo
$X'\in\mathcal C$ e $\xi'\in F(X')$, e todo monomorfismo estrito ([3], A, 2)
$u:X'\longrightarrow X$ tal que $\xi=F(u)(\xi')$, $u$ é um isomorfismo. Por
outro lado, diz-se que um par $(X,\xi)$ \emph{domina} $(X'',\xi'')$
($\xi\in F(X)$, $\xi''\in F(X'')$) se existe um morfismo
$v:X\longrightarrow X''$ tal que $\xi''=F(v)(\xi)$; se $\xi$ é minimal e $F$
é exato à esquerda, este morfismo $v$ é então único; se $\xi''$ é
minimal, $v$ é sobrejetivo. Deduz-se facilmente a proposição seguinte:

\result{PROPOSIÇÃO}{3.1} Para que $F$ seja estritamente pro-representável, é
necessário e suficiente que satisfaça as duas condições:
\begin{enumerate}
\item[(i)] $F$ é exato à esquerda;
\item[(ii)] todo par $(X,\xi)$, com $\xi\in F(X)$, é dominado por um par
minimal.
\end{enumerate}

Esta última condição é vazia se todo objeto de $\mathcal C$ é artiniano
(tomando um subobjeto $X'$ de $X$ minimal entre aqueles para os quais existe
$\xi'\in F(X')$ tal que $\xi$ seja a imagem de $\xi'$). Daí:

\result{COROLÁRIO}{} Seja $\mathcal C$ uma categoria cujos objetos são
artinianos e na qual existem os limites projetivos finitos. Então os
funtores pro-representáveis de $\mathcal C$ em $(\operatorname{Ens})$ são
exatamente os funtores exatos à esquerda, e são inclusive estritamente
pro-representáveis.

Este último fato significa também que \emph{todo pro-objeto de $\mathcal C$
é estrito}.

\subsection*{4. Exemplo: grupos de tipo galoisiano, grupos proalgébricos}
\addcontentsline{toc}{subsection}{4. Grupos de tipo galoisiano, grupos proalgébricos}

Se $\mathcal C$ é a categoria dos grupos finitos ordinários,
$\operatorname{Pro}(\mathcal C)$ é equivalente à categoria dos grupos
topológicos compactos totalmente desconexos. São grupos desse tipo e suas
generalizações, obtidas substituindo os grupos finitos ordinários por
esquemas em grupos finitos sobre um pré-esquema base dado (por exemplo, os
grupos algébricos finitos sobre um corpo $k$), que farão as vezes de
grupos fundamentais, de homotopia e de homologia absolutos e relativos para
os pré-esquemas. Em todos estes exemplos se aplica o corolário da
proposição 3.1, e é precisamente por meio do funtor associado que se deve
definir $\pi_1$ [2]. O mesmo ocorre partindo da categoria dos grupos
algébricos ou quase algébricos sobre um corpo (ou, mais geralmente, sobre um
pré-esquema noetheriano): obtêm-se os «grupos proalgébricos» de Serre [4].

\subsection*{5. Exemplo: «variedades formais»}
\addcontentsline{toc}{subsection}{5. Exemplo: variedades formais}

Sejam $\Lambda$ um anel noetheriano e $\mathcal C$ a categoria das
$\Lambda$-álgebras $A$ que são módulos artinianos de tipo finito sobre
$\Lambda$ (ou, mais brevemente,

\src{375}
as $\Lambda$-álgebras artinianas). São satisfeitas as condições do corolário da
proposição 3.1. A categoria $\operatorname{Pro}(\mathcal C)$ é aqui
equivalente à categoria das álgebras topológicas $O$ sobre $\Lambda$
isomorfas a limites projetivos topológicos
\[
  O=\varprojlim_i O_i
\]
de álgebras $O_i\in\mathcal C$, isto é, cuja topologia é linear, separada e
completa, e tal que para todo ideal aberto $J_i$ de $O$, $O/J_i$ seja uma
álgebra artiniana sobre $\Lambda$. O funtor
$\mathcal C\longrightarrow(\operatorname{Ens})$ associado a tal álgebra não é
outro que
\begin{align*}
  F(A)=h'_O(A)
  &=\{\text{homomorfismos contínuos de $\Lambda$-álgebras topológicas }
       O\longrightarrow A\}\\
  &=\varinjlim_i
    \operatorname{Hom}_{\Lambda\text{-álgebras}}(O_i,A).
\end{align*}

Observe-se, além disso, que a categoria $\mathcal C$ é essencialmente o produto
das categorias análogas correspondentes aos completamentos dos
anéis locais $\Lambda_{\mathfrak m}$, para os ideais maximais
$\mathfrak m$ de $\Lambda$; portanto, se desejado, pode-se restringir ao caso
em que $\Lambda$ seja um anel local completo dessa classe. Em qualquer caso,
$O$ se decompõe canonicamente no produto topológico de componentes
locais, correspondentes aos «pontos» do esquema formal [2] definido por
$O$. Um ponto dessa classe é definido por um elemento $\xi$ de algum
$F(K)$, onde $K\in\mathcal C$ é um corpo (por exemplo, o corpo residual
da componente local considerada); dois pares $(\xi,K)$ e $(\xi',K')$
definem o mesmo ponto se, e somente se, são dominados por um mesmo
$(\xi'',K'')$, ou se dominam um mesmo $(\xi''',K''')$. (Se os
$\Lambda/\mathfrak m$ são algebricamente fechados, basta tomar a soma
disjunta dos $F(\Lambda/\mathfrak m)$.)

É importante dar condições para que a componente local $O_\xi$ de $O$
correspondente a um $\xi\in F(K)$ seja um anel noetheriano. Quando $\Lambda$
é um anel local completo (noetheriano, recordemos), isso equivale a dizer
que $O_\xi$ é isomorfo a um anel quociente de um anel de séries formais
$\Lambda[[t_1,\ldots,t_n]]$ sobre $\Lambda$. Para dar um critério desse tipo,
introduzamos (para todo anel $A$) a $A$-álgebra $I_A$ dos «números
duais» de $A$ mediante
\[
  I_A=A[t]/t^2A[t].
\]
Seja $\varepsilon:I_A\longrightarrow A$ o homomorfismo de aumentação; este define
(se $A\in\mathcal C$) uma aplicação
\[
  F(\varepsilon):F(I_A)\longrightarrow F(A),
\]

\src{376}
e, utilizando que $F$ é exato à esquerda, define-se de maneira
intrínseca uma estrutura de $A$-módulo na parte
\[
  F(I_A,\xi)=F(\varepsilon)^{-1}(\xi)
\]
de $F(I_A)$ formada pelos $\xi'\in F(I_A)$ que «se reduzem a $\xi$»;
utilizando a forma explícita de $F$ em termos de $O$, verifica-se que este
$A$-módulo se identifica com
\[
  \operatorname{Hom}_{\Lambda}(\mathfrak m_\xi/\mathfrak m_\xi^2,A),
\]
onde $\mathfrak m_\xi$ é o núcleo do homomorfismo
$\xi:O\longrightarrow A$, isto é, se $A$ é um corpo, o ideal maximal da
componente local $O_\xi$ de $O$. Deduz-se imediatamente a proposição
seguinte:

\result{PROPOSIÇÃO}{5.1} Seja $\xi\in F(K)$, onde $K\in\mathcal C$ é um
corpo. Para que a componente local correspondente $O_\xi$ de $O$ seja um
anel noetheriano, é necessário e suficiente que o conjunto $F(I_K,\xi)$ de
elementos de $F(I_K)$ que se reduzem a $\xi$ seja um espaço vetorial
de dimensão finita sobre $K$. Nessas condições, tem-se um isomorfismo
canônico
\[
  F(I_K,\xi)=
  \operatorname{Hom}\!\left(
    \frac{\mathfrak m_\xi/\mathfrak m_\xi^2}
         {\operatorname{im}(\mathfrak n_\xi/\mathfrak n_\xi^2)},K
  \right),
\]
onde $\mathfrak n_\xi$ é o ideal maximal de $\Lambda$, núcleo do
homomorfismo $\Lambda\longrightarrow K$.
\footnote{O texto impresso dá aqui uma fórmula que só é correta quando
$\Lambda$ é um corpo. A errata de 1962 prescreve, no caso geral,
o quociente de $\mathfrak m_\xi/\mathfrak m_\xi^2$ pela imagem de
$\mathfrak n_\xi/\mathfrak n_\xi^2$.}
Em particular, a dimensão do $K$-espaço vetorial $F(I_K,\xi)$ é igual à
dimensão do espaço vetorial assim corrigido sobre o corpo
$O_\xi/\mathfrak m_\xi=K(\xi)$.

Suponhamos que $O_\xi$ seja noetheriano e, para simplificar a escrita, que
$\Lambda$ seja local completo e $O=O_\xi$. Diz-se que $O$ é \emph{liso sobre
$\Lambda$} se $O$ é uma álgebra finita e étale sobre a álgebra completada do
localizado de $\Lambda[t_1,\ldots,t_n]$ em um ideal maximal desse anel que
induz o ideal maximal de $\Lambda$; quando a extensão residual de $O$ sobre
$\Lambda$ é trivial (por exemplo, se o corpo residual de $\Lambda$ é
algebricamente fechado), isso equivale a dizer que o próprio $O$ é isomorfo a uma
álgebra de séries formais desse tipo.
\footnote{A errata de 1962 esclarece que esta definição de «liso sobre
$\Lambda$» só é correta quando a extensão residual $k'/k$ é separável;
no caso geral, remete a SGA III, 1.1.}
Finalmente, se já não se supõe necessariamente que $O$ seja noetheriano, dir-se-á
também que $O$ é liso sobre $\Lambda$ se $O$ é isomorfo a um limite
projetivo topológico de $\Lambda$-álgebras quociente que são noetherianas e
lisas sobre $\Lambda$ no sentido anterior. Isso se generaliza imediatamente ao caso
em que $\Lambda$ e $O$ já não são supostos locais. Dito isso, tem-se a
proposição seguinte:

\result{PROPOSIÇÃO}{5.2} Para que $O$ seja liso sobre $\Lambda$, é necessário
e suficiente que o funtor associado $F$ transforme epimorfismos em
epimorfismos.

Isto significa que, para todo homomorfismo sobrejetivo
$A\longrightarrow A'$ em $\mathcal C$, $F(A)\longrightarrow F(A')$ é
sobrejetivo. Naturalmente, basta verificar esta condição quando $A$ é local
e, procedendo passo a passo, quando o ideal de $A$ que é núcleo de
$A\longrightarrow A'$ é anulado pelo ideal maximal de $A$. Na prática,
isso reduz a questão a verificar que certa obstrução, ligada a invariantes
infinitesimais

\src{377}
da situação que dá origem ao funtor $F$, seja nula, problema de natureza
cohomológica.

Para terminar, digamos algumas palavras, no contexto anterior, sobre os
\emph{anéis de definição}. Seja ainda $F$ um funtor de $\mathcal C$ em
$(\operatorname{Ens})$, pro-representável mediante uma $\Lambda$-álgebra
topológica $O$. Então, para todo $A\in\mathcal C$ e todo $\xi\in F(A)$,
existe um menor subanel $A'$ de $A$ tal que $\xi$ seja imagem de um elemento
$\xi'$ de $F(A')$ (que fica então determinado de maneira única): com efeito,
basta interpretar $\xi$ mediante um homomorfismo de $O$ em $A$ e tomar como
$A'$ a imagem de $O$ por este último. Diz-se então que $A'$ é o anel
de definição do objeto $\xi\in F(A)$. Se $u:A\longrightarrow B$ é um
homomorfismo de álgebras e $\eta=F(u)(\xi)$, o anel de definição de $\eta$
é a imagem por $u$ do anel de definição de $\xi$. Quando se parte de um
funtor $F$ de $\mathcal C$ em $(\operatorname{Ens})$, a existência dos
anéis de definição e as propriedades que acabamos de dar são, com pequenas
ressalvas, equivalentes a que $F$ seja pro-representável; isto é, em
geral estão longe de ser triviais.

\fsection{B}{Os dois teoremas de existência}

Conservemos as notações de A, parágrafo 5, e partamos de um funtor covariante
\[
  F:\mathcal C\longrightarrow(\operatorname{Ens});
\]
buscam-se critérios manejáveis para que $F$ seja pro-representável, isto é,
exprimível mediante uma $\Lambda$-álgebra $O$ como acima. Em virtude do
corolário de A, proposição 3.1, para isso é necessário e suficiente que $F$ seja
exato à esquerda. No estado atual da técnica de descida (cf. as
questões formuladas em [3], p. 9), esse critério não pode ser verificado
diretamente sob essa forma nos casos mais importantes, e são necessários
critérios aparentemente menos exigentes.

\result{TEOREMA}{1} Para que o funtor $F$ seja pro-representável, é necessário
e suficiente que satisfaça as duas condições seguintes:
\begin{enumerate}
\item[(i)] $F$ comuta com os produtos finitos;
\item[(ii)] para toda álgebra $A\in\mathcal C$ e todo homomorfismo
$A\longrightarrow A'$ em $\mathcal C$ tal que o diagrama
\[
  A\longrightarrow A'\rightrightarrows A'\otimes_A A'
\]
seja exato ([3], A, definição 1.2), o diagrama transformado
\[
  F(A)\longrightarrow F(A')\rightrightarrows F(A'\otimes_A A')
\]

\src{378}
é exato.
\end{enumerate}
Além disso, basta verificar (ii) quando $A$ é local e, ademais, em um dos
dois casos seguintes:
\begin{enumerate}
\item[(a)] $A'$ é um módulo livre sobre $A$;
\item[(b)] o módulo quociente $A'/A$ é um $A$-módulo de comprimento 1.
\end{enumerate}

A demonstração deste teorema é bastante delicada e não pode ser esboçada aqui.
Limitemo-nos a assinalar que ela repousa essencialmente em um estudo das relações
de equivalência (no sentido das categorias) no espectro de uma álgebra
artiniana (estudo que ainda suscita vários problemas cuja solução parece
indispensável ao desenvolvimento ulterior da teoria).

Nas aplicações, a verificação da condição (i) é sempre trivial.
A verificação de (ii) decompõe-se em duas partes: o caso em que $A'$ é um $A$-módulo livre
corresponde à técnica de descida por morfismos planos ([3], teoremas 1, 2 e
3) e não oferece dificuldade. Para tratar o caso (b), será utilizado o resultado
seguinte:

\result{TEOREMA}{2} Sejam $A$ um anel local artiniano de ideal maximal
$\mathfrak m$ e $A'$ uma $A$-álgebra que contém a $A$, tal que
$\mathfrak mA'\subset A$ e que
\[
  A\longrightarrow A'\rightrightarrows A'\otimes_A A'
\]
seja exato (o que ocorre, em particular, se $A'/A$ é um $A$-módulo de
comprimento 1). Seja $\mathcal F$ a categoria fibrada ([3], A, definição 1.1) de
feixes quase coerentes e planos sobre pré-esquemas variáveis. Então o
morfismo $\operatorname{Spec}(A')\longrightarrow\operatorname{Spec}(A)$ é um
morfismo de $\mathcal F$-descida estrito ([3], A, definição 1.7).

Em outros termos, o dado de um $A$-módulo plano $M$ é completamente
equivalente ao dado de um $A'$-módulo plano $M'$, provido de um
$A'\otimes_AA'$-isomorfismo de $M'\otimes_AA'$ sobre $A'\otimes_AM'$ que
satisfaz a condição usual de transitividade para um dado de descida (loco
citato).

O teorema 2 demonstra-se provando primeiro que
\[
  H^i(A'/A,G_a)=0\qquad\text{para }i\geq 1
\]
([3], A, 4(e)), pois a hipótese $\mathfrak mA'\subset A$ permite reduzir-se
facilmente ao caso em que $A$ é um corpo $k$ (a saber, $A/\mathfrak m$). Se
Aplicam-se então as equivalências assinaladas em [3], página 16.

\fsection{C}{Aplicações a alguns casos particulares}

\subsection*{1. Observações gerais sobre os funtores representados por pré-esquemas}
\addcontentsline{toc}{subsection}{1. Observações gerais sobre os funtores representados por pré-esquemas}

Seja $S$ um pré-esquema localmente noetheriano. Diz-se que um pré-esquema $X$
sobre $S$ é \src{379}localmente de tipo finito sobre $S$ se, para todo
$x\in X$ que se projeta em $y\in S$,\footnote{O texto impresso diz aqui
« $y\in Y$ » e volta a empregar $Y$ ao final do parágrafo. Como o morfismo
considerado é $X\longrightarrow S$, lê-se $S$ em ambas as ocorrências.} existem
uma vizinhança afim de $y$, de anel $A$, e uma vizinhança afim de $x$ situada sobre
o anterior, de anel $B$, tais que $B$ seja uma álgebra de tipo finito sobre
$A$. Encontram-se muitos exemplos importantes de pré-esquemas localmente de
tipo finito sobre $S$ que não são de tipo finito sobre $S$, como soluções de
problemas universais clássicos; assim, importa poder considerar o esquema de
Picard de uma curva como a reunião de uma infinidade de componentes conexas
(que não devem ser confundidas com a componente conexa do elemento neutro, isto
é, a «variedade de Picard»). Em certas ocasiões, é então conveniente
trabalhar na categoria $\mathcal C$ dos pré-esquemas localmente de tipo
finito sobre $S$, para nela examinar a questão da representabilidade de
um funtor contravariante $F$. O objetivo principal destas exposições é
desenvolver uma técnica geral que permita reconhecer se tal funtor $F$ é
representável e estudar as propriedades do $S$-pré-esquema $X$ correspondente
a partir das de $F$. Assinalemos de passagem que, nesse estudo, aparecem
exemplos não patológicos de pré-esquemas sobre $S$ não separados sobre $S$, em
particular como «pré-esquemas de Picard» de excelentes $S$-esquemas; não devem,
portanto, banir-se da geometria algébrica os pré-esquemas que não são
esquemas.

Seja $X$ um pré-esquema localmente de tipo finito sobre $S$, e seja
\[
  F(Y)=\operatorname{Hom}_S(Y,X)
\]
o funtor contravariante associado. Pode-se considerar a restrição $F_0$ de
$F$ à subcategoria $\mathcal C_0$ de $\mathcal C$ formada pelos
pré-esquemas $Y$ sobre $S$ que são artinianos e finitos sobre $S$: quando
$S=\operatorname{Spec}(\Lambda)$, $\mathcal C_0$ é, pois, a categoria dual
da categoria das $\Lambda$-álgebras artinianas considerada em B. Se
$Y=\operatorname{Spec}(A)$, onde $A$ é um anel local artiniano, $Y$
reduz-se a um único ponto $y$ situado sobre um ponto fechado $s$ de $S$, e um
$S$-homomorfismo de $Y$ em $X$ (isto é, um elemento de $F(Y)$) fica definido
pelo dado de um ponto $x\in X$ sobre $s$ e de um
$\mathcal O_{s,S}$-homomorfismo de $\mathcal O_{x,X}$ em $A$. Se existe tal
homomorfismo, então $x$ é necessariamente um ponto fechado de $X$ (pois seu
corpo residual é algébrico sobre o de $s$). Isso mostra que a restrição
$F_0$ de $F$ às «$\Lambda$-álgebras artinianas» é pro-representável, e está
representada pela álgebra topológica sobre $\Lambda$ cujas componentes
locais são os completamentos $\widehat{\mathcal O}_{x,X}$ dos anéis locais
de $X$ nos pontos $x$ de $X$ que são fechados e estão situados sobre pontos
fechados de $S$. Isso mostra que apenas o conhecimento de $F_0$ fornece
informação precisa sobre a estrutura de $X$ (a saber, a estrutura dos
completamentos de seus anéis locais nos pontos indicados). Observe-se que,
mesmo quando $S$

\src{380}
é o espectro de um corpo algebricamente fechado, somente mediante a
consideração sistemática de «variedades» $Y$ tais que $\mathcal O_Y$ possa
ter elementos nilpotentes, em particular trabalhando com os espectros de
anéis artinianos locais, pode-se chegar à «formulação correta» dos
problemas universais clássicos e captar seu aspecto «infinitesimal».

Quando se parte de um funtor $F$ dado de antemão e se quer decidir se ele é
representável, o estudo do funtor $F_0$ (com o auxílio dos teoremas 1 e 2)
fornecerá indicações quase completas: ou, como ocorre com frequência
(verificando simplesmente, por exemplo, a natureza dos conjuntos
$F(I_K,\xi)$ e seu comportamento funtorial, cf. A), constata-se que $F_0$ já não
é pro-representável (o que explica o fracasso das tentativas realizadas
até agora para definir de maneira mais ou menos natural variedades de módulos
para a classificação dos fibrados vetoriais de posto $>1$); ou consegue-se
verificar que $F_0$ é representável, mas os espaços vetoriais
$F(I_K,\xi)$ não são de dimensão finita, caso em que será preciso contentar-se com
a solução «formal»; ou, por fim, $F_0$ é representado por um
produto de anéis locais completos noetherianos, o que fornece indícios muito
sólidos de que o próprio $F$ é representável e, unido a propriedades análogas,
porém de natureza mais global, que pretendemos desenvolver mais adiante, bastará
sem dúvida para implicar que assim seja efetivamente. Finalmente, encontram-se
problemas geométricos interessantes (vejam-se os parágrafos 4 e 5 abaixo) nos
quais se dispõe apenas do funtor $F_0$ (que não provém de nenhum funtor «global»
$F$), e nos quais haverá motivos para se dar por satisfeito se se conseguir associar-lhe
um «esquema formal de módulos».

Para terminar estas generalidades, indiquemos como a teoria de esquemas
explica anomalias aparentes, como a superfície de Igusa $V$, cuja «variedade
de Picard» $P$ se reduz a um ponto e para a qual, no entanto, tem-se
\[
  H^1(V,\mathcal O_V)\ne 0 ;
\]
nesse caso, $P$ é um grupo «puramente infinitesimal» que não se reduz ao
grupo unidade, isto é, está definido por uma álgebra local $O$ de posto finito
sobre o corpo base $k$ e dotado de uma aplicação diagonal correspondente à
estrutura multiplicativa de $P$; se $\mathfrak m$ é o ideal maximal de
$O$, o dual de $\mathfrak m/\mathfrak m^2$ é canonicamente isomorfo a
$H^1(V,\mathcal O_V)$ (cf. o parágrafo 3 mais abaixo). Somente quando o grupo de
Picard é um grupo algébrico no sentido clássico, isto é, liso sobre o
corpo base $k$, a dimensão de $H^1(V,\mathcal O_V)$ (sempre igual à de
$\mathfrak m/\mathfrak m^2$) é igual à do grupo de Picard.

\subsection*{2. Esquemas $\underline{\operatorname{Hom}}_S(X,Y)$,
$\Pi_{X/S}Z$, $\underline{\operatorname{Aut}}_S(X)$, etc.}
\addcontentsline{toc}{subsection}{2. Esquemas Hom, Pi e Aut}

Sejam $X,Y$ dois pré-esquemas sobre $S$; para todo pré-esquema $T$ sobre $S$,

\src{381}
sejam $X_T=X\times_ST$, $Y_T=Y\times_ST$, e tomemos o conjunto
\[
  F(T)=\operatorname{Hom}_T(X_T,Y_T)
      =\operatorname{Hom}_S(X_T,Y)
      =\operatorname{Hom}_S(X\times_ST,Y)
\]
como funtor contravariante em $T$. Se ele for representável, denota-se por
$\underline{\operatorname{Hom}}_S(X,Y)$ o pré-esquema sobre $S$ que o
representa; tem-se, portanto, um isomorfismo funtorial
\[
  \operatorname{Hom}_S
  \bigl(T,\underline{\operatorname{Hom}}_S(X,Y)\bigr)
  \xrightarrow{\sim}
  \operatorname{Hom}_S(T\times_SX,Y).
\]
Há muitas variantes deste problema universal cuja solução se reduz a ele:
o pré-esquema dos $S$-automorfismos de um $S$-pré-esquema $X$ (que será um
pré-esquema em grupos), o pré-esquema dos $S$-homomorfismos de um
$S$-pré-esquema em grupos em outro (que será um pré-esquema em grupos comutativos
se o segundo esquema em grupos é comutativo), etc. Também se pode
generalizar a definição de $\underline{\operatorname{Hom}}_S(X,Y)$
considerando um pré-esquema $Z$ sobre o pré-esquema $X$ sobre $S$, e o funtor
\[
  F(T)=\operatorname{Hom}_{X_T}(X_T,Z_T)
\]
(conjunto das «seções» de $Z_T$ sobre $X_T$); se esse funtor for
representável, o $S$-pré-esquema que o representa será denotado por
$\Pi_{X/S}Z$; ter-se-á, pois, por definição, um isomorfismo funtorial
\[
  \operatorname{Hom}_S(T,\Pi_{X/S}Z)
  =\operatorname{Hom}_{X_T}(X_T,Z_T).
\]
Tomando $Z=Y\times_SX$, recupera-se
$\underline{\operatorname{Hom}}_S(X,Y)$. Dessas definições resulta para os
novos pré-esquemas assim introduzidos um formulário tão trivial quanto útil, que não
daremos aqui (visto que é válido em toda categoria na qual existem
produtos e produtos fibrados). Mais séria é a questão da existência de
esquemas do tipo $\underline{\operatorname{Hom}}_S(X,Y)$. Constata-se primeiro
que, para $X$ fixo, $\underline{\operatorname{Hom}}_S(X,Y)$ somente pode existir
para todo $Y$ sobre $S$ (ou $\Pi_{X/S}Z$ para todo $Z$ sobre $X$) se $X$ é plano
sobre $S$. Além disso, percebe-se que só é razoável esperar a existência de
uma solução, para $Y$ suficientemente geral, se $X$ é além disso próprio sobre
$S$. Por outro lado, parece que essas condições são suficientes para a
existência de $\underline{\operatorname{Hom}}_S(X,Y)$ e $\Pi_{X/S}Z$, com a
condição de impor, quando couber, uma hipótese do tipo «quase-projetivo»
sobre $Y/S$, respectivamente sobre $Z/X$; isso pode ser verificado, em qualquer
caso, em bastantes situações (por exemplo, quando $Y$ é afim sobre $S$; ou,
mediante construções diretas elementares, quando $X$ é finito sobre $S$).
Eis o que fornecem os teoremas 1 e 2:

\src{382}
\result{PROPOSIÇÃO}{2.1} Sejam $\Lambda$ um anel noetheriano, $X$ e $Y$ dois
pré-esquemas quaisquer sobre $\Lambda$, e tomemos o funtor
\[
  F(A)=\operatorname{Hom}_A(X_A,Y_A)
\]
na categoria $\mathcal C_0$ das $\Lambda$-álgebras artinianas. Se $X$ é
plano sobre $\Lambda$, este funtor é pro-representável.

Além disso, verifica-se que para todo $A\in\mathcal C_0$ e $\xi\in F(A)$ tem-se
um isomorfismo canônico
\[
  F(I_A,\xi)=H^0\!\left(
    X_A,
    \underline{\operatorname{Hom}}_{\mathcal O_{X_A}}
    \bigl(\xi^*(\Omega^1_{Y_A/A}),\mathcal O_{X_A}\bigr)
  \right),
\]
onde $\Omega^1_{Y_A/A}$ é o feixe de $1$-diferenciais de Kähler de $Y_A$
com respeito a $A$. Tomando $A$ como um corpo, obtém-se, utilizando A.5.1 e o
teorema de finitude de [2], o corolário seguinte:

\result{COROLÁRIO}{} Suponhamos que $X$ seja plano e próprio sobre $S$, e que $Y$
seja localmente de tipo finito sobre $S$. Então $F$ é pro-representável e as
componentes locais da álgebra topológica correspondente sobre $\Lambda$ são
anéis noetherianos.

\result{OBSERVAÇÕES}{} Os problemas considerados nesta seção, e muitos
outros, eram habitualmente estudados, no âmbito da geometria algébrica
clássica, mediante as «coordenadas de Chow» dos ciclos no espaço
projetivo, que permitiam considerar esses ciclos como pontos de variedades
projetivas adequadas. Esse procedimento, e em geral o uso das
coordenadas de Chow, parece irremediavelmente insuficiente do ponto de
vista dos esquemas, pois destrói os elementos nilpotentes nas
variedades de parâmetros e, em particular, não se presta a um estudo
satisfatório das variações infinitesimais de ciclos (sem falar de sua
natureza não intrínseca, ligada ao espaço projetivo). A linguagem das
coordenadas de Chow era, infelizmente, a única utilizada por muitos
geômetras algébricos para estudar famílias de variedades ou famílias de
ciclos, o que parece ter constituído um sério obstáculo para esclarecer essas
noções, apesar de seu indubitável interesse técnico (provavelmente provisório). Se
se quiser obter o análogo das variedades de Chow na teoria dos
esquemas, chega-se ao seguinte problema universal: seja $X$ um pré-esquema sobre
$S$; para todo pré-esquema $T$ sobre $S$, considera-se o conjunto $F(T)$ dos
subesquemas fechados de $X_T=X\times_ST$ que são planos sobre $T$. Tenta-se
representar este funtor em $T$ mediante um pré-esquema sobre $S$. Mais
geralmente, pode-se dar um feixe quase coerente $G$ sobre $X$ e tomar como
$F(T)$ o conjunto dos feixes

\src{383}
quociente de $G_T$ que são planos sobre $T$. Parece que o problema tem uma
solução, mediante um esquema $C$ localmente de tipo finito sobre $S$, quando
$X$ é próprio sobre o pré-esquema localmente noetheriano $S$ e $G$ é além disso
coerente. Em qualquer caso, supondo unicamente que $S$ seja localmente
noetheriano, a restrição de $F$ às «$S$-álgebras artinianas» é
pro-representável; e se, além disso, $X$ é próprio sobre $S$ e $G$ é coerente, as
componentes locais do anel topológico correspondente $O$ são
noetherianas. Naturalmente, mesmo depois de demonstrar a existência do
«esquema de Chow» de $X$ sobre $S$, restará
encontrar uma decomposição em abertos disjuntos $C_i$ (correspondentes à
fixação de invariantes contínuos, como o grau e a dimensão dos
ciclos que se fazem variar) que sejam de tipo finito (e não somente localmente de
tipo finito) sobre $S$, precisar as relações deste esquema com as
variedades de Chow clássicas e precisar quando um $C_i$ é inclusive projetivo,
ou ao menos quase-projetivo, sobre $S$.\footnote{A errata de 1962
esclarece que o caso projetivo fica completamente resolvido pelos esquemas de
Hilbert (TDTE IV), enquanto exemplos de Nagata e Hironaka mostram que, sem
uma hipótese de projetividade, esses funtores podem não ser representáveis,
mesmo para as subvariedades de dimensão zero de uma variedade completa não
singular de dimensão três; esse fenômeno está ligado à possível
inexistência do quadrado simétrico.}

\subsection*{3. Esquemas de Picard}
\addcontentsline{toc}{subsection}{3. Esquemas de Picard}

Seja $f:X\longrightarrow S$ um $S$-pré-esquema,\footnote{Para um estudo mais
completo, a errata de 1962 remete a TDTE V.} e tomemos o feixe
multiplicativo $\mathcal O_X^*$ das unidades do feixe estrutural de $X$, e
o grupo
\[
  P(X/S)=H^0\bigl(S,R^1f_*(\mathcal O_X^*)\bigr),
\]
chamado grupo de Picard relativo de $X/S$. Um elemento deste grupo fica,
pois, definido dando uma cobertura aberta $(U_i)$ de $S$ e um feixe invertível
$L_i$ sobre cada $f^{-1}(U_i)$, de modo que, para todo par $(i,j)$,
$L_i|f^{-1}(U_i\cap U_j)$ seja isomorfo a $L_j|f^{-1}(U_i\cap U_j)$, ao menos
localmente sobre $U_i\cap U_j$ (isto é, esses dois feixes sejam «equivalentes»
no sentido de [3], B, 4). Se $X/S$ admite uma seção, então $P(X/S)$ nada mais
é que o conjunto das classes de feixes invertíveis sobre $X/S$ módulo
«equivalência» (loco citato). Ponhamos agora, para todo $T$ sobre $S$,
\[
  F(T)=P(X_T/T) ;
\]
obtém-se um funtor contravariante em $T$, que poderá ser chamado funtor de Picard de
$X/S$; se esse funtor for representável, o pré-esquema sobre $S$ que o
representa será chamado pré-esquema de Picard de $X/S$ e será denotado por
$\underline P(X/S)$. Ter-se-á, pois, um isomorfismo de funtores
\[
  \operatorname{Hom}_S\bigl(T,\underline P(X/S)\bigr)
  \xrightarrow{\sim}P(X_T/T).
\]
A formação de pré-esquemas de Picard é compatível com a extensão da
base; em particular, os pré-esquemas de Picard das fibras de $X$ sobre $S$
(que são pré-esquemas sobre os corpos residuais $K(s)$ dos $s\in S$) são
as fibras de $\underline P(X/S)$. Naturalmente, como $P(X_T/T)=F(T)$ é um
grupo comutativo, os pré-esquemas de Picard são pré-esquemas em grupos.
Observe-se também que as jacobianas generalizadas de Rosenlicht são precisamente
as componentes conexas da unidade nos esquemas de Picard de curvas
completas que podem ter singularidades, o que deveria tornar evidente a
maioria de suas propriedades (uma vez demonstrada a existência).

\result{OBSERVAÇÃO}{} A definição adotada aqui só é razoável quando
todo ponto de $S$\footnote{O texto impresso diz «$Y$»; aqui não se definiu
nenhuma base $Y$, e a condição se refere à base $S$.} admite uma vizinhança
aberta $U$ sobre a qual $X$ admite uma seção. No caso geral, é preciso
modificar ligeiramente a definição do funtor de Picard para poder obter
também um teorema de existência.

Aqui, as condições plausíveis para a existência de um pré-esquema de Picard
são as seguintes: $X$ é próprio e plano sobre $S$,
$f_*(\mathcal O_X)=\mathcal O_S$, e $X$ admite localmente uma seção sobre
$S$. Essa condição aparece de maneira natural ao aplicar a técnica de
descida, para eliminar os automorfismos de um feixe invertível $L$ sobre $X$
dotando-o de uma seção marcada sobre a seção $s$ ([3], B, 4). Obtém-se,
em particular:

\src{384}
\result{PROPOSIÇÃO}{3.1} Suponhamos que $X$ seja plano sobre
$S=\operatorname{Spec}(\Lambda)$, com $\Lambda$ noetheriano, e que, para todo
$T$ de tipo finito sobre $S$, tem-se
\[
  f_{T*}(\mathcal O_{X_T})=\mathcal O_T
\]
(se $f$ é próprio e separável,\footnote{A errata de 1962 prescreve
explicitamente a hipótese «$f$ próprio e separável».} ou se $S$ é o espectro
de um corpo, resulta de Künneth que a última condição equivale a
$f_*(\mathcal O_X)=\mathcal O_S$). Então o funtor de Picard de $X/S$ na
categoria das $\Lambda$-álgebras artinianas é pro-representável.

Além disso, tem-se aqui
\[
  F(I_A,\xi)=H^1(X_A,\mathcal O_{X_A}),
\]
em particular:

\result{COROLÁRIO}{} Se $X$ é próprio sobre $S$, então as componentes
locais da $\Lambda$-álgebra topológica $O$ correspondente ao funtor de
Picard são noetherianas.

\result{OBSERVAÇÕES}{} Podem-se generalizar as definições e os resultados
desta seção à classificação dos fibrados principais sobre $X$, tendo como
grupo estrutural um esquema em grupos $G$ sobre $S$ que seja afim, plano sobre
$S$ e comutativo. Quando $G$ não é comutativo, e portanto o fibrado em
grupos adjunto de um fibrado principal (cujas seções são os automorfismos
do fibrado principal) deixa de ser trivial, a proposição 3.1 já não é válida
tal como está. Não obstante, pode-se modificar o problema universal de modo que
continue a ter solução (ao menos, por enquanto, em geometria formal). A regra de
ouro que se deve recordar, no contexto desta \src{385}seção e das
seguintes, e sempre que se busquem «esquemas de módulos» para classes de
objetos definidos apenas a menos de isomorfismo, é sempre a seguinte: eliminar
os possíveis automorfismos dos objetos que se querem classificar mediante
a introdução, se necessário, de estruturas auxiliares (pontos ou
elementos de seções marcadas, fixação de formas diferenciais, etc.), que
serão escolhidas suficientemente inócuas para não modificar substancialmente o
problema inicial.

\subsection*{4. Módulos formais de uma variedade}
\addcontentsline{toc}{subsection}{4. Módulos formais de uma variedade}

Seja $\Lambda$ um anel local noetheriano de corpo residual $k$ (com maior
frequência, $\Lambda$ será igual a $k$, ou a um $p$-anel de Cohen), e seja
$X_0$ um pré-esquema sobre $k$. Para toda $\Lambda$-álgebra artiniana local $A$,
considera-se o conjunto $F(A)$ das classes, salvo isomorfismo, de
$A$-pré-esquemas $X$ planos sobre $A$, dotados de um isomorfismo
\[
  (*)\qquad
  X\otimes_A\kappa(A)\xleftarrow{\sim}
  X_0\otimes_k\kappa(A),
\]
onde $\kappa(A)$ é o corpo residual de $A$; naturalmente, os isomorfismos
entre tais $A$-pré-esquemas planos devem respeitar o isomorfismo precedente
incluído na estrutura. Se $A$ é uma $\Lambda$-álgebra artiniana não
necessariamente local, de componentes locais $A_i$, toma-se como $F(A)$ o
produto dos $F(A_i)$. Desse modo, $F$ torna-se um funtor
multiplicativo em $A$, que poderá ser chamado funtor de módulos para $X_0$ (e
$\Lambda$). Se esse funtor for representável, corresponde-lhe uma
$\Lambda$-álgebra topológica $O$ local, de corpo residual $k$, e o espectro
formal de $O$ será chamado esquema formal de módulos para $X_0$ (e $\Lambda$);
cf. [2] para alguns detalhes sobre essa noção.

Aqui, para aplicar a técnica de descida, os automorfismos «finitos» de $X_0$
são inofensivos, pois não influem na existência de automorfismos (no
sentido precisado acima) dos $A$-pré-esquemas $X$; quando $A$ não se
reduz a um corpo, a condição necessária e suficiente para que $X$ não tenha
nenhum $A$-automorfismo não trivial é que se tenha
\[
  H^0(X_0,\Theta_{X_0/k})=0,
\]
onde $\Theta_{X_0/k}$ denota o feixe das $k$-derivações (= feixe tangente)
de $X_0$. Por outro lado, obtém-se facilmente (ao menos se $X_0$ for liso
sobre $k$)
\[
  F(I_A,\xi)=H^1(X_A,\Theta_{X_A/A}).
\]
Conclui-se então, como de costume:

\src{386}
\result{PROPOSIÇÃO}{4.1} Suponhamos que
$H^0(X_0,\Theta_{X_0/k})=0$. Então existe o esquema formal de módulos para
$X_0$. Se, além disso, $X_0$ for próprio sobre $k$, o esquema formal de módulos é
noetheriano.

\result{OBSERVAÇÕES}{} 1\textsuperscript{o} Quando não se supõe que $X_0$ seja
liso sobre $k$, $F(I_A,\xi)$ identifica-se com um sub-$A$-módulo de
\[
  \operatorname{Ext}^1_{\mathcal O_{P_A}}
  \bigl(P_A;\mathcal I_{X_A},\mathcal O_{X_A}\bigr),
\]
onde se põe $P_A=X_A\times_\Lambda X_A$, onde $\mathcal O_{X_A}$ se
considera como um feixe coerente sobre $P_A$ mediante o morfismo diagonal
$X_A\longrightarrow P_A$, e onde $\mathcal I_{X_A}$ denota o feixe coerente
de ideais sobre $P_A$ definido pelo morfismo diagonal. Mais precisamente,
uma globalização simples da teoria de Hochschild mostra que o
$\operatorname{Ext}^1$ anterior identifica-se com o conjunto das classes,
salvo isomorfismo, de feixes de $I_A$-álgebras planas $\mathcal O$ sobre $X_A$,
dotados de um isomorfismo de aumentação
\[
  \mathcal O\otimes_{I_A}A\longrightarrow\mathcal O_{X_A}
\]
(recorde-se que se põe $I_A=A[t]/(t^2)$). O submódulo $F(I_A,\xi)$ é o
que corresponde aos feixes de álgebras comutativas. As hipóteses de lisura
não são, pois, essenciais na teoria de módulos, contrariamente ao que [2]
dava a entender.

2\textsuperscript{o} Recordemos (loco citato) que, em particular, toda curva
algébrica $X_0$ lisa e própria sobre $k$ admite um esquema formal de módulos,
liso sobre $\Lambda$, de dimensão relativa $3g-3$ se o gênero $g$ é
$\geq2$, e $g$ se $g=0$ ou $1$. Esses dois últimos casos já não entram diretamente
na proposição 4.1. No entanto, no caso das curvas elípticas
($g=1$), poderiam ser reduzidos a ela mediante as considerações que seguem.

Naturalmente, pode-se variar ad libitum a proposição 4.1 considerando
sistemas de esquemas sobre $k$ dotados de diversas estruturas. Suponhamos,
por exemplo, que $X_0$ seja um esquema abeliano sobre $k$, com ponto de origem
marcado (isto é, que $X_0$ se considere como esquema em grupos sobre $k$), e
seja $F(A)$ o conjunto das classes, salvo isomorfismo, de esquemas abelianos
sobre $A$ (isto é, de esquemas em grupos próprios e lisos sobre $A$) dotados
de um isomorfismo de esquemas abelianos $(*)$. Verifica-se que a imposição
de uma estrutura multiplicativa (e inclusive de uma mera «seção unidade»)
elimina os automorfismos infinitesimais e que, por conseguinte, existe um
esquema formal de módulos, correspondente a um anel local completo
noetheriano $O$. Por outro lado, pode-se demonstrar que, se $X$ é um esquema
próprio e liso com fibras «absolutamente conexas» sobre um pré-esquema localmente
noetheriano $S$, toda estrutura multiplicativa sobre $X$ que admita uma
seção unidade é necessariamente associativa e comutativa (sempre que o seja,
ao menos, sobre uma fibra e que $S$ seja conexo), e além disso fica determinada de
maneira única pelo conhecimento

\src{387}
da seção unidade. Além disso, suponhamos que $S$ seja o espectro de um anel
local artiniano $A$ de corpo residual $k$, que $X$ seja próprio sobre $A$ e esteja
dotado de uma seção $s$, e que, por fim, $X\otimes_Ak$ esteja dotado de uma
estrutura de esquema abeliano sobre $k$ cujo elemento neutro é o ponto de
$X\otimes_Ak$ correspondente a $s$; um cálculo simples de obstruções,
unido a um raciocínio devido a Serre, permite demonstrar que existe
efetivamente sobre $X$ uma estrutura multiplicativa que admite a seção
$s$ como seção unidade. (A partir daí, utilizando o «teorema de
existência» de [2] para passar ao caso em que $A$ é noetheriano local completo,
e depois a técnica de descida de [3] para o caso geral, pode-se demonstrar
o enunciado análogo para todo $S$ localmente noetheriano conexo.) Isso prova
que o funtor $F(A)$ considerado aqui é isomorfo ao funtor análogo definido ao
princípio desta seção, prescindindo da estrutura multiplicativa sobre
$X_0$. Segue-se, em particular, que, se $\mathfrak m$ é o ideal maximal de
$O$, $\mathfrak m/\mathfrak m^2$ é canonicamente isomorfo ao dual de
$H^1(X_0,\Theta_{X_0/k})$ e tem, portanto, dimensão $n^2$, onde
$n=\dim X_0$. Seria muito interessante determinar se $O$ é liso sobre $\Lambda$,
isto é, isomorfo a uma álgebra de séries formais em $n^2$ indeterminadas
sobre $\Lambda$. O parágrafo 1, proposição 5.2, permite dar uma formulação
equivalente deste problema como problema de existência de esquemas abelianos
que se reduzam a um esquema abeliano dado. Em qualquer caso, por métodos
transcendentes, vê-se que a resposta é afirmativa se $k$ tem característica
zero. Em característica $p\ne0$, basta evidentemente limitar-se ao caso em que
$\Lambda$ é o anel de vetores de Witt construído sobre um corpo
algebricamente fechado $k$. Talvez tenha chegado o momento para o «funtor
de Greenberg» demonstrar seu valor\ldots\footnote{A errata de 1962
anuncia que o esquema formal de módulos de uma variedade abeliana sobre um
corpo é liso sobre o anel de Witt: todo esquema abeliano sobre um quociente
artiniano local se levanta ao anel do qual é quociente. A demonstração
utiliza a variância da classe de obstrução ao levantamento introduzida na
exposição 182, p. 12. Também recorda as construções de Mumford para
as curvas de gênero $g$ e os esquemas abelianos polarizados (cf. SHMT).}

\subsection*{5. Extensão dos recobrimentos}
\addcontentsline{toc}{subsection}{5. Extensão dos recobrimentos}

Sejam $\mathfrak X$ um pré-esquema formal noetheriano [2] e $U$ uma parte
aberta de $\mathfrak X$, definida localmente pela «não anulação» de uma
seção de $\mathcal O_{\mathfrak X}$ que não seja divisor de zero não somente sobre
$\mathfrak X$, mas também sobre todo $X_n$,\footnote{Precisão imposta pela
errata de 1962.} e que seja, portanto, grande o bastante para que toda seção
de $\mathcal O_{\mathfrak X}$ sobre um aberto $V$ que seja nula sobre $U\cap V$
seja nula. Seja $\mathcal J$ um «ideal de definição» para $\mathfrak X$, e seja
\[
  X_n=(\mathfrak X,\mathcal O_{\mathfrak X}/\mathcal J^{n+1}),
\]
que é, portanto, um pré-esquema noetheriano ordinário. Se então
$\mathfrak X',\mathfrak X''$ são dois recobrimentos planos de $\mathfrak X$ (isto
é, dois pré-esquemas sobre $\mathfrak X$ definidos por feixes de álgebras que
são coerentes e localmente livres como feixes de módulos), não ramificados sobre
$U$, a aplicação evidente
\[
  \operatorname{Hom}_{\mathfrak X}(\mathfrak X',\mathfrak X'')
  \longrightarrow
  \operatorname{Hom}_{X_0}(X'_0,X''_0)
\]
é injetiva; em particular, um automorfismo de $\mathfrak X'$ que induz a
identidade sobre $X'_0$ é a identidade. Isso permite aplicar à situação a
técnica de descida. Partamos, em particular, de um recobrimento plano $X'_0$
de $X_0$, não ramificado

\src{388}
sobre $U_0$; seja $G(\mathfrak X)$ o conjunto das classes, salvo isomorfismo
(que induzam a identidade sobre $X'_0$), de recobrimentos planos $\mathfrak X'$
de $\mathfrak X$ que induzem $X'_0$ sobre $X_0$ (e são, portanto,
necessariamente não ramificados sobre $U$). Define-se do mesmo modo $G(V)$ para
toda parte aberta $V$ de $\mathfrak X$, e mais geralmente $G(\mathfrak Y)$
para todo pré-esquema formal $\mathfrak Y$ sobre $\mathfrak X$. Dito isto, os
resultados de [2] e [3] implicam em primeiro lugar os resultados seguintes:

\noindent a. Quando $V$ é um aberto variável de $\mathfrak X$, os $G(V)$
formam um feixe sobre $\mathfrak X$, seja $G_{\mathfrak X}=G$. A restrição desse
feixe a $U$ é o feixe constante cujas fibras se reduzem a um elemento.

Mais geralmente, a determinação das fibras de $G_{\mathfrak X}$ é uma
questão de anéis locais completos; mais precisamente:

\noindent b. Para todo $x\in\mathfrak X$, tem-se
\[
  G_x=G\bigl(\operatorname{Spec}(\mathcal O_{\mathfrak X,x})\bigr)
  \subset
  G\bigl(\operatorname{Spec}(\widehat{\mathcal O}_{\mathfrak X,x})\bigr),
\]
(= classes, salvo isomorfismo, de álgebras finitas e livres $B$ sobre
$\widehat{\mathcal O}_{\mathfrak X,x}$ dotadas de um isomorfismo de
$B\otimes_{\widehat{\mathcal O}_{\mathfrak X,x}}\mathcal O_{X_0,x}$ com
$\widehat{\mathcal O}'_{0x}$, onde $\mathcal O'_0$ é o feixe de álgebras sobre
$X_0$ que define $X'_0$).\footnote{A fórmula impressa identifica diretamente
$G_x$ com o termo completado e omite o acento circunflexo sobre
$\mathcal O'_{0x}$. A inclusão e o completamento são os prescritos pela
errata de 1962.}

\noindent c. Tem-se um isomorfismo canônico
\[
  G_{\mathfrak X}=\varprojlim_nG_{X_n} ;
\]
em outros termos, para todo aberto $V$ de $\mathfrak X$, tem-se
\[
  G(V)=\varprojlim_nG(V_n).
\]

\noindent d. Suponhamos que $\mathfrak X$ seja obtido de um esquema ordinário
próprio $X$ sobre um anel local noetheriano completo $\Lambda$ que possui um
ideal de definição $\mathfrak m$, tomando o completamento $\mathcal J$-ádico de
$\mathcal O_X$, onde $\mathcal J=\mathfrak m\mathcal O_X$. Então
$G(\mathfrak X)$ é canonicamente isomorfo ao conjunto das classes de
recobrimentos planos do esquema ordinário $X$ que «se reduzem a $X'_0$».

De maneira gráfica, pode-se dizer que (a) e (b) estabelecem as relações
fundamentais entre o aspecto local e o aspecto global do problema, (c) dá
as relações entre o aspecto «finito» e o aspecto «infinitesimal» e, por
último, (d) recorda (sob as condições especificadas) a identidade entre o
aspecto «formal» e o aspecto «algébrico».

Suponhamos agora que $\mathfrak X$ esteja definido sobre um anel local
noetheriano completo $\Lambda$, com
$\mathcal J=\mathfrak m\mathcal O_{\mathfrak X}$ (portanto, $X_0$ é um
pré-esquema sobre $\Lambda/\mathfrak m$). Para toda álgebra $A$ finita sobre
$\Lambda$, ponhamos
\[
  F(A)=G(\mathfrak X\times_\Lambda A),
\]
onde a marcação é tomada após a mudança de base para a álgebra reduzida $A_{\mathrm{red}}$, e não literalmente para $A/\mathfrak mA$. É um funtor covariante em $A$, com valores na categoria dos conjuntos,
e, por (c), esse funtor fica completamente determinado se for conhecido para $A$
artiniana; dizer que esse funtor é pro-representável, isto é, da forma
\[
  F(A)=\operatorname{Hom}_{\Lambda\text{-algèbres top.}}(O,A),
\]

\src{389}
onde $O$ é uma álgebra topológica sobre $\Lambda$ do tipo considerado em A,
parágrafo 5, equivale a dizer que o mesmo ocorre ao se restringir aos $A$
artinianos. A conjunção dos teoremas 1 e 2 implica então efetivamente:

\result{PROPOSIÇÃO}{5.1} Desde que toda família artiniana classificada pelo funtor seja plana sobre sua álgebra de coeficientes e que o aberto escolhido permaneça esquematicamente denso após as mudanças de base pertinentes, o funtor precedente é pro-representável.

Naturalmente, segundo (a), se $U=\mathfrak X$, $G(\mathfrak Y)$ se reduz a um
único elemento para todo $\mathfrak Y$ sobre $\mathfrak X$, e o funtor $F$
tem então escasso interesse (ter-se-á $O=\Lambda$). Parece que, em quase
qualquer outro caso, o anel local topológico $O$ não é noetheriano. Não
obstante, sua existência põe em evidência, de maneira surpreendente, a
natureza «contínua» do conjunto $G(\mathfrak X)$ das soluções (que
corresponde intuitivamente ao fato de que há uma escolha «contínua» no
modo como a ramificação se propaga ao estender $X'_0$). Compare-se esse
resultado com o ponto de vista de J.-P. Serre [4] na teoria do corpo de
classes local, que põe igualmente em relevo o caráter contínuo do grupo de
Galois topológico da extensão abeliana maximal de um corpo local
«geométrico»; ali, o grupo dual (no sentido de Pontryagin) aparece como um
limite indutivo de grupos algébricos (ou ao menos quase algébricos); também
aqui, a classificação das extensões é dada por «variedades» de
dimensão infinita. Por outro lado, no que precede, pode-se tomar como
$\mathfrak X$ o espectro formal de um anel local completo (do qual
$\Lambda$ será, por exemplo, um subanel de Cohen), e cabe esperar que os
resultados desta seção possam ser utilizados no estudo das extensões
de um anel local completo de dimensão $>1$. Tanto no caso local quanto no
global, talvez permitam formular relações precisas entre os fenômenos de
ramificação superior e fenômenos em característica zero (abordáveis por
métodos transcendentes). Em qualquer caso, a análise preliminar da
proposição 5.1 é o que permite estender ao caso «moderadamente ramificado»
os métodos expostos em [2] para estudar o grupo fundamental e resolver por
métodos transcendentes o «problema dos três pontos».

Para terminar, assinalemos que a situação se simplifica quando $X_0$ tem dimensão 1: então, em virtude de (a), complementado pela colagem formal nos esquemas locais completados perfurados,
\[
  G(\mathfrak X)\simeq
  \prod_iG\bigl(\operatorname{Spec}(\widehat{\mathcal O}_{\mathfrak X,x_i})\bigr),
\]
onde os $x_i$ são os pontos de $X_0-U$: podem-se dar arbitrariamente as
extensões «locais» nos pontos de ramificação. Além disso, a normalidade de $X_0$ por si só não implica que o esquema formal de módulos de 5.1 seja liso sobre $\operatorname{Spec}(\Lambda)$; a última afirmação impressa é falsa em geral.

\src{390}
\begin{center}
  \textbf{BIBLIOGRAFIA}
\end{center}

\begin{description}
  \item[{[1]}] Grothendieck (Alexander). --- \emph{Sur quelques points
    d'algèbre homologique}, Tohoku Math. J., t. 9, 1957, p. 119--221.
  \item[{[2]}] Grothendieck (Alexander). --- \emph{Géométrie formelle et
    géométrie algébrique}, Séminaire Bourbaki, t. 11, 1958/59, fasc. 3,
    n\textsuperscript{o} 182.
  \item[{[3]}] Grothendieck (Alexander). --- \emph{Technique de descente et
    théorèmes d'existence en géométrie algébrique, I : Généralités, Descente
    par morphismes fidèlement plats}, Séminaire Bourbaki, t. 12, 1959/60,
    fasc. 1, n\textsuperscript{o} 190.
  \item[{[4]}] Serre (Jean-Pierre). --- \emph{Corps locaux et isogénies},
    Séminaire Bourbaki, t. 11, 1958/59, fasc. 3, n\textsuperscript{o} 185.
\end{description}

\par

\def\FgaUnitKey{e195}\clearpage
\chapter*{Errata da exposição n.\textsuperscript{o} 195}
\addcontentsline{toc}{chapter}{Errata da exposição n.\textsuperscript{o} 195}
\setcounter{footnote}{0}
\markboth{Errata da exposição 195}{Errata da exposição 195}
\sid{e195}
\src{300}

\noindent\textbf{4. TDTE II: O teorema de existência na teoria formal dos
módulos.}

\begin{center}
[Seminário Bourbaki, t. 12, 1959/60, n.\textsuperscript{o} 195, 22 p.]
\end{center}

\noindent\textbf{Página 8.} --- A fórmula escrita na proposição 5.1 só é
correta quando $\Lambda$ é um corpo; no caso geral, deve-se substituir
$\mathfrak m_\xi/\mathfrak m_\xi^2$ pelo quociente deste espaço pela
imagem de $\mathfrak n_\xi/\mathfrak n_\xi^2$, onde $\mathfrak n$ é o ideal
maximal de $\Lambda$. Além disso, a definição dada para que $O$ seja liso sobre
$\Lambda$ só é correta quando a extensão residual $k'/k$ é separável.
No caso geral, cf. SGA III, 1.1.

\noindent\textbf{Página 14, observação.} --- Os problemas formulados aqui
ficam completamente resolvidos no caso projetivo mediante os «esquemas de
Hilbert» (cf. TDTE IV). Exemplos de Nagata e Hironaka mostram, em contrapartida, que
os funtores considerados não são necessariamente representáveis se não se impõe
uma hipótese de projetividade, mesmo limitando-se à classificação das
subvariedades de dimensão zero de uma variedade completa não singular de
dimensão 3; isso está ligado ao fato de que o quadrado simétrico de uma
tal variedade pode não existir.

\src{301}
\noindent\textbf{Página 15, n.\textsuperscript{o} 3.} --- Para um estudo mais
completo, veja-se TDTE V.

\noindent\textbf{Página 16, proposição 3.1.} --- Em vez de «se $f$ é
próprio», leia-se «se $f$ é próprio e separável».

\noindent\textbf{Páginas 18 e 19, observações, 2.\textsuperscript{o}} --- Demonstrei
recentemente que o esquema formal de módulos de uma variedade
abeliana sobre um corpo é efetivamente liso sobre o anel de Witt; em
outros termos, todo esquema abeliano sobre um anel artiniano local que seja
quociente de outro provém, por redução, de um esquema abeliano sobre este
último. A demonstração utiliza unicamente as propriedades de variância da
classe de obstrução ao levantamento introduzida na exposição
n.\textsuperscript{o} 182, página 12. Recordemos igualmente que os esquemas de
módulos das curvas de gênero $g$ e dos esquemas abelianos polarizados foram
construídos por Mumford (cf. SHMT).

\noindent\textbf{Página 19, parágrafo 5.} --- Deve-se supor que a seção
considerada para definir $U$ seja um não divisor de zero não somente sobre
$\mathfrak X$, mas também sobre todo $X_n$.

\noindent\textbf{Página 20, linhas 12 e 13.} --- Em vez de
\[
  G_x=G\bigl(\operatorname{Spec}(\widehat{\mathcal O}_{\mathfrak X,x})\bigr),
\]
leia-se
\[
  G_x=G\bigl(\operatorname{Spec}(\mathcal O_{\mathfrak X,x})\bigr)
  \subset
  G\bigl(\operatorname{Spec}(\widehat{\mathcal O}_{\mathfrak X,x})\bigr),
\]
e, em vez de «$\mathcal O'_{0x}$», leia-se
«$\widehat{\mathcal O}'_{0x}$».

\par

\end{document}
